% ------------------------------------------------------------------------------
% Este fichero es parte de la plantilla LaTeX para la realización de Trabajos
% Final de Grado 
% Adecuado para proyectos de desarrollo de software del Grado en Ing. Informática
% Copyright (C) Universidad de Cádiz
%% Imágenes obtenidas de Freepik (https://www.freepik.es/)
% ------------------------------------------------------------------------------


\documentclass[12pt, twoside]{report}

% Paquetes LaTeX y estilos globales
\usepackage[utf8]{inputenc}
% \usepackage{subfigure}
\usepackage{caption}
\usepackage{subcaption}
\usepackage{multicol}
\usepackage{xcolor}
\usepackage[spanish,es-tabla]{babel}
\usepackage[utf8]{inputenc}
\usepackage{graphicx}
\usepackage{titlesec}
\usepackage[bookmarks,breaklinks,colorlinks=true,allcolors=blue]{hyperref}
\usepackage{listings}
\usepackage{inconsolata}
\usepackage{float}

% \usepackage[square,numbers]{natbib}
\usepackage{natbib}
\usepackage[nottoc]{tocbibind}

%\AtBeginDocument{
%  \renewcommand{\bibsection}{\chapter{\bibname}}
%} % Bibliografia en capitulo numerado

\usepackage{geometry}
\usepackage{amsmath}
\usepackage{parskip}
\usepackage[official]{eurosym}
\usepackage[colorinlistoftodos,spanish]{todonotes} % compilar con [disable] en la versión final
\setlength{\marginparwidth}{2cm}
\usepackage{csquotes}

\usepackage{lipsum}
\usepackage{multirow}

\usepackage[es]{datetime2}
\DTMlangsetup{showdayofmonth=false}

\selectlanguage{spanish}

%Gestión de anexos
\usepackage{appendix}
\renewcommand{\appendixname}{Anexos}
\renewcommand{\appendixtocname}{Anexos}
\renewcommand{\appendixpagename}{Anexos}

% --- Paquetes añadidos para PymeKit (tablas, listas y diagramas TikZ) ---
\usepackage[T1]{fontenc}
\usepackage{booktabs}
\usepackage{tabularx}
\usepackage{array}
\usepackage{longtable}
\usepackage{enumitem}
\usepackage{tikz}
\usetikzlibrary{positioning,arrows.meta,shapes.geometric,fit,backgrounds,calc}
\usepackage{pgfgantt}
\usepackage{pdflscape}

% Formato del título de capítulos y secciones
% \titleformat{\chapter}[block]{\titlerule[0.8pt]\normalfont\huge\bfseries}{\thechapter.}{.5em}{\Huge}[{\titlerule[0.8pt]}]
% \titlespacing*{\chapter}{0pt}{-19pt}{25pt}
% \titleformat{\section}[block]{\normalfont\Large\bfseries}{\thesection.}{.5em}{\Large}

\titleformat{\chapter}[block]{\normalfont\huge\bfseries}{\thechapter.}{.5em}{\Huge}[{\titlerule[0.8pt]}]
\titlespacing*{\chapter}{0pt}{-19pt}{25pt}
\titleformat{\section}[block]{\normalfont\Large\bfseries}{\thesection.}{.5em}{\Large}

% Formato del código fuente con lstlisting
\lstset{
  basicstyle=\ttfamily,
  breaklines=true,
}

% Márgenes
\geometry{
    a4paper,
    margin=2.75cm
}

% Limite de profundidad del índice
\setcounter{tocdepth}{2}

% Indentación de párrafos
\setlength{\parindent}{1cm}

\renewcommand{\lstlistingname}{Extracto de código}
\renewcommand*{\lstlistlistingname}{Índice de extractos de código}

%------------ Caption debajo de las tablas
\DeclareCaptionLabelSeparator*{spaced}{\\[1.5ex]}
\captionsetup[table]{textfont=it,format=plain,justification=justified,
  singlelinecheck=false,labelsep=spaced,skip=0pt} 

% Con \hypersetup podremos introducir los metadatos para la descripción del documento PDF. Recuerda modificarlo cuando transformes el documento.
\hypersetup{
    pdfauthor={Enrique Mellado},
    pdftitle={Diseño e implementación de una arquitectura software reutilizable para el desarrollo ágil de aplicaciones SaaS orientadas a pymes},
    pdfsubject={Trabajo Fin de Grado. Grado en Ingeniería Informática. Universidad de Cádiz},
    pdfkeywords={SaaS, pymes, arquitectura reutilizable, multi-tenant, Supabase, CMS},
    pdflang={es},
    pdfcreator={LaTeX}
}

% Comienzo del documento
\begin{document}

    % FRONTMATTER: portadas y secciones no numeradas
    


    
\begin{center}

\includegraphics[width=0.45\textwidth]{figures/logoucaesi.jpg} 




\vspace*{2cm}
\begin{Large}
TRABAJO FIN DE GRADO
\end{Large}

\vspace*{1cm}
\begin{Large}
GRADO EN INGENIERÍA INFORMÁTICA
\end{Large}


\vspace*{3.4cm}

\textbf{\LARGE \bigskip Diseño e implementación de una arquitectura software reutilizable para el desarrollo ágil de aplicaciones SaaS orientadas a pymes} % Aquí el título de su trabajo 
% Si el título del trabajo supera un párrafo añadir el comando \bigskip antes del título para que no se solapen dichos párrafos


\vspace*{3cm}

{\Large AUTOR: ENRIQUE MELLADO}\\ % Aquí su nombre



\vspace*{1cm}

Puerto Real, \today\\


\end{center}

\thispagestyle{empty} % Impide que se incluya número de página en la portada

    \newpage{\ }
\begin{center}

\includegraphics[width=0.45\textwidth]{figures/logoucaesi.jpg} 




\vspace*{2cm}
\begin{Large}
TRABAJO FIN DE GRADO
\end{Large}

\vspace*{1cm}
\begin{Large}
GRADO EN INGENIERÍA INFORMÁTICA
\end{Large}


\vspace*{3.4cm}

\textbf{\LARGE \bigskip Diseño e implementación de una arquitectura software reutilizable para el desarrollo ágil de aplicaciones SaaS orientadas a pymes} % Aquí el título de su trabajo 
% Si el título del trabajo supera un párrafo añadir el comando \bigskip antes del título para que no se solapen dichos párrafos


\vspace*{3cm}

{\Large TUTOR: JUAN CARLOS DE LA TORRE MACÍAS}\\ % Aquí su nombre
{\Large AUTOR: ENRIQUE MELLADO}\\ % Aquí su nombre



\vspace*{1cm}

Puerto Real, \today \\


\end{center}

\thispagestyle{empty} % Impide que se incluya número de página en la portada





    \newpage{\ }
\thispagestyle{empty}
\chapter*{Agradecimientos}

\todo{El autor personaliza los agradecimientos}
Quiero expresar mi agradecimiento a mi tutor, Juan Carlos de la Torre
Macías, por su orientación y sus consejos a lo largo de este trabajo, y al
profesorado de la Escuela Superior de Ingeniería de la Universidad de Cádiz
por la formación recibida durante el grado.

Gracias también a mi familia y a mis amigos, por su apoyo y su paciencia
durante todo este tiempo.

    \newpage{\ }
\thispagestyle{empty}
\chapter*{Resumen}

\noindent Las pequeñas y medianas empresas (pymes) necesitan cada vez más
aplicaciones web propias, pero rara vez pueden asumir el coste de construir
desde cero los componentes que casi cualquier aplicación SaaS
(\textit{Software as a Service}) repite: autenticación segura, gestión de
usuarios y equipos con roles, cobro de suscripciones y una herramienta para
administrar los datos sin programar. Este Trabajo Fin de Grado diseña,
implementa y valida PymeKit, una arquitectura software reutilizable que reúne
esos componentes en una base común sobre la que desarrollar de forma ágil
nuevas aplicaciones SaaS orientadas a pymes.

% PENDIENTE (P-06): forma de presentar la parte reutilizada, a confirmar con el tutor.
PymeKit no parte de cero. Integra dos bases de código existentes con
licencia ---una base de aplicación SaaS y un gestor de contenidos (CMS) para
la base de datos--- y las adapta y extiende hasta formar un único producto. Las
aportaciones del trabajo son la integración del CMS dentro de la consola de
administración de la aplicación, con un solo servicio, un solo inicio de
sesión y una sola pila de librerías; la reorganización de la consola por
áreas de negocio comprensibles para una pyme; la interfaz íntegramente en
español sobre una infraestructura de internacionalización extensible; la
centralización de la configuración reutilizable; la corrección documentada de
numerosos fallos de seguridad heredados; y un proceso de desarrollo asistido
por agentes de inteligencia artificial, controlado mediante directrices,
comprobaciones automáticas y revisiones adversariales.

La solución se organiza como un monorepo de paquetes desacoplados con
TypeScript estricto. La aplicación web usa TanStack Start y React; los datos,
la autenticación y el almacenamiento se apoyan en Supabase (PostgreSQL); los
pagos, en Stripe; y la API del CMS, en Hono y Drizzle, montada dentro del
propio servidor web. La autorización se aplica en la base de datos mediante
seguridad a nivel de fila (RLS), de modo que el aislamiento entre clientes
(\textit{multi-tenancy}) no depende solo del código de la aplicación. El
acceso a la consola exige un segundo factor de autenticación, y las acciones
realizadas desde el CMS quedan registradas en un registro de auditoría que el
propio personal no puede falsificar.

La calidad se verifica con pruebas unitarias, más de dos mil pruebas de base
de datos con pgTAP ---incluidas pruebas de ataque entre clientes--- y pruebas
de extremo a extremo con Playwright, ejecutadas en integración continua. El
proceso dejó documentadas en una bitácora más de cincuenta incidencias, buena
parte de ellas fallos de seguridad del código de partida que se detectaron y
corrigieron durante la integración. La memoria describe la planificación, el
análisis, el diseño, la construcción, las pruebas y el despliegue de la
plataforma, la evaluación de su reutilización y los manuales de usuario y de
desarrollador.

\vspace{.5cm}

% Añadir una serie de palabras clave que aparezcan en el texto
\noindent \textbf{Palabras clave:} SaaS, pymes, arquitectura reutilizable,
\textit{multi-tenancy}, seguridad a nivel de fila (RLS), Supabase, PostgreSQL,
CMS, TanStack Start, desarrollo asistido por IA.

    
    % Índice del documento y de figuras
    \begingroup
    % Los enlaces son normalmente azules, pero en los índices se configuran a negro
    % para que no aparezca todo azul
    \hypersetup{linkcolor=black}
    \newpage{\ }
    \thispagestyle{empty}
    \tableofcontents
    \listoffigures
    \listoftables
    \endgroup
    
    \newpage{\ }
    \cleardoublepage

 
     \setcounter{page}{1} % Comienza a incluir números de página a partir de aquí
     % Cambia el estilo de números de página a normal
    \clearpage\pagenumbering{arabic}
    
    % PROLEGÓMENO
    \part{Prolegómeno}
     \null\vfill
    \noindent La primera parte de la memoria del Trabajo Fin de Grado (TFG) debe contener una introducción y una planificación del proyecto. La introducción es un capítulo que, a modo de resumen, debe contener una breve descripción del contexto de la disciplina en la que el proyecto tiene aplicación y la motivación para su desarrollo, así como del alcance previsto. En el segundo capítulo deberán describirse los antecedentes del proyecto, esto es, su contexto y la situación actual. Además, se desarrollará el estado de la técnica en relación con la propuesta de proyecto. El tercer capítulo debe incluir el plan de gestión del proyecto. La planificación deberá ajustarse a las prácticas de ingeniería en general, y de la ingeniería del software en particular. Deberá tener en cuenta los plazos, los entregables (documentos y software), los recursos (humanos y de equipamiento inventariable) y el método de ingeniería de software a emplear.
    \\

\chapter{Introducción}
\label{cap:introduccion}
% ------------------------------------------------------------------------------
% Capítulo 1 · Introducción
% Fuentes: propuesta aprobada del TFG (objetivo, descripción, alcance y EDT),
% docs/tfg/{PLAN,REQUISITOS,DECISIONES}.md.
% ------------------------------------------------------------------------------

Este capítulo presenta la motivación del trabajo, sus objetivos y su alcance,
tal como se fijaron en la propuesta aprobada, y describe la organización del
resto del documento.

\section{Motivación}
\label{sec:intro-motivacion}

Las pequeñas y medianas empresas constituyen la inmensa mayoría del tejido
empresarial europeo~\cite{ec-sme-definition}. Para muchas de ellas, disponer
de una aplicación web propia ---una tienda, una zona de clientes, un sistema de
reservas o una herramienta interna--- ha dejado de ser opcional. El modelo de
\textit{Software as a Service} (SaaS)~\cite{nist-cloud}, en el que la aplicación
se ofrece a través de la red y se paga por uso o por suscripción, es la forma
habitual de construir y explotar este tipo de productos.

Sin embargo, cualquier aplicación SaaS, por sencilla que sea su lógica de
negocio, necesita un conjunto de piezas que se repiten de un proyecto a otro:
registro e inicio de sesión con verificación de correo y segundo factor,
cuentas de usuario y de organización, invitaciones y roles, cobro de
suscripciones, una consola de administración y alguna forma de consultar y
corregir los datos sin escribir consultas a mano. Ninguna de ellas aporta
valor diferencial al producto, pero todas tienen que estar bien hechas: un
error en el control de acceso o en el aislamiento entre clientes compromete
los datos de todas las empresas que usan la aplicación.

Para un equipo pequeño, o para un desarrollador que da servicio a varias
pymes, construir estas piezas desde cero en cada proyecto es lento y
arriesgado. Lo razonable es partir de una base común, ya probada, y dedicar
el esfuerzo a lo que distingue a cada producto. Esa es la idea que motiva este
trabajo: disponer de una \textbf{arquitectura software reutilizable} que
reduzca el tiempo de puesta en marcha de nuevos productos SaaS, facilite la
reutilización de componentes comunes y estandarice buenas prácticas en
seguridad, gestión del acceso, persistencia de datos y administración del
sistema.

\section{Objetivos}
\label{sec:intro-objetivos}

\subsection{Objetivo general}

El objetivo general, tal como figura en la propuesta aprobada, es
\emph{diseñar, implementar y validar una arquitectura software reutilizable
que sirva como base para el desarrollo ágil de aplicaciones SaaS orientadas a
pequeñas y medianas empresas, integrando los componentes esenciales de una
solución moderna: \textit{frontend}, \textit{backend}, sistema de
autenticación, gestión de usuarios y roles, integración de pagos y un módulo
CMS para la administración de contenidos y datos.}

\subsection{Objetivos específicos}

A partir de la descripción y del alcance de la propuesta, el objetivo general
se descompone en los objetivos específicos de la
tabla~\ref{tab:intro-objetivos}. Cada uno se relaciona con los requisitos del
capítulo de Requisitos del Sistema y se retoma en las conclusiones
(capítulo~\ref{cap:conclusiones}) para valorar su grado de cumplimiento.

\begin{table}[htbp]
  \centering
  \small
  \caption{Objetivos específicos del trabajo y requisitos relacionados.}
  \label{tab:intro-objetivos}
  \begin{tabularx}{\textwidth}{@{}lXl@{}}
    \toprule
    \textbf{ID} & \textbf{Objetivo específico} & \textbf{Requisitos} \\
    \midrule
    OE-1 & Analizar el contexto de las aplicaciones SaaS para pymes, identificar los componentes comunes reutilizables y definir los requisitos funcionales y no funcionales de la plataforma base. & Todos \\
    OE-2 & Diseñar una arquitectura modular que separe \textit{frontend}, \textit{backend}, base de datos y servicios auxiliares, con paquetes desacoplados. & RF-01, RF-02, RNF-04 \\
    OE-3 & Diseñar e implementar la autenticación, la autorización y la gestión de usuarios, equipos y roles, aplicando el control de acceso en la propia base de datos. & RF-03--RF-06, RNF-02, RNF-03 \\
    OE-4 & Integrar una pasarela de pagos para la gestión de suscripciones. & RF-07 \\
    OE-5 & Integrar un módulo CMS para la administración de contenidos y de los datos almacenados, con permisos propios y auditoría. & RF-08--RF-11 \\
    OE-6 & Diseñar un modelo de datos multicliente (\textit{multi-tenant}) con relaciones, restricciones y mecanismos de acceso orientados a la reutilización. & RNF-02, RNF-03 \\
    OE-7 & Preparar la plataforma para su reutilización: configuración centralizada, estructura de proyecto reproducible e interfaz en español extensible a otros idiomas. & RF-13, RNF-01, RNF-05 \\
    OE-8 & Validar la solución con pruebas de los principales casos de uso y de la integración entre módulos, y evaluar su reutilización en un escenario de ejemplo. & RNF-01, RNF-06 \\
    OE-9 & Definir el despliegue y elaborar la documentación: memoria, documentación técnica y manuales de instalación y de uso. & RNF-07 \\
    \bottomrule
  \end{tabularx}
\end{table}

\section{Alcance}
\label{sec:intro-alcance}

La propuesta organiza el trabajo en una Estructura de Descomposición del
Trabajo (EDT) con seis paquetes: (1)~análisis del problema y definición de
requisitos; (2)~diseño de la arquitectura software; (3)~diseño del modelo de
datos; (4)~implementación del prototipo funcional; (5)~validación de la
solución; y (6)~documentación. Los entregables comprometidos son la memoria,
el software desarrollado, los ficheros de configuración y despliegue, el
manual de instalación y puesta en marcha y un manual básico de uso. La
relación entre los paquetes de la EDT y las fases reales del proyecto se
detalla en el plan de gestión del proyecto (capítulo~\ref{cap:plan}).

El prototipo debe incluir, como mínimo, gestión de usuarios, autenticación,
persistencia de datos, administración desde un panel CMS e integración con un
flujo de pago. Dentro de ese marco, a lo largo del proyecto se fijaron estas
decisiones de alcance:

\begin{itemize}
  \item \textbf{Plataforma genérica.} PymeKit es una base reutilizable, no
        una aplicación de un sector concreto. Su configuración por defecto
        corresponde a la web de \emph{una} pyme: página pública y blog,
        acceso de clientes y personal, pagos y una consola de administración
        (ADR-005, ADR-020). Los equipos y las suscripciones se conservan
        como módulos que cada instalación activa si los necesita.
  \item \textbf{Un único CMS, de datos.} La propuesta pide un CMS para la
        administración de contenidos y de datos. Se cubre con un CMS que
        trabaja directamente sobre las tablas de la base de datos. El
        contenido editorial (el blog) se guarda también en la base de datos y
        se gestiona desde ese mismo CMS, en lugar de mantener un segundo gestor
        de contenidos basado en ficheros (ADR-004, ADR-017).
  \item \textbf{CMS integrado en la consola.} El CMS no es una aplicación
        aparte: forma parte de la consola de administración de la web, con un
        solo servicio y un solo inicio de sesión (ADR-011).
  \item \textbf{Interfaz solo en español}, sin textos escritos directamente
        en el código, de modo que añadir otro idioma sea un punto de
        extensión documentado (ADR-021).
  \item \textbf{Fuera de alcance:} aplicaciones móviles nativas, la lógica
        de negocio de un sector concreto, la facturación fiscal (emisión de
        facturas con validez legal) y la explotación en producción con
        clientes reales. Los pagos funcionan en el modo de pruebas de la
        pasarela.
\end{itemize}

\paragraph{Aportaciones del trabajo.}
Las aportaciones del trabajo son:

\begin{itemize}
  \item la \textbf{integración del CMS} en la consola de administración de la
        web (ADR-011): montaje de su API dentro del servidor web, pantallas
        implementadas como rutas de la aplicación, sesión única y control de
        acceso coordinado entre los permisos de la plataforma y los del CMS
        (ADR-014, ADR-016) y una sola base de datos para ambos (ADR-012);
  \item la \textbf{unificación de la pila} de librerías de formularios,
        textos y componentes de interfaz en toda la web (ADR-013);
  \item la \textbf{interfaz en español} sobre una infraestructura de
        internacionalización extensible (ADR-021);
  \item la \textbf{consola por áreas de negocio}, con nombres en español
        para las tablas, ayudas contextuales en los formularios y un resumen
        en cada pantalla (ADR-020);
  \item la \textbf{auditoría adversarial y corrección documentada de fallos
        de seguridad}, fijados con pruebas de regresión (ADR-015 y bitácora
        de incidencias);
  \item el \textbf{proceso de desarrollo asistido por IA}, con un arnés de
        directrices, comprobaciones automáticas y revisiones;
  \item la \textbf{configuración reutilizable} centralizada, el
        \textbf{despliegue} y la \textbf{documentación} técnica y de usuario.
\end{itemize}

\section{Glosario de términos}
\label{sec:intro-glosario}

\begin{description}[style=nextline,leftmargin=1.2cm]
  \item[AAL (\textit{Authenticator Assurance Level})] Nivel de garantía de
        una sesión: \texttt{aal1} tras la contraseña o el enlace mágico;
        \texttt{aal2} tras verificar además un segundo factor.
  \item[ADR (\textit{Architecture Decision Record})] Registro breve de una
        decisión de arquitectura con su contexto, alternativas y
        consecuencias~\cite{nygard2011adr}.
  \item[BaaS (\textit{Backend as a Service})] Servicio que ofrece ya
        construidos la base de datos, la autenticación, el almacenamiento y
        una API, de modo que la aplicación no necesita un \textit{backend}
        propio para esas funciones.
  \item[CMS (\textit{Content Management System})] Sistema para gestionar
        contenidos. En este trabajo, un CMS \emph{de datos}: una interfaz para
        consultar y editar las tablas de la base de datos con permisos.
  \item[E2E (\textit{end-to-end})] Prueba que ejercita la aplicación completa
        a través de un navegador, como lo haría un usuario.
  \item[EDT] Estructura de Descomposición del Trabajo (\textit{Work Breakdown
        Structure}).
  \item[Harness (arnés)] Conjunto de directrices, procedimientos y
        comprobaciones automáticas que acota y verifica el trabajo de los
        agentes de IA en el repositorio.
  \item[i18n] Internacionalización: preparar el software para mostrar sus
        textos en varios idiomas sin modificar el código.
  \item[MFA (\textit{Multi-Factor Authentication})] Autenticación con más de
        un factor; aquí, contraseña más un código temporal (TOTP).
  \item[Monorepo] Repositorio único que contiene varias aplicaciones y
        paquetes con dependencias explícitas entre ellos.
  \item[Multi-tenant (multicliente)] Arquitectura en la que una sola
        instalación da servicio a varios clientes (\textit{tenants}) cuyos
        datos deben quedar aislados entre sí.
  \item[pgTAP] Marco de pruebas unitarias para PostgreSQL~\cite{pgtap}.
  \item[Pyme] Pequeña y mediana empresa: menos de 250 personas empleadas y un
        volumen de negocio anual de hasta 50 millones de euros o un balance
        de hasta 43 millones~\cite{ec-sme-definition}.
  \item[RBAC (\textit{Role-Based Access Control})] Control de acceso basado
        en roles a los que se asignan permisos.
  \item[RLS (\textit{Row Level Security})] Mecanismo de PostgreSQL que filtra
        las filas que cada usuario puede leer o modificar mediante
        políticas~\cite{postgresql-rls}.
  \item[SaaS (\textit{Software as a Service})] Modelo en el que el proveedor
        ofrece una aplicación a través de la red sobre una infraestructura
        que gestiona él mismo~\cite{nist-cloud}.
  \item[\textit{Server function}] Función que se ejecuta en el servidor pero
        se invoca desde el cliente como una llamada tipada.
  \item[Super-admin] Administrador de la plataforma, con acceso a la consola
        y a todos los permisos del CMS.
\end{description}

\section{Organización del documento}
\label{sec:intro-organizacion}

La memoria sigue la estructura de la plantilla de la Escuela Superior de
Ingeniería y se divide en cuatro partes:

\begin{itemize}
  \item \textbf{Prolegómeno.} Este capítulo de introducción; el
        capítulo~\ref{cap:contexto}, de antecedentes, con la situación de las
        pymes ante el software SaaS, el estado de la técnica y la
        justificación de la solución; y el capítulo~\ref{cap:plan}, con el
        plan de gestión del proyecto: metodología, tecnologías,
        planificación, organización, costes, riesgos, gestión de la
        configuración y aseguramiento de la calidad.
  \item \textbf{Desarrollo.} Los capítulos de requisitos, análisis, diseño,
        construcción, pruebas y despliegue del sistema, que recorren las
        etapas del proceso de ingeniería con los diagramas, las decisiones y
        los fragmentos de código más significativos.
  \item \textbf{Epílogo.} Las conclusiones (capítulo~\ref{cap:conclusiones}),
        con el grado de cumplimiento de los objetivos, las lecciones
        aprendidas y el trabajo futuro; la información sobre la licencia y la
        bibliografía.
  \item \textbf{Anexos.} El manual del desarrollador
        (anexo~\ref{cap:manual-desarrollador}), el manual de usuario
        (anexo~\ref{cap:manual-usuario}) y la evaluación de la reutilización
        de la plataforma (anexo~\ref{cap:evaluacion}).
\end{itemize}

El software se entrega como un repositorio git con el monorepo completo:
código de la aplicación y del CMS, esquemas, migraciones y pruebas de la base
de datos, pruebas de extremo a extremo, configuración de integración continua
y de despliegue, y la documentación interna del proyecto (plan, requisitos,
decisiones, trazabilidad y bitácora). El repositorio es privado (véase la información sobre
licencia).


\newpage{\ }
\cleardoublepage

\chapter{Antecedentes}
\label{cap:contexto}
% ------------------------------------------------------------------------------
% Capítulo 2 · Antecedentes
% Fuentes: propuesta del TFG, docs/tfg/DECISIONES.md (ADR-001, 004, 005, 011,
% 014, 017, 020), bibliografía. Por ADR-007 no se nombran los componentes base
% utilizados.
% ------------------------------------------------------------------------------

Este capítulo describe la situación de la que parte el trabajo: qué necesitan
las pymes cuando quieren ofrecer o usar software SaaS propio, qué piezas
comparten casi todas esas aplicaciones y qué soluciones existen hoy para no
tener que construirlas desde cero. A partir de ese análisis se justifica la
solución adoptada.

\section{Contexto}
\label{sec:ctx-contexto}

\subsection{Las pymes y el software como servicio}
\label{sec:ctx-pymes-saas}

La Comisión Europea define como pyme a la empresa que emplea a menos de 250
personas y cuyo volumen de negocio anual no supera los 50 millones de euros o
cuyo balance general no supera los 43 millones~\cite{ec-sme-definition}. Son
la gran mayoría de las empresas, pero sus recursos técnicos son limitados:
pocas tienen un departamento de informática, y las que encargan una aplicación
a medida lo hacen con presupuestos y plazos ajustados.

El modelo SaaS encaja con esa situación. Según la definición del
NIST~\cite{nist-cloud}, en este modelo el proveedor ofrece una aplicación que
se ejecuta en una infraestructura en la nube gestionada por él; el cliente la
usa a través del navegador sin preocuparse de servidores, actualizaciones ni
copias de seguridad, y paga por suscripción. Para una pyme, el SaaS aparece en
dos papeles:

\begin{itemize}
  \item como \textbf{consumidora}, cuando contrata una herramienta en lugar de
        instalarla; y
  \item como \textbf{proveedora}, cuando su propio negocio se ofrece en
        forma de aplicación web: una zona privada para sus clientes, un
        servicio de reservas, un catálogo con suscripción o, simplemente, la
        web de la empresa con un acceso para clientes y otro para el
        personal.
\end{itemize}

Este trabajo se sitúa en el segundo caso, y en particular en el de quien
\emph{desarrolla} esas aplicaciones para pymes: un estudio pequeño o un
desarrollador independiente que debe entregar, con poco presupuesto, una
aplicación segura y mantenible, y que probablemente tendrá que repetir el
ejercicio para otro cliente poco después.

\subsection{Componentes comunes de una aplicación SaaS}
\label{sec:ctx-componentes}

Por diferente que sea el negocio, una aplicación SaaS repite un conjunto de
componentes transversales. La tabla~\ref{tab:ctx-componentes} los resume; son,
precisamente, los que la propuesta del trabajo enumera como esenciales.

\begin{table}[htbp]
  \centering
  \small
  \caption{Componentes que se repiten en casi cualquier aplicación SaaS.}
  \label{tab:ctx-componentes}
  \begin{tabularx}{\textwidth}{@{}>{\raggedright\arraybackslash}p{3.6cm}X@{}}
    \toprule
    \textbf{Componente} & \textbf{Qué debe resolver} \\
    \midrule
    Interfaz pública y privada & Página de presentación, precios y contenido; área privada con navegación y ajustes. \\
    Autenticación & Registro, verificación del correo, inicio de sesión con contraseña o enlace, recuperación de la contraseña, segundo factor (MFA) y, opcionalmente, proveedores externos. \\
    Cuentas y organizaciones & Cuenta personal, cuentas de organización con miembros, invitaciones, roles y permisos por funcionalidad. \\
    Aislamiento entre clientes & Garantizar que los datos de una organización no son visibles ni modificables desde otra (\textit{multi-tenancy}). \\
    Pagos & Planes, alta de la suscripción, portal del cliente y sincronización del estado con la pasarela de pago. \\
    Administración & Consola para el personal: buscar usuarios, bloquearlos, darles soporte, consultar métricas. \\
    Gestión de datos y contenidos & Consultar y corregir datos, publicar contenidos y gestionar archivos sin programar, con permisos y trazabilidad. \\
    Transversales & Correos transaccionales, notificaciones, internacionalización, registro de errores, pruebas y despliegue. \\
    \bottomrule
  \end{tabularx}
\end{table}

Ninguno de estos componentes diferencia un producto de otro, pero su coste
de desarrollo es alto y su margen de error, muy pequeño. La literatura sobre
aplicaciones multicliente insiste en que el aislamiento, la configurabilidad
por cliente y la seguridad son las preocupaciones de arquitectura que más
condicionan estos sistemas y las que más complican su
mantenimiento~\cite{bezemer2010multitenant,krebs2012multitenancy}. La
ingeniería del software lleva décadas proponiendo la misma respuesta a este
tipo de problema: identificar lo común a una familia de productos y
construirlo una vez, como una base reutilizable sobre la que se derivan los
productos concretos~\cite{clements2001spl}.

\section{Estado de la técnica}
\label{sec:ctx-estado-tecnica}

Se revisan a continuación las cuatro familias de soluciones que cubren, total
o parcialmente, los componentes anteriores: las plantillas o kits de
aplicaciones SaaS, los servicios de \textit{backend} gestionado, los gestores
de contenidos y paneles de administración de bases de datos, y las técnicas de
aislamiento multicliente en la base de datos.

\subsection{Plantillas y kits de aplicaciones SaaS}
\label{sec:ctx-kits}

Desde hace años existe un mercado de plantillas de arranque
(\textit{boilerplates} o \textit{starter kits}) para aplicaciones SaaS. Son
proyectos que ya incluyen la autenticación, las cuentas de equipo, la
facturación y una página de presentación, construidos sobre un
\textit{framework} web concreto, y que el desarrollador copia y adapta. Se
distribuyen tanto con licencias libres como con licencias comerciales de pago
único, y estas últimas suelen ser más completas: incorporan cuentas de equipo
con roles, panel de administración, pruebas automatizadas y documentación.

Su principal ventaja es el ahorro de tiempo inicial. Sus limitaciones, desde
el punto de vista de este trabajo, son tres:

\begin{enumerate}
  \item rara vez incluyen una herramienta para que el personal no técnico
        administre los datos, más allá de un panel de super-administración
        con unas pocas operaciones sobre usuarios;
  \item están pensadas para el mercado anglosajón: interfaz, correos,
        comentarios y documentación en inglés; y
  \item su calidad de seguridad no está garantizada. Al ser código ajeno,
        quien lo adopta asume también sus fallos, y la revisión corre de su
        cuenta.
\end{enumerate}

\subsection{Servicios de \textit{backend} gestionado (BaaS)}
\label{sec:ctx-baas}

Los servicios de \textit{Backend as a Service} ofrecen ya construidos la base
de datos, la autenticación, el almacenamiento de archivos y una API de acceso a
los datos. La tabla~\ref{tab:ctx-baas} compara los más extendidos en sus
características relevantes para este trabajo.

\begin{table}[htbp]
  \centering
  \small
  \caption{Comparación de servicios de \textit{backend} gestionado.}
  \label{tab:ctx-baas}
  \begin{tabularx}{\textwidth}{@{}l>{\raggedright\arraybackslash}X>{\raggedright\arraybackslash}X>{\raggedright\arraybackslash}p{2.6cm}@{}}
    \toprule
    \textbf{Servicio} & \textbf{Modelo de datos} & \textbf{Autorización} & \textbf{Código / alojamiento} \\
    \midrule
    Firebase~\cite{firebase-docs} & Documental (Cloud Firestore) & Reglas de seguridad propias, declaradas aparte de los datos & Propietario; solo en la nube del proveedor \\
    Supabase~\cite{supabase-docs} & Relacional (PostgreSQL) con API generada & Políticas RLS de PostgreSQL sobre la identidad del JWT~\cite{supabase-rls} & Abierto; nube o autoalojado \\
    Appwrite~\cite{appwrite-docs} & Colecciones de documentos sobre una base relacional interna & Permisos por colección y documento & Abierto; nube o autoalojado \\
    PocketBase~\cite{pocketbase} & SQLite embebido & Reglas de API por colección & Abierto; un único ejecutable autoalojado \\
    \bottomrule
  \end{tabularx}
\end{table}

Para una aplicación multicliente con relaciones entre cuentas, miembros,
roles, suscripciones y datos de negocio, un modelo relacional con integridad
referencial es una ventaja clara. Supabase añade otra: su autorización no es
un sistema paralelo, sino la propia seguridad a nivel de fila de
PostgreSQL~\cite{postgresql-rls}. Cada petición llega a la base de datos con
la identidad del usuario, y las políticas deciden qué filas puede ver o
modificar. Al ser PostgreSQL estándar, además, el esquema, las políticas y sus
pruebas pueden versionarse como código y probarse con herramientas como
pgTAP~\cite{pgtap}. Su contrapartida es que la seguridad depende de que esas
políticas estén bien escritas, y un error en ellas no se ve en la interfaz.

\subsection{Gestores de contenidos y paneles de administración de datos}
\label{sec:ctx-cms}

El componente de administración de datos admite varias soluciones:

\begin{itemize}
  \item \textbf{CMS \textit{headless}} como Strapi~\cite{strapi-docs}. El
        desarrollador define tipos de contenido y el CMS genera una API y un
        panel para editarlos. Son excelentes para contenidos editoriales,
        pero gestionan su propio esquema y su propio sistema de usuarios: no
        están pensados para administrar las tablas de negocio de otra
        aplicación ni para reutilizar sus políticas de acceso.
  \item \textbf{Plataformas de datos sobre una base existente} como
        Directus~\cite{directus-docs}, que se conectan a una base de datos SQL
        ya creada y ofrecen sobre ella una API y una interfaz de edición con
        roles y permisos. Se acercan más a lo que se busca, pero añaden un
        segundo servicio de \textit{backend} con su propio modelo de usuarios
        y permisos que hay que sincronizar con el de la aplicación.
  \item \textbf{Herramientas para desarrolladores}, como el editor de tablas
        del panel de Supabase~\cite{supabase-docs} o pgAdmin~\cite{pgadmin}.
        Son potentes, pero trabajan con privilegios de administrador de la
        base de datos, sin permisos por tabla para el personal y sin
        auditoría orientada al negocio. No se pueden entregar a quien gestiona
        una pyme.
  \item \textbf{CMS específicos para un BaaS}, que se instalan junto a un
        proyecto de Supabase y ofrecen un explorador de datos, de usuarios y de
        almacenamiento con un RBAC propio guardado en la misma base de datos.
        Son la opción más cercana al objetivo, aunque normalmente se
        distribuyen como una aplicación independiente, con su propio inicio de
        sesión y su propia pila tecnológica.
\end{itemize}

\subsection{Aislamiento multicliente en la base de datos}
\label{sec:ctx-multitenancy}

Existen tres estrategias clásicas para aislar los datos de los clientes de una
aplicación SaaS~\cite{chong2006multitenant}: una base de datos por cliente; un
esquema por cliente dentro de una base compartida; y tablas compartidas en las
que cada fila lleva el identificador del cliente al que pertenece. La primera
aísla mejor, pero multiplica el coste de operación; la última es la más
económica y la que mejor escala en número de clientes, a cambio de que el
aislamiento dependa por completo de que \emph{todas} las consultas filtren por
cliente~\cite{krebs2012multitenancy}.

La seguridad a nivel de fila resuelve esa debilidad trasladando el filtro de
la aplicación a la base de datos~\cite{postgresql-rls,postgresql-createpolicy}.
Cada tabla declara políticas del tipo «un usuario solo puede leer las filas de
las cuentas de las que es miembro», y PostgreSQL las aplica a cualquier
consulta, venga de donde venga. Un fallo en el código de la aplicación ya no
basta para filtrar datos de otro cliente: tendría que fallar también la
política. Es el enfoque que recomienda Supabase~\cite{supabase-rls} y el que
adopta este trabajo, complementado con pruebas automáticas que atacan las
políticas suplantando a usuarios de distintos clientes.

\section{Justificación de la solución}
\label{sec:ctx-justificacion}

Del análisis anterior se desprenden las decisiones que dan forma a PymeKit:

\begin{enumerate}
  \item \textbf{Una base reutilizable, no una aplicación.} Lo que se repite
        de un SaaS a otro es la infraestructura de la
        tabla~\ref{tab:ctx-componentes}. Construirla una vez, bien probada y
        configurable, es lo que acelera los proyectos siguientes (RNF-01).
  \item \textbf{PostgreSQL con RLS como frontera de seguridad.} Un modelo
        relacional encaja con los datos de un SaaS multicliente, y la RLS hace
        que el aislamiento entre clientes no dependa solo del código. Se
        adopta Supabase porque ofrece autenticación, almacenamiento y
        PostgreSQL estándar, es de código abierto y permite un entorno local
        completo para el desarrollo y las pruebas.
  \item \textbf{Un CMS de datos dentro de la propia aplicación.} Ni los CMS
        \textit{headless} ni las herramientas para desarrolladores cubren la
        necesidad de que el personal de una pyme administre los datos de
        \emph{su} aplicación con permisos y auditoría. Un CMS que trabaja
        sobre la misma base de datos sí lo hace, y si además se integra en la
        consola de administración se evita un segundo servicio, un segundo
        inicio de sesión y una segunda pila tecnológica (ADR-011).
  \item \textbf{Español e internacionalización.} Una base dirigida a pymes
        españolas debe estar en español, pero sin cerrar la puerta a otros
        idiomas (ADR-021).
\end{enumerate}

\paragraph{Partir de componentes base.}
Como decisión de ingeniería, PymeKit se construye a partir de componentes base
con licencia comercial ---una base de aplicación SaaS y un CMS para la base de
datos---, una práctica habitual en la industria y coherente con el propio
objetivo de reutilización de este trabajo (ADR-001). La aportación del trabajo
está en la arquitectura, la integración en una sola aplicación, la adaptación
a pymes españolas, la verificación de su seguridad y su reutilizabilidad. Que una base reutilizable deba verificarse
antes de reutilizarse es, en sí mismo, uno de los resultados que se analizan en
el capítulo de conclusiones.


\newpage{\ }
\cleardoublepage

\chapter{Plan de gestión de proyecto}
\label{cap:planificación}
% ------------------------------------------------------------------------------
% Capítulo 3 · Plan de gestión del proyecto
% Fuentes: docs/tfg/{PLAN,PROGRESO,BITACORA,DECISIONES}.md, AGENTS.md,
% .claude/, .github/workflows/ci.yml, git log (54 commits, 29/09–20/10/2026)
% y el borrador memoria/borradores/02-planificacion.tex.
% Fecha de corte de los datos: 20/10/2026.
% ------------------------------------------------------------------------------
\label{cap:plan}

Este capítulo recoge la gestión del proyecto: la metodología seguida, las
tecnologías elegidas, la planificación y su comparación con lo realmente
ejecutado, la organización, los costes, los riesgos, la gestión de la
configuración y el aseguramiento de la calidad.

% ---------------------------------------------------------------------
\section{Metodología de desarrollo}
\label{sec:plan-metodologia}

\subsection{Desarrollo incremental por fases}
\label{sec:plan-incremental}

El proyecto sigue un modelo de desarrollo \textbf{iterativo e
incremental}~\cite{larman2003iterative}. El trabajo se divide en nueve fases
(F0--F8) y, dentro de las fases grandes, en incrementos más pequeños (por
ejemplo, F2.1--F2.9 en la integración del CMS). Cada fase o incremento tiene,
fijados de antemano en el plan del proyecto:

\begin{itemize}
  \item un \textbf{objetivo} y una lista de entregables;
  \item un \textbf{criterio de «hecho»} verificable (por ejemplo, «la
        aplicación arranca en local, CU-01 funciona a mano y la integración
        continua está en verde»); y
  \item sus \textbf{dependencias} con otras fases.
\end{itemize}

Un incremento solo se cierra cuando supera su criterio de «hecho» y la cadena
de verificación común descrita en la sección~\ref{sec:plan-calidad}. Así, al
final de cada incremento el sistema queda en un estado funcional y comprobado,
y un retraso afecta al alcance de los incrementos siguientes, no a la
estabilidad de lo ya construido. Esta forma de trabajar encaja con el objetivo
de una arquitectura reutilizable: cada incremento añade un módulo que debe
integrarse sin romper los anteriores (RNF-04, RNF-06).

No se ha adoptado un marco ágil formal como Scrum: con un único desarrollador
no tienen sentido sus roles ni sus ceremonias. Sí se han tomado de los métodos
ágiles la entrega de incrementos funcionales cortos, la definición explícita
de «hecho», la integración continua y la revisión frecuente con el cliente,
papel que aquí desempeña el propio autor como responsable del producto: varias
decisiones de alcance (la integración del CMS en la consola, la consola por
áreas o la interfaz solo en español) surgieron de su revisión de un incremento
ya terminado.

Las decisiones de arquitectura o de alcance se registran en el momento en que
se toman como \textit{Architecture Decision Records}~\cite{nygard2011adr},
con su contexto, las alternativas consideradas y sus consecuencias. A la fecha
de corte hay veintiuno (ADR-001 a ADR-021); los relevantes para el diseño se
resumen en el capítulo de Diseño del Sistema. Las incidencias ---fallos
de código, problemas de entorno y de herramientas--- se anotan
también en el momento, en una bitácora con el formato \emph{qué pasó, causa,
solución, evidencia y lección}. La bitácora es la fuente de los riesgos
materializados (sección~\ref{sec:plan-riesgos}) y de las lecciones aprendidas
del capítulo~\ref{cap:conclusiones}.

\subsection{Desarrollo asistido por agentes de IA}
\label{sec:plan-ia}

% PENDIENTE (autor): comprobar si la ESI/UCA exige una declaración formal del
% uso de herramientas de IA generativa en el TFG (formato y ubicación).

Una parte sustancial del código, de las pruebas y de la documentación se ha
producido con la ayuda de un asistente de programación basado en modelos de
lenguaje que actúa como \textit{agente}: un programa que, a partir de una
instrucción, lee ficheros, ejecuta órdenes y propone o aplica cambios en el
repositorio. Se ha usado Claude Code~\cite{claudecode}. Los
\textit{commits} del repositorio registran esa coautoría en su mensaje.
\todo{Confirmar el modelo o modelos de IA empleados y el periodo de uso para
la declaración de uso de IA, si la normativa de la ESI la exige.}

Para que este modo de trabajo sea controlable y evaluable, antes de escribir
código de la aplicación (fase F0) se construyó un \textit{harness} o arnés de
directrices y controles automáticos que acota lo que hacen los agentes y
comprueba el resultado. Es una aportación propia del trabajo y tiene cinco
componentes:

\begin{itemize}
  \item \textbf{Directrices escritas.} Un fichero raíz de directrices para
        agentes y desarrolladores y otros por paquete, que fijan la pila, las
        convenciones, las reglas de seguridad de la base de datos, el idioma
        de los comentarios y de la documentación y una lista de verificación
        obligatoria al terminar cualquier cambio.
  \item \textbf{Documentación de seguimiento}: plan, requisitos, decisiones,
        trazabilidad, progreso con la dedicación, y bitácora. Un agente debe
        leer el progreso y el plan de la fase antes de empezar, y no puede
        reabrir una decisión sin registrar un ADR nuevo.
  \item \textbf{Procedimientos reutilizables} (\textit{skills}) para tareas
        repetitivas o delicadas: incorporar un módulo,
        reescribir sus comentarios, registrar una decisión, revisar un cambio
        de forma adversarial, auditar las políticas de seguridad de la base de
        datos o redactar un borrador de la memoria.
  \item \textbf{Controles automáticos} (\textit{hooks}): al inicio de cada
        sesión se inyectan la fase actual y sus tareas pendientes; tras cada
        edición se ejecutan sobre el fichero los controles automáticos del
        repositorio; y antes de terminar se recuerda actualizar el
        progreso si hay cambios sin registrar.
  \item \textbf{Permisos explícitos}: los agentes no pueden escribir fuera del
        repositorio, leer ficheros de secretos ni
        hacer \textit{commit} o \textit{push} sin orden expresa del autor.
\end{itemize}

\paragraph{Reparto de responsabilidades.}
\textbf{El autor} fija el alcance y los requisitos, toma las decisiones de
arquitectura ---la integración del CMS en la consola (ADR-011) o la
organización por áreas (ADR-020) fueron propuestas suyas---, revisa y acepta o
rechaza los cambios, hace pruebas manuales exploratorias y es el único que
autoriza los \textit{commits}. \textbf{Los agentes} analizan el código
existente, proponen alternativas, implementan los incrementos, escriben las
pruebas, ejecutan la verificación, redactan la documentación de seguimiento y
los borradores de la memoria, y hacen las revisiones adversariales.
\todo{El autor completa este reparto con su propia valoración: qué partes ha
escrito o reescrito a mano, qué criterio ha seguido para aceptar cambios y
cuánto tiempo ha dedicado a la revisión.}

La revisión humana sigue siendo imprescindible. La incidencia B-14 (la consola
se cargaba sin segundo factor) la detectó el autor con una prueba manual que
ninguna prueba automática cubría, y la B-11 muestra que un test en verde no
garantiza la corrección: un test existente afirmaba precisamente el
comportamiento inseguro.

\paragraph{Modo de ahorro.}
La verificación exhaustiva con varios agentes en paralelo encontró brechas
reales, pero su coste en uso de la herramienta resultó excesivo (B-48). Desde
la F2.7c se trabaja en un \emph{modo de ahorro}: un único agente por paso,
pruebas E2E acotadas a la zona tocada y la suite completa solo al cerrar una
fase, antes de fusionar. El nivel de revisión se decide por riesgo y no por
defecto; las incidencias B-51 y B-54 muestran el precio de ese ajuste y
justifican ejecutar siempre la suite completa antes de abrir una
\textit{pull request}.

% ---------------------------------------------------------------------
\section{Tecnologías}
\label{sec:plan-tecnologias}

La pila tecnológica se elige por usar un único lenguaje tipado (TypeScript) en
el cliente, el servidor y el acceso a la base de datos, por apoyarse en
Supabase para los datos, la autenticación y el almacenamiento con la
autorización aplicada en la propia base de datos, y por emplear una sola pila
de librerías en toda la web (ADR-013). La tabla~\ref{tab:plan-tecnologias} resume cada
tecnología y su justificación.

\begin{table}[htbp]
  \centering
  \footnotesize
  \caption{Tecnologías empleadas y justificación.}
  \label{tab:plan-tecnologias}
  \begin{tabularx}{\textwidth}{@{}>{\raggedright\arraybackslash}p{3.3cm}>{\raggedright\arraybackslash}p{3.1cm}X@{}}
    \toprule
    \textbf{Tecnología} & \textbf{Uso} & \textbf{Justificación} \\
    \midrule
    TypeScript~\cite{typescript} (modo estricto) & Lenguaje de todo el código & Tipos compartidos entre cliente, servidor y base de datos; errores detectados en compilación (RNF-05). \\
    React~\cite{react} y TanStack Start~\cite{tanstack-start} & Aplicación web (\textit{frontend} y \textit{server functions}) & Renderizado en servidor, rutas por fichero con TanStack Router~\cite{tanstack-router} y funciones de servidor tipadas y validadas (RF-01, RF-02). \\
    TanStack Query~\cite{tanstack-query} & Caché de datos en el cliente & Carga en servidor y caché de consultas (RNF-09). \\
    Supabase~\cite{supabase-docs} (PostgreSQL, Auth, Storage) & Base de datos, autenticación y archivos & PostgreSQL estándar con RLS como frontera de seguridad~\cite{supabase-rls}; MFA con niveles de garantía~\cite{supabase-mfa}; entorno local completo. \\
    Hono~\cite{hono} y Drizzle ORM~\cite{drizzle} & API del CMS en \texttt{/api/cms} & API ligera sobre estándares web, montable dentro del servidor de la aplicación; consultas tipadas ejecutadas con la identidad del usuario. \\
    Stripe~\cite{stripe-docs} & Suscripciones y pagos & Pasarela con modo de pruebas completo, \textit{checkout}, portal del cliente y \textit{webhooks} (RF-07). \\
    Tailwind CSS~\cite{tailwind} y shadcn/ui~\cite{shadcn} & Estilos y componentes & Componentes accesibles, copiados al repositorio y adaptables; tema claro y oscuro (RNF-08). \\
    TanStack Form y Zod & Formularios y validación & Una sola forma de construir formularios en toda la web, incluido el CMS (ADR-013). \\
    use-intl & Internacionalización & Todos los textos en catálogos de mensajes; base del punto de extensión de idiomas (ADR-021). \\
    Turborepo~\cite{turborepo} y pnpm & Monorepo & Paquetes desacoplados con dependencias explícitas y caché de tareas (RNF-04). \\
    Vitest~\cite{vitest}, pgTAP~\cite{pgtap} y Playwright~\cite{playwright} & Pruebas & Unitarias, de base de datos (incluidas pruebas de ataque a las políticas RLS) y de extremo a extremo (RNF-06). \\
    Docker~\cite{docker-docs} & Entorno local & Ejecuta la instancia local de Supabase. \\
    Git, GitHub y GitHub Actions~\cite{github-actions} & Versiones e integración continua & Repositorio privado y CI con controles del TFG, tipos, \textit{lint}, pruebas unitarias, pgTAP y E2E. \\
    Claude Code~\cite{claudecode} & Asistente de programación & Desarrollo asistido por agentes con el arnés de la sección~\ref{sec:plan-ia}. \\
    \LaTeX{} & Memoria & Plantilla oficial de la ESI. \\
    \bottomrule
  \end{tabularx}
\end{table}

% ---------------------------------------------------------------------
\section{Planificación del proyecto}
\label{sec:plan-planificacion}

\subsection{Fases y relación con la EDT}
\label{sec:plan-fases}

La propuesta define una EDT con seis paquetes de trabajo: (1)~Análisis,
(2)~Arquitectura, (3)~Modelo de datos, (4)~Implementación, (5)~Validación y
(6)~Documentación. La tabla~\ref{tab:plan-fases} relaciona con ellos cada fase
del proyecto y su criterio de «hecho».

\begin{table}[htbp]
  \centering
  \small
  \caption{Fases del proyecto, paquetes de la EDT que cubren y criterio de «hecho».}
  \label{tab:plan-fases}
  \begin{tabularx}{\textwidth}{@{}l>{\raggedright\arraybackslash}p{3.4cm}cX@{}}
    \toprule
    \textbf{Fase} & \textbf{Nombre} & \textbf{EDT} & \textbf{Hecho cuando\ldots} \\
    \midrule
    F0 & Harness y planificación & 1 & Los controles automáticos pasan, los \textit{hooks} funcionan y existe el primer \textit{commit} \\
    F1 & Base de la aplicación web funcionando & 4 & La aplicación arranca en local, CU-01 funciona a mano y la CI está en verde \\
    F2 & Integración del CMS & 2, 3, 4 & Cada incremento F2.1--F2.9 cumple su criterio (BD, API, interfaz, explorador, usuarios y almacenamiento, auditoría, ajustes y RBAC, paneles, cierre) \\
    F3 & Interfaz en español e i18n & 4 & Interfaz en español por defecto y controles automáticos sin excepciones \\
    F4 & Comentarios didácticos en español & 4, 6 & Sin comentarios en inglés en los módulos críticos y revisión aprobada \\
    F5 & Adaptación a pymes y configurabilidad & 2, 3, 4 & CU-02 a CU-05 funcionan y la configuración está documentada \\
    F6 & Validación & 5 & CI en verde y resultados documentados \\
    F7 & Despliegue & 4, 6 & Despliegue de demostración accesible y manual probado desde cero \\
    F8 & Memoria y manuales & 1, 2, 3, 6 & Continua; se cierra tras F6 y F7 \\
    \bottomrule
  \end{tabularx}
\end{table}

Las dependencias principales son lineales (F0 $\rightarrow$ F1 $\rightarrow$
F2 $\rightarrow$ F3 $\rightarrow$ F5 $\rightarrow$ F6), con dos excepciones: F4
puede solaparse con F3 y F5, y F8 se desarrolla de forma continua, de modo que
cada fase cerrada deja un borrador de su parte de la memoria.

La fase F2 es la de mayor riesgo. Al comenzarla se evaluaron tres formas de
integrar el CMS ---aplicaciones separadas, integración híbrida e integración
total--- y se eligió la integración total en la consola de administración
(ADR-011), a pesar de ser la de mayor coste estimado, por dar lugar a un único
servicio, un único inicio de sesión y un único estilo de código. Para
controlar ese riesgo, la fase se dividió en incrementos y se identificó de
antemano el candidato a recortar si el plazo lo exigía: los paneles
configurables (RF-11, requisito deseable). Finalmente no hizo falta
recortarlo.

\subsection{Calendario real}
\label{sec:plan-calendario}

La figura~\ref{fig:plan-gantt} muestra el diagrama de Gantt de lo ejecutado
hasta la fecha de corte (20/10/2026), reconstruido a partir de las fechas de
los \textit{commits} y del registro de progreso. La
tabla~\ref{tab:plan-hitos} recoge los hitos.

\begin{figure}[htbp]
  \centering
  \resizebox{\textwidth}{!}{%
  \begin{ganttchart}[
      time slot format=isodate,
      x unit=0.55cm,
      y unit chart=0.62cm,
      y unit title=0.55cm,
      vgrid={*1{gray!25, dotted}},
      title label font=\footnotesize,
      bar label font=\footnotesize,
      group label font=\footnotesize\bfseries,
      milestone label font=\footnotesize\itshape,
      bar/.append style={fill=blue!45},
      group/.append style={fill=gray!60},
      milestone/.append style={fill=red!70},
      bar height=0.55,
      group peaks height=0.2,
    ]{2026-09-29}{2026-10-20}
    \gantttitle{sep.}{2}
    \gantttitle{octubre de 2026}{20} \\
    \gantttitlecalendar{day} \\
    \ganttbar{F0 · Harness y plan}{2026-09-29}{2026-09-29} \\
    \ganttbar{F1 · Base de la web}{2026-09-29}{2026-09-29} \\
    \ganttmilestone{Fusión de F1 en \texttt{main}}{2026-09-29} \\
    \ganttgroup{F2 · Integración del CMS}{2026-09-29}{2026-10-02} \\
    \ganttbar{F2.1--F2.4 · BD, API, interfaz, explorador}{2026-09-29}{2026-09-29} \\
    \ganttbar{F2.5--F2.7b · usuarios, archivos, auditoría, ajustes}{2026-09-30}{2026-09-30} \\
    \ganttbar{F2.7c--F2.9 · recursos, paneles, cierre}{2026-10-02}{2026-10-02} \\
    \ganttmilestone{Fusión de F2}{2026-10-02} \\
    \ganttgroup{F3 · Español y consola por áreas}{2026-09-30}{2026-10-05} \\
    \ganttbar{Identidad visual (adelantada)}{2026-09-30}{2026-09-30} \\
    \ganttbar{F3a--F3c · español, consola por áreas}{2026-10-02}{2026-10-05} \\
    \ganttmilestone{Fusión de F3 (PR \#2)}{2026-10-02} \\
    \ganttbar{Bloque académico en la landing}{2026-10-20}{2026-10-20} \\
    \ganttbar[bar/.append style={fill=green!40}]{F8 · Memoria (continua)}{2026-09-29}{2026-10-20}
  \end{ganttchart}%
  }
  \caption[Diagrama de Gantt de la ejecución real]{Diagrama de Gantt de la
    ejecución real hasta el 20/10/2026. Las barras abarcan días naturales; el
    trabajo efectivo se concentró en los días 29/09, 30/09, 02/10, 05/10 y
    20/10. Las fases F4--F7 no se habían iniciado en la fecha de corte.}
  \label{fig:plan-gantt}
\end{figure}

\begin{table}[htbp]
  \centering
  \small
  \caption{Hitos del proyecto según el historial del repositorio.}
  \label{tab:plan-hitos}
  \begin{tabularx}{\textwidth}{@{}llX@{}}
    \toprule
    \textbf{Fecha} & \textbf{\textit{Commit}} & \textbf{Hito} \\
    \midrule
    29/09/2026 & \texttt{0053dd7} & Primer \textit{commit}: directrices, plan y arnés del TFG (cierre de F0) \\
    29/09/2026 & \texttt{d9ff441} & Base de la aplicación web operativa y fusionada en \texttt{main} (cierre de F1) \\
    29/09/2026 & \texttt{48ddd1c} & Decisión de integrar el CMS en la consola (ADR-011 a ADR-014) \\
    29/09/2026 & \texttt{fd9a16a} & Primeros borradores de la memoria (planificación, análisis y diseño) \\
    02/10/2026 & \texttt{8d74ca3} & CMS integrado en la consola, fusionado en \texttt{main} (cierre de F2) \\
    02/10/2026 & \texttt{dc75f6e} & Español por defecto, fusionado (cierre de F3a--F3b) \\
    02/10/2026 & \texttt{ab1ee4c} & Interfaz solo en español (ADR-021) \\
    05/10/2026 & \texttt{6899a70} & Resumen en cada pantalla de la consola \\
    20/10/2026 & \texttt{75ff3fa} & Bloque académico y recorrido de la demostración en la landing \\
    \bottomrule
  \end{tabularx}
\end{table}

\todo{Completar el calendario con las fases F4--F8 cuando se ejecuten (fechas
de inicio y de cierre) y actualizar la figura~\ref{fig:plan-gantt}.}

\subsection{Dedicación registrada}
\label{sec:plan-dedicacion}

La dedicación se registra en el documento de progreso del proyecto con dos
medidas: las \textbf{horas del autor}, que son su dedicación real declarada y
las que cuentan para la planificación y los costes; y el \textbf{tiempo de
reloj}, estimado a partir de la hora de los \textit{commits}, que incluye
periodos en los que los agentes trabajan mientras el autor no está dedicado a
la tarea y es, por tanto, una cota superior. La
tabla~\ref{tab:plan-dedicacion} recoge los datos disponibles.

\begin{table}[htbp]
  \centering
  \small
  \caption{Dedicación registrada por jornada.}
  \label{tab:plan-dedicacion}
  \begin{tabularx}{\textwidth}{@{}lrrX@{}}
    \toprule
    \textbf{Fecha} & \textbf{Autor (h)} & \textbf{Reloj (h)} & \textbf{Fases trabajadas} \\
    \midrule
    29/09/2026 & 8,0 & $\approx$ 8,0 & F0, F1, F2.1--F2.4c, bitácora y borradores de la memoria \\
    30/09/2026 & 9,0 & --- & F2.5, F2.6, F2.6b, F2.7a, F2.7b, identidad visual (F3 adelantada) \\
    02/10/2026 & 6,0 & --- & F2.7c, F2.8, F2.9 (cierre de F2), F3a, F3b y F3c \\
    05/10/2026 & \textit{pendiente: h} & --- & Resumen en cada pestaña de la consola \\
    20/10/2026 & \textit{pendiente: h} & --- & Bloque académico de la landing y redacción de la memoria \\
    \midrule
    \textbf{Total registrado} & \textbf{23,0} & & F0--F3 completas \\
    \bottomrule
  \end{tabularx}
\end{table}

El 29/09 se desglosó además por tramos a partir de los \textit{commits}: F0
ocupó unos 45 minutos de reloj, F1 otros 45 y la integración del CMS
(F2.1--F2.4c) las seis horas y media restantes.

\todo{Registrar la dedicación de las jornadas posteriores al 02/10 y de las
fases F4--F8, e incluir el tiempo que no deja \textit{commits}: revisión,
pruebas manuales, reuniones con el tutor y redacción de la memoria.}

\subsection{Comparación entre lo estimado y lo real}
\label{sec:plan-estimado-real}

El plan del proyecto fijó desde el principio el orden de las fases, sus
entregables y sus criterios de «hecho», pero no una estimación en horas de
cada una. La única estimación cuantitativa registrada es la del análisis de
alternativas para integrar el CMS (ADR-011), hecha al comenzar la F2 en
términos de esfuerzo de desarrollo convencional. La
tabla~\ref{tab:plan-estimado-real} la compara con lo ocurrido.

\begin{table}[htbp]
  \centering
  \small
  \caption{Estimación registrada de la integración del CMS frente a la ejecución real.}
  \label{tab:plan-estimado-real}
  \begin{tabularx}{\textwidth}{@{}>{\raggedright\arraybackslash}p{4.2cm}>{\raggedright\arraybackslash}p{3cm}X@{}}
    \toprule
    \textbf{Alternativa (ADR-011)} & \textbf{Estimación} & \textbf{Real} \\
    \midrule
    (a) Aplicaciones separadas & 3--6 días & No ejecutada \\
    (c) Integración híbrida & 1--2 semanas & No ejecutada \\
    (b) Integración total (elegida) & 6--12 o más semanas & Cuatro días naturales (29/09 a 02/10) con tres jornadas de trabajo, que en total suman 23 h del autor e incluyen también F0, F1 y la mayor parte de F3 \\
    \bottomrule
  \end{tabularx}
\end{table}

La diferencia es de un orden de magnitud y requiere una lectura cuidadosa.
Primero, la estimación suponía un desarrollo convencional, mientras que la
ejecución fue asistida por agentes que trabajaban en paralelo y de forma
continuada durante cada jornada. Segundo, las horas del autor miden su
dedicación, no el trabajo total realizado por los agentes, por lo que no son
comparables uno a uno con una estimación de esfuerzo humano. Tercero, la
integración no está completa en todos sus detalles: el registro de progreso
mantiene una lista de pendientes menores de cada incremento, y la fase de
comentarios didácticos (F4), que afecta también al código del CMS, sigue
abierta. Aun con esas salvedades, el dato indica que el desarrollo asistido
desplaza el cuello de botella desde la escritura del código hacia la
verificación y la revisión, que es donde se concentraron las incidencias de la
bitácora. Este efecto se valora en el capítulo~\ref{cap:conclusiones}.

\todo{Añadir la carga prevista del TFG según el plan de estudios (créditos
ECTS y su equivalencia en horas, entre 25 y 30 h por crédito) y, si el autor
dispone de ella, la planificación inicial por fases (fechas o semanas
previstas) para completar la comparación.}

% ---------------------------------------------------------------------
\section{Organización}
\label{sec:plan-organizacion}

\subsection{Personas y roles}

El proyecto lo ejecuta una sola persona, que asume varios roles, con la
supervisión del tutor académico (tabla~\ref{tab:plan-roles}).

\begin{table}[htbp]
  \centering
  \small
  \caption{Roles del proyecto.}
  \label{tab:plan-roles}
  \begin{tabularx}{\textwidth}{@{}>{\raggedright\arraybackslash}p{3.8cm}>{\raggedright\arraybackslash}p{3.2cm}X@{}}
    \toprule
    \textbf{Rol} & \textbf{Quién} & \textbf{Responsabilidades} \\
    \midrule
    Director del proyecto y responsable del producto & Autor & Alcance, prioridades, decisiones de arquitectura (ADR) y aceptación de cada incremento \\
    Analista y arquitecto & Autor (asistido) & Requisitos, casos de uso, arquitectura, modelo de datos \\
    Desarrollador y probador & Autor y agentes de IA & Implementación, pruebas unitarias, de base de datos y E2E, revisión adversarial \\
    Redactor & Autor (asistido) & Documentación técnica, manuales y memoria \\
    Tutor académico & Juan Carlos de la Torre Macías & Supervisión, validación del alcance (punto abierto P-01) y revisión de la memoria \\
    \bottomrule
  \end{tabularx}
\end{table}

\subsection{Recursos materiales}

\begin{itemize}
  \item \textbf{Servidor de desarrollo remoto} con Linux y Docker, sobre el
        que se ejecutan la instancia local de Supabase, la aplicación y los
        agentes. Se accede a él mediante túnel SSH; sus servicios locales están
        aislados de Internet (ADR-010).
        \todo{Indicar las características del servidor (CPU, memoria,
        proveedor) y del equipo personal del autor.}
  \item \textbf{Servicios en la nube}: GitHub (repositorio privado e
        integración continua), Stripe en modo de pruebas y, para el
        despliegue, una plataforma de aplicaciones y un proyecto de Supabase
        en la nube (véase el capítulo de Despliegue del Sistema).
  \item \textbf{Software}: Node.js, pnpm, Turborepo, Supabase CLI, Docker,
        Vitest, pgTAP, Playwright y \LaTeX{}, todos de código abierto, y el
        asistente de programación con IA.
\end{itemize}

% ---------------------------------------------------------------------
\section{Costes}
\label{sec:plan-costes}

El coste del proyecto se estima como la suma del coste de personal y del de
los recursos materiales. El de personal se obtiene multiplicando las horas
dedicadas por el coste por hora de un ingeniero informático de perfil
\textit{junior} en España, incluidas las cotizaciones a cargo de la empresa. El
de los recursos materiales incluye la amortización del equipo, los servicios
contratados durante el proyecto y las licencias. La
tabla~\ref{tab:plan-costes} recoge la estructura del presupuesto.

\begin{table}[htbp]
  \centering
  \small
  \caption{Presupuesto del proyecto.}
  \label{tab:plan-costes}
  \begin{tabularx}{\textwidth}{@{}X>{\raggedleft\arraybackslash}p{3.8cm}r@{}}
    \toprule
    \textbf{Concepto} & \textbf{Base de cálculo} & \textbf{Importe (€)} \\
    \midrule
    Personal: horas del autor $\times$ coste por hora & \textit{pendiente: horas totales y tarifa} & \textit{pendiente} \\
    Equipo de desarrollo (amortización) & \textit{pendiente: valor y meses} & \textit{pendiente} \\
    Servidor de desarrollo remoto & \textit{pendiente: meses} & \textit{pendiente} \\
    Licencias de componentes de software & \textit{pendiente: importe} & \textit{pendiente} \\
    Asistente de programación con IA & \textit{pendiente: meses y plan} & \textit{pendiente} \\
    Servicios para el despliegue (plataforma y base de datos en la nube) & \textit{pendiente: plan} & \textit{pendiente} \\
    Otros (dominio, etc.) & \textit{pendiente} & \textit{pendiente} \\
    \midrule
    \textbf{Total} & & \textit{pendiente} \\
    \bottomrule
  \end{tabularx}
\end{table}

Las herramientas de desarrollo empleadas son de código abierto y no tienen
coste de licencia, y Stripe se usa en modo de pruebas, sin coste. Conviene
subrayar que, en un escenario profesional, el ahorro que justifica una base
reutilizable no se aprecia en este presupuesto sino en el de los proyectos
siguientes: cada SaaS que se construye sobre PymeKit evita el coste de
desarrollar y verificar de nuevo sus componentes comunes (véase la evaluación
de la reutilización en el anexo~\ref{cap:evaluacion}).
\todo{Completar los importes de la tabla~\ref{tab:plan-costes}, indicar la fuente de la tarifa horaria (convenio o estudio salarial) y
confirmar los planes de GitHub, Supabase y la plataforma de despliegue.}

% ---------------------------------------------------------------------
\section{Riesgos}
\label{sec:plan-riesgos}

La tabla~\ref{tab:plan-riesgos} recoge los riesgos identificados, su
mitigación y si se han materializado, con las incidencias de la bitácora que
lo documentan (identificadores B-xx).

\begin{table}[htbp]
  \centering
  \footnotesize
  \caption{Riesgos del proyecto, mitigación y materialización.}
  \label{tab:plan-riesgos}
  \begin{tabularx}{\textwidth}{@{}l>{\raggedright\arraybackslash}p{3.5cm}>{\raggedright\arraybackslash}X>{\raggedright\arraybackslash}p{2.3cm}@{}}
    \toprule
    \textbf{ID} & \textbf{Riesgo} & \textbf{Mitigación} & \textbf{Materializado} \\
    \midrule
    R-01 & Desbordamiento del plazo por la integración del CMS &
      Incrementos funcionales F2.1--F2.9; RF-11 identificado como primer recorte (ADR-011) & No \\
    R-02 & Fallos de seguridad en el control de acceso &
      Auditoría adversarial de RLS con pruebas pgTAP de ataque y refutación independiente; pruebas de regresión (ADR-015) & Sí: B-05--B-09, B-19--B-25, B-27, B-28, B-42, B-47, B-50 \\
    R-03 & Exposición del entorno de desarrollo (servidor con IP pública) &
      Filtrado en la cadena \texttt{DOCKER-USER}, acceso por túnel SSH y revisión de cada servicio nuevo (ADR-010) & Sí: B-04 \\
    R-04 & Errores introducidos por cambios masivos o por los agentes &
      Controles automáticos que cubren las variantes del error, revisión adversarial y prueba manual del autor & Sí: B-01, B-14, B-16, B-45 \\
    R-05 & Deriva entre el esquema declarativo y las migraciones &
      Diferencia de esquemas en la verificación de todo cambio de BD; propiedades de seguridad comprobadas con pgTAP & Sí: B-10, B-35 \\
    R-06 & Pruebas que no prueban lo que dicen o que quedan desactualizadas &
      Leer las aserciones de las pruebas existentes; suite completa antes de fusionar & Sí: B-11, B-51, B-54 \\
    R-07 & Interrupción de las herramientas o de la conexión durante una tarea &
      Cerrar cada incremento con su \textit{commit}; tareas de seguridad formuladas como verificación defensiva y por etapas & Sí: B-12, B-15 \\
    R-08 & Falsos positivos de las herramientas (secretos, pruebas inestables) &
      Documentar las credenciales públicas de desarrollo; E2E contra la \textit{build} de test & Sí: B-03, B-13 \\
    R-09 & Coste excesivo de la verificación asistida por agentes &
      Modo de ahorro: un agente por paso y suite completa al cerrar la fase & Sí: B-48 \\
    R-10 & Incumplimiento de las licencias de las dependencias &
      Repositorio privado y revisión de las condiciones de licencia (ADR-001, ADR-006) & No \\
    R-11 & Validación de la reutilización sin un dominio de negocio concreto &
      Medir el arranque de un SaaS nuevo a partir de la plataforma (ADR-005, P-01) & Pendiente \\
    \bottomrule
  \end{tabularx}
\end{table}

\todo{Asignar probabilidad e impacto (bajo, medio o alto) a cada riesgo según
el criterio del autor y, si se desea, una matriz de riesgos.}

Tres riesgos materializados merecen comentario por su enseñanza para la
planificación:

\begin{itemize}
  \item \textbf{B-04, exposición del entorno.} Los servicios locales de la
        base de datos publicaban sus puertos en todas las interfaces y el
        motor de contenedores insertaba reglas que se saltaban el cortafuegos,
        de modo que la base de datos de desarrollo, con credenciales conocidas
        y solo datos de prueba, fue accesible desde Internet durante
        aproximadamente una hora. Se corrigió con una regla persistente, se
        recrearon los volúmenes y se pasó a trabajar mediante un túnel SSH
        (ADR-010).
  \item \textbf{Fallos de control de acceso.} Las auditorías de la base de
        datos del CMS encontraron brechas reales, entre ellas una
        escalada de privilegios en el RBAC (B-05), comprobaciones que fallaban
        en abierto (B-06, B-23), una inyección SQL (B-24) y un escape de las
        denegaciones explícitas mediante delegación (B-47). El riesgo estaba
        previsto, y la mitigación ---auditoría obligatoria en todo cambio de
        base de datos--- es la que permitió detectarlos.
  \item \textbf{B-15, corte de conexión con un agente trabajando.} El trabajo
        quedó a medias y mezclado con una corrección de seguridad pendiente. La
        lección pasó al proceso: cada incremento se cierra con su propio
        \textit{commit}.
\end{itemize}

% ---------------------------------------------------------------------
\section{Gestión de la configuración}
\label{sec:plan-configuracion}

\begin{itemize}
  \item \textbf{Control de versiones.} Git, con un repositorio privado en
        GitHub (ADR-001). La
        rama principal es \texttt{main}; se trabaja en ramas por fase
        (\texttt{fase-1/base}, \texttt{fase-2/cms}, \texttt{fase-3/espanol}) que
        se fusionan en \texttt{main} al cerrar la fase, mediante
        \textit{pull request} cuando se requiere ejecutar la suite E2E en la
        CI. Los mensajes de \textit{commit} están en español y siguen
        \textit{Conventional Commits}~\cite{conventionalcommits} con el módulo
        como ámbito, por ejemplo \texttt{feat(cms): auditoría y búsqueda
        global (F2.6)}.
  \item \textbf{Base de datos como código.} El esquema se mantiene de forma
        declarativa en ficheros numerados, y las migraciones se generan a
        partir de ellos y se versionan en el repositorio. Nunca se modifica una
        migración ya aplicada: cada cambio es una migración nueva y no
        destructiva. Las pruebas pgTAP viven junto al esquema.
  \item \textbf{Configuración y secretos.} La configuración reutilizable está
        en ficheros \texttt{*.config.ts} validados con Zod y en variables de
        entorno públicas documentadas; los secretos solo se guardan en
        ficheros locales excluidos del repositorio o en las variables de la
        plataforma de despliegue.
  \item \textbf{Dependencias.} Un catálogo único de versiones para todo el
        monorepo, con el fichero de bloqueo versionado.
  \item \textbf{Entregas.} Cada fusión de una fase en \texttt{main} constituye
        una versión estable del producto (tabla~\ref{tab:plan-hitos}). El
        despliegue de demostración se construye desde \texttt{main}.
        \todo{Decidir si se etiquetan versiones (\texttt{git tag}) para la
        entrega final del TFG.}
\end{itemize}

% ---------------------------------------------------------------------
\section{Aseguramiento de calidad}
\label{sec:plan-calidad}

\subsection{Estándares y normas}

\begin{itemize}
  \item \textbf{Código}: TypeScript estricto, sin \texttt{any} sin
        justificar; \textit{lint} y formato automáticos; convenciones del
        repositorio para \textit{server functions}, formularios, clientes de
        base de datos y separación entre módulos de cliente y de servidor.
  \item \textbf{Comentarios}: guía propia de comentarios didácticos en
        español, con cabecera por fichero, documentación de las funciones
        exportadas y justificación obligatoria de toda decisión de seguridad.
  \item \textbf{Base de datos}: RLS obligatoria en toda tabla nueva,
        privilegios mínimos por verbo y por columna, y ninguna función con
        privilegios elevados sin controles explícitos.
  \item \textbf{Seguridad}: el estándar OWASP ASVS~\cite{owasp-asvs} sirve de
        referencia para los controles de autenticación, sesión y control de
        acceso.
  \item \textbf{Calidad del producto}: los requisitos no funcionales se
        alinean con las características del modelo ISO/IEC
        25010~\cite{iso25010}: seguridad (RNF-02, RNF-03), mantenibilidad y
        modularidad (RNF-04, RNF-05), portabilidad e instalabilidad (RNF-07),
        usabilidad (RNF-08) y eficiencia (RNF-09).
  \item \textbf{Trazabilidad}: una matriz relaciona cada requisito con el
        código que lo implementa, las pruebas que lo verifican y la sección de
        la memoria que lo describe. Los puntos clave del código llevan
        etiquetas \texttt{[TFG]} que un script extrae a esa matriz.
\end{itemize}

\subsection{Verificación y validación}

Todo cambio pasa por la misma cadena de comprobaciones, declarada obligatoria
en las directrices del repositorio:

\begin{enumerate}
  \item comprobación de tipos, \textit{lint} y formato;
  \item controles automáticos del repositorio, tanto sobre el contenido
        como sobre las rutas de los ficheros;
  \item pruebas afectadas: unitarias, de base de datos (pgTAP) y E2E si se
        toca un flujo de usuario;
  \item revisión adversarial del cambio en dos pistas: calidad y
        arquitectura, y búsqueda de errores y fallos de seguridad;
  \item si el cambio afecta a la base de datos, auditoría adversarial de las
        políticas RLS en cuatro etapas: revisión estática, pruebas de ataque
        ejecutables entre clientes, refutación independiente de los hallazgos
        y corrección con pruebas de regresión;
  \item revisión de los comentarios del código; y
  \item actualización del progreso, la trazabilidad y, si procede, las
        decisiones y la bitácora.
\end{enumerate}

Las comprobaciones automatizables se repiten en la integración continua, en un
\textit{workflow} con tres trabajos: (1)~controles propios del TFG
(controles automáticos del repositorio, bloqueantes; informe de comentarios en inglés, informativo hasta
cerrar la F4); (2)~tipos, \textit{lint}, formato y pruebas unitarias; y
(3)~base de datos local, pruebas pgTAP y E2E con Playwright contra la
\textit{build} de test, que se ejecuta en las \textit{pull requests} hacia
\texttt{main} y bajo demanda.

\subsection{Criterios de aceptación y acciones correctoras}

Un incremento se acepta cuando cumple su criterio de «hecho», la cadena
anterior está en verde y el autor lo ha revisado. Se rechaza ---y se corrige
antes de continuar--- si una auditoría encuentra una brecha confirmada, si
alguna prueba deja de pasar o si los controles automáticos del repositorio
emiten un aviso. Cada incidencia se registra en la bitácora con su lección, y
cuando la lección es generalizable se convierte en una acción preventiva: una
regla nueva en las directrices, una variante más en un control automático
(B-01) o una prueba de regresión. A la fecha de corte la bitácora
recoge 55 incidencias, de las que 29 son de seguridad; todas las de gravedad alta están corregidas.



  
    % DESARROLLO
    \part{Desarrollo}
    \null\vfill
    \noindent En esta parte se debe describir el desarrollo del proyecto siguiendo la metodología empleada. Sus capítulos no deben ser una descripción exhaustiva de todos los documentos, diagramas, código fuente y, en general, entregables generados, sino más bien una explicación resumida del desarrollo, estructurada según las etapas principales del proceso de ingeniería. Deben seleccionarse aquellos diagramas, fragmentos de código y secciones de los entregables que sean más significativos para dicha explicación. La totalidad de los entregables resultado del proyecto se ubicarán en los anexos y/o en el material digital que acompañe al proyecto.

\chapter{Requisitos del Sistema}
% =====================================================================
% Capítulo: Requisitos del Sistema
% Fuentes: propuesta del TFG (Objetivo, Descripción y Alcance),
% docs/tfg/REQUISITOS.md, DECISIONES.md (ADR-004, 005, 011, 017, 020, 021)
% y borradores/04-analisis.tex. Estado de las fuentes: F3c cerrada.
% =====================================================================

En este capítulo se presentan los objetivos de negocio y del sistema, los
actores que interactúan con PymeKit y el catálogo de requisitos funcionales,
no funcionales y de información. Al final se relaciona cada requisito con el
apartado de la propuesta del TFG del que procede y se resume la brecha entre
esos requisitos y las bases de código de partida.

Los identificadores (OBN, OBS, RF, RNF, RI) son estables: se citan en el
código mediante etiquetas \texttt{[TFG]}, en la matriz de trazabilidad del
repositorio y en esta memoria, y no se renumeran. Si un requisito cambia, se
modifica su redacción y se registra la decisión que lo justifica.

% ---------------------------------------------------------------------
\section{Objetivos de negocio}
\label{sec:req-objetivos-negocio}

Quien desarrolla un producto SaaS para pymes debe resolver, antes de escribir
su lógica de negocio, un conjunto de problemas que se repiten en casi todos
los productos de este tipo: identidad, aislamiento entre clientes, equipos y
permisos, cobro, administración interna y calidad transversal. Estos
componentes no aportan valor diferencial, pero consumen gran parte del
esfuerzo y concentran el riesgo de seguridad. De ahí se derivan los objetivos
de negocio de la tabla~\ref{tab:req-obn}.

\begin{table}[htbp]
  \centering
  \footnotesize
  \caption{Objetivos de negocio.}
  \label{tab:req-obn}
  \begin{tabularx}{\textwidth}{@{}lX@{}}
    \toprule
    \textbf{ID} & \textbf{Objetivo} \\
    \midrule
    OBN-01 & Reducir el tiempo de puesta en marcha de un nuevo producto SaaS para pymes, partiendo de una base que ya resuelve los componentes comunes. \\
    OBN-02 & Estandarizar buenas prácticas de seguridad, control de acceso y persistencia, de modo que cada producto nuevo las herede sin reimplementarlas. \\
    OBN-03 & Permitir que una pyme administre su web pública, su contenido y sus datos desde una consola, sin acceder directamente a la base de datos ni depender de un desarrollador para cada cambio. \\
    OBN-04 & Facilitar la monetización mediante suscripciones con una pasarela de pago estándar. \\
    OBN-05 & Servir como material didáctico y técnico documentado para quien reutilice la plataforma (código comentado en español, decisiones registradas y manuales). \\
    \bottomrule
  \end{tabularx}
\end{table}

Durante el desarrollo se precisó el enfoque del producto (ADR-020): PymeKit es
la base de la web de \emph{una} pyme, con una página pública, un blog
gestionado desde la consola y un inicio de sesión que puede corresponder a un
super-administrador, a personal interno o a un cliente. Los módulos SaaS
(equipos, invitaciones y suscripciones) se conservan como opcionales, porque
permiten tanto ese escenario como el de un SaaS multi-cliente.

% ---------------------------------------------------------------------
\section{Objetivos del sistema}
\label{sec:req-objetivos-sistema}

Los objetivos del sistema concretan los componentes que enumera el objetivo de
la propuesta (tabla~\ref{tab:req-obs}).

\begin{table}[htbp]
  \centering
  \footnotesize
  \caption{Objetivos del sistema y objetivos de negocio a los que contribuyen.}
  \label{tab:req-obs}
  \begin{tabularx}{\textwidth}{@{}lXl@{}}
    \toprule
    \textbf{ID} & \textbf{Objetivo del sistema} & \textbf{OBN} \\
    \midrule
    OBS-01 & Ofrecer un \textit{frontend} web con página pública, precios, blog y área privada. & 01, 03 \\
    OBS-02 & Ofrecer un \textit{backend} con API tipada y validada para la lógica de negocio y el acceso a datos. & 01 \\
    OBS-03 & Autenticar a los usuarios y autorizar cada acceso, con segundo factor y aislamiento entre cuentas aplicado en la base de datos. & 02 \\
    OBS-04 & Gestionar usuarios, cuentas de equipo, invitaciones, roles y permisos. & 01, 02 \\
    OBS-05 & Integrar una pasarela de pago para suscripciones. & 04 \\
    OBS-06 & Proporcionar una consola de administración con un CMS de datos y contenidos, con permisos y auditoría propios. & 03 \\
    OBS-07 & Ser reutilizable por configuración y estar documentado y verificado de forma automática. & 01, 05 \\
    \bottomrule
  \end{tabularx}
\end{table}

% ---------------------------------------------------------------------
\section{Actores}
\label{sec:req-actores}

\begin{description}[style=nextline]
  \item[Visitante.] Persona no identificada que consulta la web pública
        (inicio, precios y blog) y puede registrarse.
  \item[Usuario.] Persona registrada. Dispone de una cuenta personal y puede
        pertenecer a una o varias cuentas de equipo.
  \item[Propietario del equipo.] Usuario con el rol de mayor nivel en una
        cuenta de equipo; gestiona sus miembros y su suscripción.
  \item[Miembro.] Usuario de un equipo con un rol de menor nivel y, por
        tanto, con permisos limitados.
  \item[Super-administrador.] Personal del proveedor con control total de la
        plataforma y, por extensión, del CMS (ADR-014). Debe usar siempre
        segundo factor.
  \item[Personal del CMS.] Personal de la pyme o del proveedor (gestor de
        contenidos, soporte) que accede a la consola con los permisos
        limitados que le asigna el control de acceso propio del CMS
        (ADR-014, ADR-016).
  \item[Desarrollador.] Quien reutiliza PymeKit para construir un producto
        nuevo; interactúa con el repositorio, la configuración y los
        manuales, no con la interfaz.
  \item[Pasarela de pago.] Sistema externo (Stripe) que procesa los cobros y
        notifica los cambios de suscripción mediante \textit{webhooks}
        (peticiones HTTP que el sistema externo envía a la aplicación cuando
        ocurre un evento).
\end{description}

% ---------------------------------------------------------------------
\section{Requisitos funcionales}
\label{sec:req-rf}

La prioridad es \textbf{M} (imprescindible: forma parte del prototipo mínimo
que exige la propuesta) o \textbf{D} (deseable). La última columna indica el
objetivo del sistema que desarrolla cada requisito.

\begin{table}[htbp]
  \centering
  \footnotesize
  \caption{Requisitos funcionales.}
  \label{tab:req-rf}
  \begin{tabularx}{\textwidth}{@{}lXcl@{}}
    \toprule
    \textbf{ID} & \textbf{Requisito} & \textbf{Prior.} & \textbf{OBS} \\
    \midrule
    RF-01 & Interfaz web para el usuario final: página de inicio pública (\textit{landing}), precios, blog público y área privada con navegación. & M & 01 \\
    RF-02 & API de \textit{backend} para la lógica de negocio y el acceso a datos mediante funciones de servidor (\textit{server functions}) tipadas y validadas. & M & 02 \\
    RF-03 & Registro e inicio de sesión (contraseña y enlace mágico), recuperación de contraseña y verificación del correo electrónico. & M & 03 \\
    RF-04 & Segundo factor de autenticación (MFA) y proveedores OAuth configurables. & D & 03 \\
    RF-05 & Gestión de la cuenta personal: perfil, correo, contraseña y baja. & M & 04 \\
    RF-06 & Cuentas de equipo: crear equipos, invitar miembros, asignar roles y aplicar permisos por funcionalidad. & M & 04 \\
    RF-07 & Suscripciones con pasarela de pago: planes, pago (\textit{checkout}), portal de cliente y sincronización mediante \textit{webhooks}. & M & 05 \\
    RF-08 & Panel de super-administración: listar y gestionar usuarios y cuentas, bloquear, suplantar y restablecer. & M & 06 \\
    RF-09 & CMS de datos y contenidos integrado en la consola de administración: explorar y editar tablas, gestionar usuarios y almacenamiento, publicar el contenido del blog, con permisos propios. & M & 06 \\
    RF-10 & Registro de auditoría de las acciones realizadas desde el CMS. & D & 06 \\
    RF-11 & Paneles (\textit{dashboards}) configurables en el CMS. & D & 06 \\
    RF-12 & Notificaciones dentro de la aplicación y correos transaccionales. & D & 04 \\
    RF-13 & Interfaz en español, con la infraestructura de internacionalización preparada para añadir otros idiomas. & M & 07 \\
    \bottomrule
  \end{tabularx}
\end{table}

Dos requisitos han cambiado de redacción durante el proyecto, y ambos cambios
están registrados como decisiones:

\begin{itemize}
  \item \textbf{RF-09} se planteó al principio como un CMS en una aplicación
        separada que compartía la base de datos. Tras evaluar alternativas se
        decidió integrarlo en la consola de administración (ADR-011) y,
        después, usarlo también para gestionar el contenido del blog que
        publica la web (ADR-017), que es el caso de uso de «administración de
        contenidos» que pide la propuesta.
  \item \textbf{RF-13} pedía en origen una interfaz en español por defecto y
        en inglés. A petición del autor, la interfaz pasó a ofrecer
        \emph{solo} español, conservando claves, mensajes y configuración de
        idiomas como punto de extensión documentado (ADR-021).
\end{itemize}

RF-11 se identificó desde el principio como el primer candidato a recortar si
el plazo lo exigía; finalmente se implementó (ADR-019).

% ---------------------------------------------------------------------
\section{Requisitos no funcionales}
\label{sec:req-rnf}

Los requisitos no funcionales se clasifican según las características de
calidad de la norma ISO/IEC 25010~\cite{iso25010}. Cada uno se acompaña del
modo en que se verifica, para que sea comprobable y no una mera declaración
de intenciones.

\begin{table}[htbp]
  \centering
  \footnotesize
  \caption{Requisitos no funcionales, característica ISO/IEC 25010 y forma de verificarlos.}
  \label{tab:req-rnf}
  \begin{tabularx}{\textwidth}{@{}l>{\raggedright\arraybackslash}p{5.4cm}>{\raggedright\arraybackslash}p{2.2cm}X@{}}
    \toprule
    \textbf{ID} & \textbf{Requisito} & \textbf{ISO 25010} & \textbf{Verificación} \\
    \midrule
    RNF-01 & \textbf{Reutilización y configurabilidad}: la plataforma se adapta a un producto nuevo cambiando la configuración (ficheros \texttt{*.config.ts} y variables de entorno), sin tocar el núcleo. & Mantenibilidad (reusabilidad), flexibilidad & Evaluación de la reutilización: pasos y tiempo para arrancar un proyecto nuevo \\
    RNF-02 & \textbf{Seguridad}: la autorización se aplica en la base de datos (RLS) además de en la aplicación; secretos fuera del repositorio; validación de todas las entradas. & Seguridad & Pruebas pgTAP y revisiones adversariales de RLS y del código \\
    RNF-03 & \textbf{Aislamiento multi-tenant}: los datos de una cuenta no son accesibles desde otra. & Seguridad (confidencialidad) & Pruebas pgTAP de aislamiento \\
    RNF-04 & \textbf{Modularidad}: monorepo de paquetes desacoplados con dependencias explícitas y exportaciones separadas para cliente y servidor. & Mantenibilidad (modularidad) & Estructura del monorepo, comprobación de dependencias y de tipos \\
    RNF-05 & \textbf{Mantenibilidad}: TypeScript estricto, \textit{lint} y formato automáticos y código comentado en español didáctico. & Mantenibilidad (analizabilidad) & Integración continua y revisión de comentarios \\
    RNF-06 & \textbf{Calidad verificable}: pruebas unitarias, de base de datos y E2E de los casos de uso principales. & Mantenibilidad (capacidad de prueba) & Integración continua en verde \\
    RNF-07 & \textbf{Desplegabilidad}: despliegue reproducible con contenedores y en la nube, con manual de instalación. & Flexibilidad (instalabilidad) & Despliegue de prueba \\
    RNF-08 & \textbf{Usabilidad y accesibilidad}: interfaz adaptable, componentes accesibles, tema claro/oscuro y ayuda contextual en la consola. & Capacidad de interacción & Revisión manual y E2E \\
    RNF-09 & \textbf{Rendimiento razonable}: carga de datos en el servidor, caché de consultas y paginación en el CMS. & Eficiencia de desempeño & Revisión manual y pruebas puntuales \\
    \bottomrule
  \end{tabularx}
\end{table}
% VERIFICAR: nombres de las características según la edición de la norma que
% se cite (2011 o 2023; en 2023 «usabilidad» pasa a «capacidad de
% interacción» y «portabilidad» a «flexibilidad»).

% ---------------------------------------------------------------------
\section{Requisitos de información}
\label{sec:req-ri}

La tabla~\ref{tab:req-ri} recoge la información que el sistema debe almacenar.
Su modelo conceptual se desarrolla en el capítulo de Análisis
(sección~\ref{sec:ana-modelo-datos}).

\begin{table}[htbp]
  \centering
  \footnotesize
  \caption{Requisitos de información.}
  \label{tab:req-ri}
  \begin{tabularx}{\textwidth}{@{}l>{\raggedright\arraybackslash}p{3.1cm}Xl@{}}
    \toprule
    \textbf{ID} & \textbf{Información} & \textbf{Datos principales} & \textbf{RF} \\
    \midrule
    RI-01 & Usuarios & Identidad, correo, factores MFA, estado de bloqueo (gestionados por el servicio de autenticación) & 03--05 \\
    RI-02 & Cuentas & Cuenta personal o de equipo, nombre, identificador legible (\textit{slug}), propietario principal & 05, 06 \\
    RI-03 & Membresías y roles & Pertenencia de un usuario a una cuenta, rol con nivel jerárquico y permisos por funcionalidad & 06 \\
    RI-04 & Invitaciones & Correo invitado, rol ofrecido, testigo (\textit{token}), caducidad, quién invita & 06 \\
    RI-05 & Facturación & Cliente de la pasarela, suscripciones y sus líneas, pedidos de pago único & 07 \\
    RI-06 & Notificaciones & Destinatario, tipo, texto, enlace, leída o no & 12 \\
    RI-07 & Contenido del blog & Entradas (título, \textit{slug}, cuerpo en Markdown, estado, fecha de publicación, autor), categorías y etiquetas & 01, 09 \\
    RI-08 & Personal y permisos del CMS & Cuentas del CMS, roles con rango, grupos de permisos, permisos de datos, de sistema y de almacenamiento, con vigencia & 09 \\
    RI-09 & Metadatos del CMS & Por tabla: nombre visible, descripción, área de navegación, formato de visualización, configuración de columnas y diseño de la ficha; vistas guardadas & 09 \\
    RI-10 & Auditoría & Operación, tabla, registro, datos antes y después, gravedad, instantánea del autor & 10 \\
    RI-11 & Paneles & Panel, propietario, compartición por rol, \textit{widgets} y su posición en la cuadrícula & 11 \\
    \bottomrule
  \end{tabularx}
\end{table}

% ---------------------------------------------------------------------
\section{Trazabilidad con la propuesta}
\label{sec:req-trazabilidad}

La tabla~\ref{tab:req-traza-propuesta} relaciona los apartados de la
propuesta del TFG con los requisitos que los cubren y el capítulo de la
memoria donde se desarrollan. La relación inversa, de cada requisito con el
código y las pruebas que lo verifican, se presenta en el capítulo de Pruebas
(sección~\ref{sec:pru-resultados}).

\begin{table}[htbp]
  \centering
  \footnotesize
  \caption{Trazabilidad entre la propuesta del TFG, los requisitos y la memoria.}
  \label{tab:req-traza-propuesta}
  \begin{tabularx}{\textwidth}{@{}>{\raggedright\arraybackslash}X>{\raggedright\arraybackslash}p{3.6cm}>{\raggedright\arraybackslash}p{3.2cm}@{}}
    \toprule
    \textbf{Apartado de la propuesta} & \textbf{Requisitos} & \textbf{Capítulo} \\
    \midrule
    Objetivo: \textit{frontend}, \textit{backend}, autenticación, usuarios y roles, pagos y CMS & RF-01--RF-09 & Requisitos, Diseño \\
    Descripción: estructura desacoplada y configurable & RNF-01, RNF-04 & Diseño, Construcción \\
    Descripción: prototipo con gestión de usuarios, autenticación, persistencia, CMS y flujo de pago & RF-03, RF-05, RF-06, RF-07, RF-09 & Construcción, Pruebas \\
    Alcance 1: análisis y requisitos & Todo el catálogo & Requisitos, Análisis \\
    Alcance 2: arquitectura, autenticación, roles, pagos y CMS & RF-02--RF-04, RF-06--RF-09, RNF-02 & Diseño \\
    Alcance 3: modelo de datos y mecanismos de acceso & RI-01--RI-11, RNF-02, RNF-03 & Análisis, Diseño \\
    Alcance 4: implementación del prototipo y configuración para su reutilización & RF-01--RF-13, RNF-01, RNF-05 & Construcción \\
    Alcance 5: verificación de casos de uso, integración y evaluación de la reutilización & RNF-06, RNF-01 & Pruebas \\
    Alcance 6: memoria, documentación técnica y manuales & RNF-05, RNF-07 & Despliegue, anexos \\
    \bottomrule
  \end{tabularx}
\end{table}

% ---------------------------------------------------------------------
\section{Análisis GAP}
\label{sec:req-gap}

% PENDIENTE (P-06): redacción neutra acordada con el tutor; no se nombran las
% bases de código de partida (ADR-007).
Como se justifica en el capítulo de Planificación (ADR-001), PymeKit no parte
de cero: se apoya en dos bases de código existentes, utilizadas con licencia.
La primera es una base SaaS sobre la misma pila tecnológica, que ya resolvía
la autenticación, el modelo de cuentas multi-tenant, los equipos y la
facturación. La segunda es un CMS de administración de datos para el mismo
tipo de base de datos, distribuido como aplicación de página única con su
propia API. La tabla~\ref{tab:req-gap} resume qué cubría cada una frente a
los requisitos y qué ha tenido que aportar el proyecto.

\begin{table}[htbp]
  \centering
  \footnotesize
  \caption{Brecha entre los requisitos y las bases de partida («sí»: cubierto; «parcial»: cubierto en parte; «no»: no cubierto).}
  \label{tab:req-gap}
  \begin{tabularx}{\textwidth}{@{}lccX@{}}
    \toprule
    \textbf{Req.} & \textbf{Base SaaS} & \textbf{Base CMS} & \textbf{Aportación del proyecto} \\
    \midrule
    RF-01 & parcial & no & Blog en la base de datos con lectura pública por RLS; marca e identidad propias \\
    RF-02--RF-05 & sí & no & Comprobación de sesión viva en las rutas protegidas; textos en español \\
    RF-06, RF-07 & sí & no & Integración y pruebas; configuración como módulos opcionales \\
    RF-08 & sí & no & Consola compartida con el CMS y MFA obligatorio en ella \\
    RF-09 & no & parcial & Integración total en la web, reescritura de las pantallas, consola por áreas, nombres en español, endurecimiento de seguridad \\
    RF-10, RF-11 & no & parcial & Auditoría no falsificable, autor conservado, paneles sin dependencias nuevas \\
    RF-13 & parcial & parcial & Un único sistema de traducción y toda la interfaz en español \\
    RNF-02, RNF-03 & parcial & parcial & Revisiones adversariales y corrección de brechas heredadas \\
    RNF-05 & no & no & Comentarios didácticos en español y controles automáticos del TFG \\
    \bottomrule
  \end{tabularx}
\end{table}

La brecha más importante no está en la funcionalidad, sino en la
\emph{integración} y en la \emph{seguridad}: las dos bases tenían modelos de
autorización, sistemas de formularios y de traducción y enrutadores distintos,
y el CMS de partida contenía fallos de control de acceso que se describen en
los capítulos de Diseño y Pruebas. La distinción detallada entre lo
reutilizado, lo adaptado y lo propio se presenta en el capítulo de
Construcción (sección~\ref{sec:imp-reutilizacion}).


\newpage{\ }
\cleardoublepage

\chapter{Análisis del Sistema}
% =====================================================================
% Capítulo: Análisis del Sistema
% Fuentes: borradores/04-analisis.tex, docs/tfg/REQUISITOS.md,
% apps/web/supabase/schemas/*.sql, apps/e2e/tests, DECISIONES.md.
% Estado de las fuentes: F3c cerrada.
% =====================================================================

Este capítulo analiza el comportamiento esperado del sistema a partir de los
requisitos del capítulo anterior. Se presentan el modelo de casos de uso, con
la especificación de los principales, el modelo conceptual de datos y las
reglas de negocio que el sistema debe respetar.

% ---------------------------------------------------------------------
\section{Modelo de casos de uso}
\label{sec:ana-casos-uso}

Se han seleccionado ocho casos de uso principales que recorren todos los
módulos de la plataforma y que sirven de base para su validación
(tabla~\ref{tab:ana-cu} y figura~\ref{fig:ana-casos-uso}). Cada uno tiene una
o varias pruebas E2E que lo ejercitan (capítulo de Pruebas).

\begin{table}[htbp]
  \centering
  \footnotesize
  \caption{Casos de uso principales.}
  \label{tab:ana-cu}
  \begin{tabularx}{\textwidth}{@{}l>{\raggedright\arraybackslash}p{2.6cm}X>{\raggedright\arraybackslash}p{2.4cm}@{}}
    \toprule
    \textbf{CU} & \textbf{Actor principal} & \textbf{Descripción} & \textbf{Requisitos} \\
    \midrule
    CU-01 & Visitante & Se registra, verifica el correo y accede al área privada & RF-01, RF-03 \\
    CU-02 & Usuario & Crea un equipo e invita a un empleado con el rol \textit{member} & RF-06 \\
    CU-03 & Propietario del equipo & Contrata un plan y el estado de la suscripción se refleja en la aplicación & RF-07 \\
    CU-04 & Miembro sin permiso & Intenta gestionar la facturación y se le deniega en la interfaz, en el servidor y en la base de datos & RF-06, RNF-02, RNF-03 \\
    CU-05 & Super-administrador & Busca un usuario, lo bloquea y lo suplanta para darle soporte & RF-08 \\
    CU-06 & Personal del CMS & Edita registros de una tabla según sus permisos; la acción queda auditada & RF-09, RF-10 \\
    CU-07 & Desarrollador & Arranca un producto nuevo a partir de PymeKit cambiando solo la configuración & RNF-01, RNF-07 \\
    CU-08 & Personal del CMS & Publica una entrada del blog desde la consola, que el visitante lee en la web pública & RF-01, RF-09 \\
    \bottomrule
  \end{tabularx}
\end{table}

\begin{figure}[htbp]
  \centering
  \resizebox{\textwidth}{!}{%
  \begin{tikzpicture}[
      font=\footnotesize,
      uc/.style={ellipse, draw, minimum width=4.6cm, minimum height=0.85cm,
                 align=center, inner sep=1pt},
      lbl/.style={align=center, font=\footnotesize\bfseries},
      link/.style={-},
      actor/.pic={
        \draw (0,0.42) circle (0.13);
        \draw (0,0.29) -- (0,-0.05);
        \draw (-0.22,0.17) -- (0.22,0.17);
        \draw (0,-0.05) -- (-0.17,-0.38);
        \draw (0,-0.05) -- (0.17,-0.38);
      }
    ]
    % Casos de uso
    \node[uc] (cu1) at (0, 4.2)  {CU-01 Registrarse y\\acceder al área privada};
    \node[uc] (cu2) at (0, 3.0)  {CU-02 Crear equipo\\e invitar a un miembro};
    \node[uc] (cu3) at (0, 1.8)  {CU-03 Contratar\\un plan};
    \node[uc] (cu4) at (0, 0.6)  {CU-04 Gestionar la\\facturación (denegado)};
    \node[uc] (cu5) at (0, -0.6) {CU-05 Bloquear y\\suplantar a un usuario};
    \node[uc] (cu6) at (0, -1.8) {CU-06 Editar registros\\en el CMS (auditado)};
    \node[uc] (cu8) at (0, -3.0) {CU-08 Publicar una\\entrada del blog};
    \node[uc] (cu7) at (0, -4.2) {CU-07 Arrancar un\\producto nuevo};

    \begin{scope}[on background layer]
      \node[draw, rounded corners, fit=(cu1)(cu7), inner sep=12pt,
            label={[font=\small\bfseries]above:PymeKit}] (sys) {};
    \end{scope}

    % Actores de la izquierda (clientes)
    \pic at (-6.2, 4.0) {actor};  \node[lbl] (vis) at (-6.2, 3.3)  {Visitante};
    \pic at (-6.2, 2.0) {actor};  \node[lbl] (usu) at (-6.2, 1.3)  {Usuario};
    \pic at (-6.2, 0.0) {actor};  \node[lbl] (pro) at (-6.2, -0.75) {Propietario\\del equipo};
    \pic at (-6.2, -2.2) {actor}; \node[lbl] (mie) at (-6.2, -2.95) {Miembro\\sin permiso};

    % Actores de la derecha (personal y sistemas externos)
    \pic at (6.2, 2.4) {actor};   \node[lbl] (pas) at (6.2, 1.7)  {Pasarela\\de pago};
    \pic at (6.2, 0.2) {actor};   \node[lbl] (sa)  at (6.2, -0.5) {Super-\\administrador};
    \pic at (6.2, -2.0) {actor};  \node[lbl] (ges) at (6.2, -2.7) {Personal\\del CMS};
    \pic at (6.2, -4.2) {actor};  \node[lbl] (dev) at (6.2, -4.9) {Desarrollador};

    % Puntos de anclaje junto a cada actor
    \coordinate (avis) at (-5.8, 4.1);
    \coordinate (ausu) at (-5.8, 2.1);
    \coordinate (apro) at (-5.8, 0.1);
    \coordinate (amie) at (-5.8, -2.1);
    \coordinate (apas) at (5.8, 2.5);
    \coordinate (asa)  at (5.8, 0.3);
    \coordinate (ages) at (5.8, -1.9);
    \coordinate (adev) at (5.8, -4.1);

    \draw[link] (avis) -- (cu1.west);
    \draw[link] (avis) -- (cu8.west);
    \draw[link] (ausu) -- (cu2.west);
    \draw[link] (apro) -- (cu3.west);
    \draw[link] (amie) -- (cu4.west);
    \draw[link] (apas) -- (cu3.east);
    \draw[link] (asa)  -- (cu5.east);
    \draw[link] (ages) -- (cu6.east);
    \draw[link] (ages) -- (cu8.east);
    \draw[link] (adev) -- (cu7.east);
  \end{tikzpicture}}
  \caption{Diagrama de casos de uso principales de PymeKit.}
  \label{fig:ana-casos-uso}
  % Texto alternativo: el sistema PymeKit contiene ocho casos de uso. A la
  % izquierda, los actores cliente: el visitante (CU-01 y CU-08, como lector
  % del blog), el usuario (CU-02), el propietario del equipo (CU-03) y el
  % miembro sin permiso (CU-04). A la derecha, la pasarela de pago como actor
  % secundario de CU-03, el super-administrador (CU-05), el personal del CMS
  % (CU-06 y CU-08) y el desarrollador (CU-07).
\end{figure}

El propietario y el miembro son especializaciones del actor \emph{usuario}, y
el super-administrador lo es del \emph{personal del CMS}, porque es la raíz
del control de acceso de la consola (ADR-014). Para no recargar la figura,
esas generalizaciones no se dibujan.

\subsection{Especificación de los casos de uso principales}
\label{sec:ana-fichas}

Las fichas siguientes describen el flujo normal y los flujos alternativos más
relevantes de cada caso de uso. Los alternativos se numeran con el paso del
flujo normal en el que se producen.

% --- CU-01 -----------------------------------------------------------
\begin{table}[H]
  \centering
  \footnotesize
  \caption{CU-01 · Registrarse y acceder al área privada.}
  \label{tab:ana-cu01}
  \begin{tabularx}{\textwidth}{@{}>{\bfseries}p{2.6cm}X@{}}
    \toprule
    Actor principal & Visitante \\
    Precondiciones & El registro con contraseña está activado en la configuración de autenticación. \\
    \midrule
    Flujo normal &
      \begin{enumerate}[nosep,leftmargin=*]
        \item El visitante abre la página de registro e introduce su correo y una contraseña.
        \item El sistema valida los datos, crea el usuario y envía un correo de verificación.
        \item El visitante abre el enlace del correo.
        \item El sistema confirma el correo, crea su cuenta personal (con el mismo identificador que el usuario) e inicia la sesión.
        \item El sistema muestra el área privada en el contexto de la cuenta personal.
      \end{enumerate} \\
    \midrule
    Flujos alternativos &
      \begin{itemize}[nosep,leftmargin=*]
        \item[1a.] El visitante elige el enlace mágico: recibe un correo con un enlace de un solo uso y continúa en el paso 3.
        \item[2a.] El correo ya está registrado o la contraseña no cumple la política: se muestra el error y no se crea nada.
        \item[3a.] El enlace ha caducado: se ofrece reenviar el correo.
        \item[4a.] El usuario tiene un segundo factor configurado: se le pide el código y solo tras verificarlo (sesión \texttt{aal2}) accede.
        \item[5a.] La sesión deja de existir (por ejemplo, se revoca): al navegar se limpian las \textit{cookies} y se redirige al inicio de sesión (B-53).
      \end{itemize} \\
    \midrule
    Postcondición & Existe un usuario verificado con su cuenta personal y una sesión activa. \\
    Requisitos & RF-01, RF-03, RF-04 \\
    \bottomrule
  \end{tabularx}
\end{table}

% --- CU-02 -----------------------------------------------------------
\begin{table}[H]
  \centering
  \footnotesize
  \caption{CU-02 · Crear un equipo e invitar a un miembro.}
  \label{tab:ana-cu02}
  \begin{tabularx}{\textwidth}{@{}>{\bfseries}p{2.6cm}X@{}}
    \toprule
    Actor principal & Usuario (pasa a ser propietario del equipo) \\
    Precondiciones & El usuario ha iniciado sesión. Los equipos están activados (\texttt{VITE\_ACCOUNT\_MODE} distinto de \texttt{personal-only}). \\
    \midrule
    Flujo normal &
      \begin{enumerate}[nosep,leftmargin=*]
        \item El usuario elige «Crear equipo» e introduce el nombre.
        \item El sistema crea la cuenta de equipo con un \textit{slug} único, lo registra como propietario principal y le asigna el rol de mayor nivel.
        \item El propietario abre «Miembros», introduce el correo de un empleado y elige el rol \textit{member}.
        \item El sistema comprueba el permiso \texttt{invites.manage} y que el rol ofrecido no supera el suyo, crea la invitación con un testigo y una caducidad de siete días y envía el correo.
        \item El invitado abre el enlace, inicia sesión o se registra y acepta la invitación.
        \item El sistema crea la membresía con el rol de la invitación y elimina la invitación.
      \end{enumerate} \\
    \midrule
    Flujos alternativos &
      \begin{itemize}[nosep,leftmargin=*]
        \item[1a.] La creación de equipos está desactivada por configuración: la opción no se ofrece.
        \item[4a.] El correo ya pertenece a un miembro, o se supera el límite de invitaciones (20 por hora y 50 pendientes por cuenta): se rechaza.
        \item[4b.] El rol ofrecido es superior al del que invita: se rechaza en el servidor y en la base de datos.
        \item[5a.] La invitación ha caducado o ya se usó: se informa y no se crea la membresía.
      \end{itemize} \\
    \midrule
    Postcondición & El equipo tiene dos miembros: el propietario y el empleado con el rol \textit{member}. \\
    Requisitos & RF-06, RF-12 \\
    \bottomrule
  \end{tabularx}
\end{table}

% --- CU-03 -----------------------------------------------------------
\begin{table}[H]
  \centering
  \footnotesize
  \caption{CU-03 · Contratar un plan.}
  \label{tab:ana-cu03}
  \begin{tabularx}{\textwidth}{@{}>{\bfseries}p{2.6cm}X@{}}
    \toprule
    Actor principal & Propietario del equipo \\
    Actor secundario & Pasarela de pago \\
    Precondiciones & La facturación de equipos está activada y la pasarela está configurada (en modo de pruebas durante el desarrollo). El usuario tiene el permiso \texttt{billing.manage}. \\
    \midrule
    Flujo normal &
      \begin{enumerate}[nosep,leftmargin=*]
        \item El propietario abre «Facturación» y elige un plan de los definidos en la configuración.
        \item El sistema crea una sesión de pago en la pasarela y muestra el formulario de pago embebido.
        \item El propietario completa el pago.
        \item La pasarela notifica el resultado mediante un \textit{webhook} a la ruta de facturación de la web.
        \item El sistema verifica la firma del \textit{webhook} y crea o actualiza el cliente, la suscripción y sus líneas.
        \item La página de facturación muestra el plan y el estado de la suscripción.
      \end{enumerate} \\
    \midrule
    Flujos alternativos &
      \begin{itemize}[nosep,leftmargin=*]
        \item[1a.] La facturación está desactivada para ese tipo de cuenta: se redirige a los ajustes.
        \item[3a.] El pago se rechaza: la pasarela lo indica y no se crea la suscripción.
        \item[5a.] La firma del \textit{webhook} no es válida: se rechaza la petición sin modificar datos.
        \item[6a.] El propietario abre el portal de cliente para cambiar o cancelar el plan; los cambios llegan de nuevo por \textit{webhook}.
      \end{itemize} \\
    \midrule
    Postcondición & La cuenta tiene una suscripción activa sincronizada con la pasarela. \\
    Requisitos & RF-07 \\
    \bottomrule
  \end{tabularx}
\end{table}

% --- CU-04 -----------------------------------------------------------
\begin{table}[H]
  \centering
  \footnotesize
  \caption{CU-04 · Gestionar la facturación sin permiso (acceso denegado).}
  \label{tab:ana-cu04}
  \begin{tabularx}{\textwidth}{@{}>{\bfseries}p{2.6cm}X@{}}
    \toprule
    Actor principal & Miembro sin el permiso \texttt{billing.manage} \\
    Precondiciones & El miembro pertenece a un equipo con el rol \textit{member}. \\
    \midrule
    Flujo normal &
      \begin{enumerate}[nosep,leftmargin=*]
        \item El miembro abre la página de facturación del equipo.
        \item El sistema muestra el estado de la suscripción, pero no ofrece contratar, cambiar de plan ni abrir el portal.
      \end{enumerate} \\
    \midrule
    Flujos alternativos &
      \begin{itemize}[nosep,leftmargin=*]
        \item[2a.] El miembro invoca directamente la función de servidor de pago: el \textit{middleware} de permisos la rechaza antes de llegar a la pasarela.
        \item[2b.] El miembro intenta escribir en las tablas de facturación con su sesión a través de la API de datos: la base de datos lo impide, porque esas tablas solo las escribe la sincronización con la pasarela.
        \item[2c.] El miembro intenta leer la facturación de un equipo al que no pertenece: las políticas RLS no devuelven ninguna fila.
      \end{itemize} \\
    \midrule
    Postcondición & No se modifica ningún dato; el intento no tiene efecto en ninguna capa. \\
    Requisitos & RF-06, RNF-02, RNF-03 \\
    \bottomrule
  \end{tabularx}
\end{table}

% --- CU-05 -----------------------------------------------------------
\begin{table}[H]
  \centering
  \footnotesize
  \caption{CU-05 · Bloquear y suplantar a un usuario.}
  \label{tab:ana-cu05}
  \begin{tabularx}{\textwidth}{@{}>{\bfseries}p{2.6cm}X@{}}
    \toprule
    Actor principal & Super-administrador \\
    Precondiciones & El actor es super-administrador y tiene un segundo factor configurado. \\
    \midrule
    Flujo normal &
      \begin{enumerate}[nosep,leftmargin=*]
        \item El super-administrador inicia sesión y completa el segundo factor.
        \item Abre la consola y, en el área Cuentas, «Gestión de cuentas».
        \item Busca al usuario por nombre o correo y abre su ficha.
        \item Elige «Bloquear» y confirma la acción.
        \item El sistema bloquea al usuario, que ya no puede iniciar sesión.
        \item Elige «Suplantar» sobre otro usuario y confirma.
        \item El sistema abre una sesión como ese usuario para reproducir su problema.
      \end{enumerate} \\
    \midrule
    Flujos alternativos &
      \begin{itemize}[nosep,leftmargin=*]
        \item[2a.] La sesión no es \texttt{aal2}: si tiene factor, se redirige a la verificación y se vuelve a la página pedida; si no, la consola responde 404 (B-14).
        \item[2b.] El actor es personal del CMS, no super-administrador: se le redirige al CMS; las páginas de la plataforma no le están disponibles.
        \item[4a.] El super-administrador desbloquea más tarde al usuario o restablece su contraseña desde la misma ficha.
      \end{itemize} \\
    \midrule
    Postcondición & El usuario queda bloqueado; la suplantación se usa solo durante la sesión de soporte. \\
    Requisitos & RF-08, RNF-02 \\
    \bottomrule
  \end{tabularx}
\end{table}

% --- CU-06 -----------------------------------------------------------
\begin{table}[H]
  \centering
  \footnotesize
  \caption{CU-06 · Editar registros en el CMS con auditoría.}
  \label{tab:ana-cu06}
  \begin{tabularx}{\textwidth}{@{}>{\bfseries}p{2.6cm}X@{}}
    \toprule
    Actor principal & Personal del CMS (por ejemplo, con el rol Soporte) \\
    Precondiciones & El actor tiene una cuenta activa del CMS, un rol con permiso de lectura y actualización sobre la tabla y sesión \texttt{aal2}. \\
    \midrule
    Flujo normal &
      \begin{enumerate}[nosep,leftmargin=*]
        \item El actor abre la consola; la barra lateral solo muestra las áreas y tablas para las que tiene permiso.
        \item Abre una tabla del área correspondiente; el sistema muestra el listado paginado con nombres en español, filtros y búsqueda.
        \item Abre un registro y elige «Editar».
        \item El sistema genera el formulario a partir de los metadatos de la tabla, con su validación y una ayuda «?» por campo.
        \item El actor modifica un valor y guarda.
        \item El sistema comprueba el permiso, ejecuta la escritura con la identidad del actor y registra en la misma transacción una entrada de auditoría con los datos anteriores y nuevos.
        \item La ficha muestra el valor actualizado.
      \end{enumerate} \\
    \midrule
    Flujos alternativos &
      \begin{itemize}[nosep,leftmargin=*]
        \item[1a.] El actor escribe en la barra de direcciones la URL de una tabla sin permiso: la API responde 403 y la interfaz lo indica.
        \item[5a.] Un valor no es válido para su tipo o falta un campo obligatorio: se marca el campo y no se envía; si la base de datos lo rechaza, se muestra un mensaje con un código estable, nunca el texto interno de PostgreSQL (B-20).
        \item[6a.] La condición enviada no identifica exactamente un registro: se rechaza con 400 (B-19).
        \item[6b.] Falla la escritura de la auditoría: se deshace también la modificación (B-41).
      \end{itemize} \\
    \midrule
    Postcondición & El registro está modificado y la auditoría conserva quién, cuándo y qué cambió. \\
    Requisitos & RF-09, RF-10, RNF-02 \\
    \bottomrule
  \end{tabularx}
\end{table}

% --- CU-08 -----------------------------------------------------------
\begin{table}[H]
  \centering
  \footnotesize
  \caption{CU-08 · Publicar una entrada del blog.}
  \label{tab:ana-cu08}
  \begin{tabularx}{\textwidth}{@{}>{\bfseries}p{2.6cm}X@{}}
    \toprule
    Actor principal & Personal del CMS con permiso sobre las tablas del blog \\
    Actor secundario & Visitante (lector) \\
    Precondiciones & Igual que en CU-06, para las tablas del área Blog. \\
    \midrule
    Flujo normal &
      \begin{enumerate}[nosep,leftmargin=*]
        \item El actor abre el área Blog, tabla «Entradas», y elige «Crear».
        \item Introduce título, \textit{slug}, resumen, cuerpo en Markdown y categoría, y deja el estado en borrador.
        \item El sistema guarda la entrada (no visible en la web) y la firma con la cuenta del actor si no se indica autor.
        \item El actor cambia el estado a «publicada».
        \item El sistema fija la fecha de publicación si no la tiene.
        \item El visitante abre \texttt{/blog}, ve la entrada en el listado y la abre; el Markdown se presenta saneado.
      \end{enumerate} \\
    \midrule
    Flujos alternativos &
      \begin{itemize}[nosep,leftmargin=*]
        \item[2a.] El \textit{slug} ya existe: se rechaza con un error de conflicto.
        \item[5a.] La fecha de publicación es futura: la entrada queda programada y no aparece en la web hasta esa fecha.
        \item[6a.] El visitante pide el \textit{slug} de un borrador o de una entrada archivada: la web responde 404, porque la política de lectura pública no la devuelve.
      \end{itemize} \\
    \midrule
    Postcondición & La entrada es visible en la web pública y en el mapa del sitio. \\
    Requisitos & RF-01, RF-09, RNF-02 \\
    \bottomrule
  \end{tabularx}
\end{table}

El caso CU-07 (\emph{arrancar un producto nuevo}) no tiene interacción con la
interfaz: el desarrollador clona el repositorio, ajusta los ficheros de
configuración y las variables de entorno descritos en el capítulo de Diseño
(sección~\ref{sec:dis-configuracion}), crea la base de datos con las
migraciones y despliega siguiendo el manual. Sirve para evaluar RNF-01 y
RNF-07.
% PENDIENTE (P-01): escenario y métricas concretas de la evaluación de la
% reutilización.

% ---------------------------------------------------------------------
\section{Modelo conceptual de datos}
\label{sec:ana-modelo-datos}

Toda la información reside en una única base de datos PostgreSQL. Las
entidades de la plataforma están en el esquema \texttt{public}; las del CMS, en
el esquema \texttt{cms}; y los usuarios, en el esquema \texttt{auth}, que
gestiona el servicio de autenticación. Las figuras~\ref{fig:ana-datos-public}
y~\ref{fig:ana-datos-cms} muestran las entidades principales y sus
relaciones, con los atributos más significativos. Se omiten las columnas de
trazabilidad (\texttt{created\_at}, \texttt{updated\_by}\ldots) presentes en
casi todas las tablas.

\begin{figure}[htbp]
  \centering
  \resizebox{\textwidth}{!}{%
  \begin{tikzpicture}[
      font=\scriptsize,
      ent/.style={draw, rectangle, align=left, inner sep=3pt,
                  minimum width=2.9cm, fill=white},
      ext/.style={ent, dashed},
      rel/.style={-},
      card/.style={font=\tiny, inner sep=1pt, fill=white}
    ]
    % Identidad y cuentas
    \node[ext] (users) at (-6, 0) {\textbf{auth.users}\\id\\email\\factores MFA};
    \node[ent] (acc) at (0, 0) {\textbf{accounts}\\id\\primary\_owner\_user\_id\\name, slug\\is\_personal\_account};
    \node[ent] (mem) at (-3, -3) {\textbf{accounts\_memberships}\\user\_id, account\_id\\account\_role};
    \node[ent] (roles) at (-6, -5.6) {\textbf{roles}\\name\\hierarchy\_level};
    \node[ent] (rp) at (-6, -8) {\textbf{role\_permissions}\\role\\permission};
    \node[ent] (inv) at (0, -5.6) {\textbf{invitations}\\email, role\\invite\_token\\expires\_at};

    % Facturación
    \node[ent] (bc) at (4.6, -3) {\textbf{billing\_customers}\\account\_id\\provider, customer\_id};
    \node[ent] (sub) at (4.6, -5.6) {\textbf{subscriptions}\\status, active\\period\_ends\_at};
    \node[ent] (si) at (4.6, -8) {\textbf{subscription\_items}\\product\_id, variant\_id\\price\_amount};
    \node[ent] (ord) at (8.6, -5.6) {\textbf{orders}\\status, total\_amount};
    \node[ent] (oi) at (8.6, -8) {\textbf{order\_items}\\product\_id, variant\_id};

    % Notificaciones
    \node[ent] (not) at (8.6, 0) {\textbf{notifications}\\account\_id, type\\body, link, dismissed};

    % Blog
    \node[ent] (bp) at (0, 3.4) {\textbf{blog\_posts}\\title, slug, body\\status, published\_at};
    \node[ent] (bcat) at (-6, 3.4) {\textbf{blog\_categories}\\name, slug};
    \node[ent] (bpt) at (4.6, 3.4) {\textbf{blog\_post\_tags}\\post\_id, tag\_id};
    \node[ent] (btag) at (8.6, 3.4) {\textbf{blog\_tags}\\name, slug};

    % Relaciones
    \draw[rel] (users) -- node[card, pos=0.15]{1} node[card, pos=0.85]{0..1} node[above, font=\tiny, pos=0.5]{cuenta personal} (acc);
    \draw[rel] (users) -- node[card, pos=0.15]{1} node[card, pos=0.85]{N} (mem);
    \draw[rel] (acc) -- node[card, pos=0.15]{1} node[card, pos=0.85]{N} (mem);
    \draw[rel] (mem) -- node[card, pos=0.15]{N} node[card, pos=0.85]{1} (roles);
    \draw[rel] (roles) -- node[card, pos=0.15]{1} node[card, pos=0.85]{N} (rp);
    \draw[rel] (acc) -- node[card, pos=0.1]{1} node[card, pos=0.9]{N} (inv);
    \draw[rel] (inv.west) -- node[card, pos=0.15]{N} node[card, pos=0.85]{1} (roles.east |- inv.west);
    \draw[rel] (acc) -- node[card, pos=0.15]{1} node[card, pos=0.85]{N} (bc);
    \draw[rel] (bc) -- node[card, pos=0.15]{1} node[card, pos=0.85]{N} (sub);
    \draw[rel] (sub) -- node[card, pos=0.15]{1} node[card, pos=0.85]{N} (si);
    \draw[rel] (bc.east) -| node[card, pos=0.1]{1} node[card, pos=0.9]{N} (ord.north);
    \draw[rel] (ord) -- node[card, pos=0.15]{1} node[card, pos=0.85]{N} (oi);
    \draw[rel] (acc) -- node[card, pos=0.08]{1} node[card, pos=0.92]{N} (not);
    \draw[rel] (bp) -- node[card, pos=0.15]{N} node[card, pos=0.8]{0..1} node[right, font=\tiny, pos=0.5]{autor} (acc);
    \draw[rel] (bp) -- node[card, pos=0.12]{N} node[card, pos=0.88]{0..1} (bcat);
    \draw[rel] (bp) -- node[card, pos=0.15]{1} node[card, pos=0.85]{N} (bpt);
    \draw[rel] (bpt) -- node[card, pos=0.15]{N} node[card, pos=0.85]{1} (btag);
  \end{tikzpicture}}
  \caption{Modelo conceptual de las entidades de la plataforma (esquema
    \texttt{public}). La entidad de trazo discontinuo pertenece al servicio
    de autenticación.}
  \label{fig:ana-datos-public}
  % Texto alternativo: la cuenta (accounts) es el centro del modelo. Cada
  % usuario de auth.users tiene a lo sumo una cuenta personal con su mismo
  % identificador y puede pertenecer a varias cuentas mediante
  % accounts_memberships, que asigna un rol (roles) con sus permisos
  % (role_permissions). Las cuentas tienen invitaciones (con un rol),
  % notificaciones y clientes de facturación, de los que cuelgan
  % suscripciones con sus líneas y pedidos con sus líneas. Las entradas del
  % blog tienen una categoría opcional, un autor (cuenta) opcional y varias
  % etiquetas a través de blog_post_tags.
\end{figure}

La \textbf{cuenta} (\texttt{accounts}) es la unidad de aislamiento
(\textit{tenant}): toda entidad de negocio se enlaza a ella mediante una
columna \texttt{account\_id}. Una cuenta es personal (una por usuario, con su
mismo identificador) o de equipo (con un propietario principal y varios
miembros). Los roles forman una jerarquía mediante
\texttt{hierarchy\_level} y cada uno tiene un conjunto de permisos por
funcionalidad (\texttt{roles.manage}, \texttt{billing.manage},
\texttt{settings.manage}, \texttt{members.manage} e \texttt{invites.manage}).
Por defecto hay tres roles: propietario (\textit{owner}), con todos los
permisos; encargado (\textit{manager}), que gestiona al personal y las
invitaciones pero no la facturación; y miembro (\textit{member}) (ADR-024).
El contenido del blog es la única información pensada para lectura
anónima.

\begin{figure}[htbp]
  \centering
  \resizebox{\textwidth}{!}{%
  \begin{tikzpicture}[
      font=\scriptsize,
      ent/.style={draw, rectangle, align=left, inner sep=3pt,
                  minimum width=2.9cm, fill=white},
      ext/.style={ent, dashed},
      rel/.style={-},
      card/.style={font=\tiny, inner sep=1pt, fill=white},
      lab/.style={font=\tiny\itshape, fill=white, inner sep=1pt}
    ]
    \node[ext] (users) at (-6, 2.8) {\textbf{auth.users}\\id};
    \node[ent] (cacc) at (-6, 0) {\textbf{cms.accounts}\\auth\_user\_id\\is\_active, preferences};
    \node[ent] (croles) at (0, 0) {\textbf{cms.roles}\\name\\rank (0--100, único)\\valid\_from, valid\_until};
    \node[ent] (cperm) at (6, 0) {\textbf{cms.permissions}\\permission\_type\\schema, table, column\\system\_resource, action};
    \node[ent] (cgroups) at (3, 3.4) {\textbf{cms.permission\_groups}\\name, valid\_until};
    \node[ent] (audit) at (-6, -3.2) {\textbf{cms.audit\_logs}\\operation, schema, table\\old\_data, new\_data\\severity, actor\_*};
    \node[ent] (dash) at (0, -3.2) {\textbf{cms.dashboards}\\name, created\_by};
    \node[ent] (widg) at (0, -5.8) {\textbf{cms.dashboard\_widgets}\\widget\_type, config\\posición en la cuadrícula};
    \node[ent] (meta) at (6, -3.2) {\textbf{cms.table\_metadata}\\schema\_name, table\_name\\display\_name, description\\ui\_config (área, ficha)};
    \node[ent] (views) at (6, -5.8) {\textbf{cms.saved\_views}\\filtros, orden, columnas};
    \node[ent] (conf) at (-6, -5.8) {\textbf{cms.configuration}\\key, value\\(p.\,ej. requires\_mfa)};

    \draw[rel] (users) -- node[card, pos=0.2]{1} node[card, pos=0.8]{0..1} (cacc);
    \draw[rel] (cacc) -- node[card, pos=0.1]{N} node[card, pos=0.9]{M} node[lab, above, pos=0.5]{account\_roles} (croles);
    \draw[rel] (croles) -- node[card, pos=0.1]{N} node[card, pos=0.9]{M} node[lab, above, pos=0.5]{role\_permissions} (cperm);
    \draw[rel] (croles.north) |- node[card, pos=0.1]{N} node[card, pos=0.9]{M} node[lab, pos=0.25, left]{role\_permission\_groups} (cgroups.west);
    \draw[rel] (cgroups.east) -| node[card, pos=0.1]{N} node[card, pos=0.9]{M} node[lab, pos=0.75, right]{permission\_group\_permissions} (cperm.north);
    \draw[rel] (cacc.south east) to[bend right=12] node[card, pos=0.08]{N} node[card, pos=0.92]{M}
      node[lab, below, pos=0.5]{account\_permissions (concesión o denegación)} (cperm.south west);
    \draw[rel] (audit) -- node[card, pos=0.2]{N} node[card, pos=0.8]{0..1} (cacc);
    \draw[rel] (dash.north west) -- node[card, pos=0.2]{N} node[card, pos=0.8]{1} (cacc.south east);
    \draw[rel] (dash) -- node[card, pos=0.2]{1} node[card, pos=0.8]{N} (widg);
    \draw[rel] (dash.north) -- node[card, pos=0.2]{N} node[card, pos=0.8]{M} node[lab, right, pos=0.5]{compartido con} (croles.south);
    \draw[rel] (meta) -- node[card, pos=0.2]{1} node[card, pos=0.8]{N} (views);
  \end{tikzpicture}}
  \caption{Modelo conceptual del esquema \texttt{cms}: control de acceso del
    personal, auditoría, paneles y metadatos de las tablas administradas.}
  \label{fig:ana-datos-cms}
  % Texto alternativo: cada usuario puede tener una cuenta del CMS. Las
  % cuentas tienen roles (con rango único) y los roles tienen permisos,
  % directamente o mediante grupos de permisos; además, una cuenta puede
  % recibir concesiones o denegaciones individuales. La auditoría, los
  % paneles (con sus widgets y compartidos con roles) y las vistas guardadas
  % se relacionan con las cuentas y con los metadatos de tabla, que guardan
  % el nombre visible, la descripción, el área de navegación y el diseño de
  % la ficha de cada tabla. La configuración guarda opciones globales, como
  % la obligatoriedad del MFA.
\end{figure}

El esquema \texttt{cms} añade un segundo modelo de autorización, pensado para
el personal y no para los clientes. Sus permisos son de tres clases:
\emph{de datos} (acción sobre un esquema, una tabla o una columna, con
comodines), \emph{de sistema} (por ejemplo, gestionar miembros o
configuración) y \emph{de almacenamiento} (por contenedor y prefijo de
ruta). Los roles tienen un rango único entre 0 y 100, y todos los permisos,
roles y grupos pueden tener un periodo de vigencia. Los metadatos de tabla
(\texttt{cms.table\_metadata}) son el mecanismo de reutilización del CMS: el
nombre visible, la descripción, el área de la barra lateral y el diseño de
la ficha de cada tabla son \emph{datos}, no código, y se cambian desde la
propia consola (ADR-020).

% ---------------------------------------------------------------------
\section{Reglas de negocio}
\label{sec:ana-reglas}

La tabla~\ref{tab:ana-reglas} recoge las reglas de negocio que el sistema
aplica. Todas se hacen cumplir en la base de datos (restricciones,
\textit{triggers}, políticas o funciones), además de en la interfaz cuando
procede, y cada una está cubierta por pruebas pgTAP.

\begin{longtable}{@{}l>{\raggedright\arraybackslash}p{10.2cm}l@{}}
  \caption{Reglas de negocio.}\label{tab:ana-reglas}\\
  \toprule
  \textbf{ID} & \textbf{Regla} & \textbf{Req.} \\
  \midrule
  \endfirsthead
  \toprule
  \textbf{ID} & \textbf{Regla} & \textbf{Req.} \\
  \midrule
  \endhead
  \bottomrule
  \endfoot
  RN-01 & Cada usuario tiene exactamente una cuenta personal, con su mismo identificador. & RF-05 \\
  RN-02 & Una cuenta de equipo tiene un único propietario principal, que solo cambia mediante una transferencia explícita de la propiedad. & RF-06 \\
  RN-03 & El \textit{slug} de una cuenta de equipo y el de una entrada del blog son únicos. & RF-06, RF-09 \\
  RN-04 & Un miembro no puede invitar, asignar ni modificar un rol de nivel superior al suyo. & RF-06 \\
  RN-05 & Las escrituras destructivas (borrar invitaciones, expulsar miembros, gestionar la facturación) exigen el permiso correspondiente, no la mera pertenencia a la cuenta. & RF-06, RNF-02 \\
  RN-06 & Una invitación caduca a los siete días; cada cuenta puede crear como máximo 20 invitaciones por hora y tener 50 pendientes. & RF-06 \\
  RN-07 & Los datos de suscripciones y pedidos solo los escribe la sincronización con la pasarela de pago; ningún usuario los modifica con su sesión. & RF-07 \\
  RN-08 & La consola de administración y el CMS exigen sesión con segundo factor (\texttt{aal2}). & RF-08, RF-09 \\
  RN-09 & Un super-administrador es automáticamente la raíz del CMS (rol de rango máximo); solo existe un rol raíz y retirar la condición de super-administrador retira el acceso, conservando la cuenta para la auditoría. & RF-08, RF-09 \\
  RN-10 & En el CMS nadie actúa sobre su propia cuenta ni sobre cuentas de rango igual o superior, y las cuentas raíz no se desactivan ni pierden su rol desde la consola. & RF-09 \\
  RN-11 & Una denegación explícita prevalece sobre cualquier concesión, y solo se puede delegar un permiso que se tiene y que no se solapa con ninguna denegación propia. & RF-09, RNF-02 \\
  RN-12 & El MFA del CMS es obligatorio por defecto; solo una cuenta raíz con sesión \texttt{aal2} puede desactivarlo, y el cambio queda auditado. & RF-09, RF-10 \\
  RN-13 & La auditoría solo la escriben las funciones del sistema, con el autor tomado de la sesión; una escritura cuya auditoría falla se deshace. & RF-10 \\
  RN-14 & La auditoría conserva una instantánea del autor aunque se borre su cuenta (ADR-018). & RF-10 \\
  RN-15 & Un \textit{widget} de un panel solo muestra datos de tablas que puede leer quien lo ve, aunque el panel esté compartido. & RF-11 \\
  RN-16 & En la web pública solo es visible una entrada del blog publicada cuya fecha de publicación ya se ha alcanzado. & RF-01 \\
  RN-17 & Los esquemas del sistema (\texttt{auth}, catálogos de PostgreSQL y otros protegidos) nunca se exploran ni se buscan desde el CMS. & RF-09, RNF-02 \\
\end{longtable}


\newpage{\ }
\cleardoublepage

\chapter{Diseño del Sistema}
% =====================================================================
% Capítulo: Diseño del Sistema
% Fuentes: borradores/05-diseno.tex, docs/tfg/DECISIONES.md (ADR-001…021),
% AGENTS.md, apps/web/AGENTS.md, apps/web/supabase/AGENTS.md,
% apps/web/src/config/*.config.ts, packages/cms/*,
% packages/function-middleware. Estado de las fuentes: F3c cerrada.
% =====================================================================

Este capítulo describe el diseño de PymeKit: la arquitectura general, el
modelo multi-tenant y de control de acceso, la estrategia de seguridad en
profundidad, la integración con la pasarela de pago, el diseño de la interfaz
y los mecanismos de configuración que hacen la plataforma reutilizable. Cada
decisión remite al registro de decisiones de arquitectura (ADR,
\textit{Architecture Decision Record}~\cite{nygard2011adr}) del proyecto, que
se resume en la sección~\ref{sec:dis-decisiones}. La infraestructura de
ejecución (contenedores, nube y puesta en marcha) se trata en el capítulo de
Despliegue del Sistema.

% ---------------------------------------------------------------------
\section{Arquitectura general}
\label{sec:dis-arquitectura}

\subsection{Monorepo de paquetes}
\label{sec:dis-monorepo}

El código se organiza en un único \textit{monorepo} (un repositorio con
varias aplicaciones y paquetes que se versionan juntos) gestionado con
Turborepo~\cite{turborepo} y pnpm (ADR-002):

\begin{itemize}
  \item \texttt{apps/web}: la aplicación web, que incluye la web pública, el
        área privada, la consola de administración y el CMS;
  \item \texttt{apps/web/supabase}: esquemas, migraciones, datos de prueba y
        pruebas pgTAP de \emph{toda} la base de datos;
  \item \texttt{apps/e2e}: las pruebas E2E con Playwright~\cite{playwright};
  \item \texttt{packages/*}: los paquetes reutilizables de la plataforma
        (\texttt{@pymekit/*}: interfaz, cliente de base de datos,
        \textit{middleware}, autenticación, cuentas, equipos, facturación,
        administración, i18n\ldots);
  \item \texttt{packages/cms/*}: los paquetes del CMS
        (\texttt{@pymekit/cms-*}), separados en paquetes de servidor
        (servicios y rutas de la API) y paquetes de cliente (componentes de
        las pantallas).
\end{itemize}

El monorepo contiene 51 paquetes, 27 de ellos del CMS, sin ciclos de
dependencias (RNF-04). Cada paquete declara sus dependencias de forma
explícita y separa las exportaciones de cliente y de servidor. Esta
separación evita que código con credenciales de servidor acabe en el
navegador ---la \textit{build} de la web falla si el cliente de base de datos
del servidor aparece en un fragmento destinado al navegador--- y, como se
comprobó en la incidencia B-16, también evita ciclos entre paquetes.

\subsection{Capas y pila tecnológica}
\label{sec:dis-capas}

La aplicación sigue una arquitectura en capas en la que cada capa solo
depende de la inferior y la base de datos es la última barrera de
autorización:

\begin{enumerate}
  \item \textbf{Presentación.} Rutas por fichero de TanStack
        Router~\cite{tanstack-router} con React 19, renderizadas en el
        servidor (SSR). Componentes accesibles de \texttt{@pymekit/ui}
        con Tailwind CSS. La caché de datos del cliente usa TanStack
        Query~\cite{tanstack-query}.
  \item \textbf{Aplicación.} \textit{Server functions} de TanStack
        Start~\cite{tanstack-start}: funciones que se ejecutan en el servidor
        y se invocan desde el cliente con tipos compartidos, validadas con
        Zod y protegidas por una cadena de \textit{middleware} (RF-02). Para el
        CMS, una API HTTP construida con Hono~\cite{hono}.
  \item \textbf{Servicios.} Lógica de negocio en servicios sin dependencia de
        la interfaz, con el cliente de datos inyectado: el cliente de
        Supabase~\cite{supabase-docs} en la web y Drizzle ORM~\cite{drizzle}
        en el CMS.
  \item \textbf{Datos.} PostgreSQL con políticas de seguridad a nivel de fila
        (RLS), funciones y \textit{triggers}; servicio de autenticación con
        JWT y MFA; almacenamiento de ficheros.
\end{enumerate}

Para que el conjunto sea coherente se usa una sola pila de librerías en toda
la web, CMS incluido: un único sistema de formularios, uno de
internacionalización y una sola biblioteca de componentes (ADR-013). Así cada
cosa se hace de una única forma (RNF-05).

\subsection{Una única base de datos}
\label{sec:dis-bd-unica}

La aplicación y el CMS comparten \textbf{una sola instancia} de PostgreSQL
(ADR-002). Los datos de la plataforma viven en el esquema \texttt{public} y
los del CMS en un esquema propio, \texttt{cms} (ADR-012). Compartir la base de
datos es lo que permite que el CMS administre los datos de la propia
plataforma ---incluido el blog que publica la web (ADR-017)--- y que las
mismas políticas de seguridad protejan ambas partes. Los esquemas
declarativos se numeran (del \texttt{00} al \texttt{55}) para fijar el orden
de creación, y las migraciones se generan a partir de ellos.

\subsection{CMS integrado en la consola de administración}
\label{sec:dis-cms-integrado}

La decisión de mayor calado del diseño es cómo integrar el CMS (ADR-011). Un
CMS de este tipo se construye habitualmente como una aplicación de página
única con su propio enrutador y una API independiente. Se evaluaron tres opciones:

\begin{enumerate}[label=(\alph*)]
  \item \textbf{aplicaciones separadas} (la web, la interfaz del CMS y su
        API como tres servicios): la de menor esfuerzo, pero con tres
        despliegues, dos inicios de sesión y las \textit{cookies} de sesión
        duplicadas;
  \item \textbf{integración híbrida} (la API dentro de la web y la interfaz
        del CMS servida como aplicación estática): convivirían dos
        enrutadores y dos formas de programar;
  \item \textbf{integración total}: la API se monta dentro del servidor de
        la web y las pantallas se reescriben como rutas de la web.
\end{enumerate}

Se eligió la opción (c), aun siendo la de mayor esfuerzo estimado, porque da
un único servicio, un único inicio de sesión y un único estilo de código:

\begin{itemize}
  \item La aplicación Hono del CMS se monta en \texttt{/api/cms/*} mediante
        una ruta de servidor de la web que le reenvía la petición tal cual.
        El CMS comparte así proceso, dominio y \textit{cookies} con la web.
  \item Las pantallas del CMS son rutas de la web bajo \texttt{/admin/cms/*},
        con \textit{loaders} y TanStack Query sobre un cliente RPC tipado de
        la API.
  \item Durante el SSR, los \textit{loaders} no hacen una petición de red a la
        propia API: invocan la aplicación Hono dentro del mismo proceso,
        reenviando solo la cabecera \texttt{cookie} (ADR-016).
\end{itemize}

La figura~\ref{fig:dis-arquitectura} resume la arquitectura resultante.

\begin{figure}[htbp]
  \centering
  \resizebox{\textwidth}{!}{%
  \begin{tikzpicture}[
      font=\footnotesize,
      box/.style={draw, rounded corners, align=center, minimum height=0.9cm,
                  inner sep=4pt, fill=white},
      group/.style={draw, dashed, rounded corners, inner sep=8pt},
      arr/.style={-{Stealth[length=2mm]}}
    ]
    \node[box, minimum width=11cm] (browser) at (0, 5.2)
      {Navegador (React 19): web pública, área privada, consola \texttt{/admin} y CMS \texttt{/admin/cms}};

    \node[box, minimum width=3.0cm] (ssr) at (-3.8, 3.3)
      {Rutas y \textit{loaders}\\(SSR, TanStack Router)};
    \node[box, minimum width=3.0cm] (sfn) at (0, 3.3)
      {\textit{Server functions}\\+ \textit{middleware}};
    \node[box, minimum width=3.0cm] (hono) at (3.8, 3.3)
      {Ruta \texttt{/api/cms/*}\\$\rightarrow$ aplicación Hono};
    \node[box, minimum width=3.0cm] (svc) at (-1.9, 1.6)
      {Servicios + cliente Supabase\\(sesión del usuario)};
    \node[box, minimum width=3.0cm] (drizzle) at (3.8, 1.6)
      {Servicios Drizzle\\(\textit{claims} del usuario)};
    \node[box, minimum width=2.4cm] (wh) at (-6.6, 1.6)
      {Ruta \texttt{/api/billing}\\\textit{webhook}};

    \begin{scope}[on background layer]
      \node[group, fit=(ssr)(sfn)(hono)(svc)(drizzle)(wh),
            label={[font=\small\bfseries, anchor=north east]south east:\texttt{apps/web} (TanStack Start, un único proceso)}] (web) {};
    \end{scope}

    \node[box, minimum width=2.6cm] (auth) at (-3.8, -0.6) {Auth\\(JWT, MFA)};
    \node[box, minimum width=4.6cm] (pg) at (0.4, -0.6)
      {PostgreSQL + RLS\\esquemas \texttt{public}, \texttt{cms} y \texttt{auth}};
    \node[box, minimum width=2.2cm] (st) at (4.4, -0.6) {Storage};
    \begin{scope}[on background layer]
      \node[group, fit=(auth)(pg)(st),
            label={[font=\small\bfseries]below:Supabase (una única instancia)}] (sb) {};
    \end{scope}

    \node[box] (stripe) at (-10.6, 1.6) {Pasarela\\de pago};

    \draw[arr] (browser) -- (ssr);
    \draw[arr] (browser) -- (sfn);
    \draw[arr] (browser) -- node[right, font=\scriptsize]{RPC} (hono);
    \draw[arr] (ssr) -- (sfn);
    \draw[arr, dashed] (ssr.south) -- ++(0,-0.3) coordinate (ssrq) -- (ssrq -| hono.south) node[midway, below, font=\scriptsize, fill=white, inner sep=1pt]{SSR en proceso} -- (hono.south);
    \draw[arr] (sfn) -- (svc);
    \draw[arr] (hono) -- (drizzle);
    \draw[arr] (svc) -- (pg);
    \draw[arr] (svc) -- (auth);
    \draw[arr] (drizzle) -- (pg);
    \draw[arr] (stripe) -- node[above, font=\scriptsize]{\textit{webhooks}} (wh);
    \draw[arr] (wh) |- (pg.west);
    \draw[arr, dashed] (sfn.north west) to[bend right=25] node[above, font=\scriptsize, pos=0.6]{\textit{checkout}} (stripe.north);
  \end{tikzpicture}}
  \caption{Arquitectura general de PymeKit: una única aplicación web con el
    CMS integrado y una única instancia de la base de datos.}
  \label{fig:dis-arquitectura}
  % Texto alternativo: el navegador se comunica con la aplicación web, que
  % contiene las rutas con SSR, las server functions y la ruta /api/cms que
  % monta la aplicación Hono del CMS. Las server functions usan servicios con
  % el cliente de Supabase y la sesión del usuario; la API del CMS usa
  % servicios Drizzle con los claims del usuario. Ambos acceden a una única
  % instancia (Auth, PostgreSQL con RLS y Storage). La pasarela de pago
  % recibe las sesiones de pago desde las server functions y notifica los
  % cambios a la ruta de webhook, que actualiza la base de datos.
\end{figure}

% ---------------------------------------------------------------------
\section{Modelo multi-tenant y control de acceso}
\label{sec:dis-multitenant}

\subsection{Cuentas como unidad de aislamiento}
\label{sec:dis-cuentas}

PymeKit sigue un modelo multi-tenant de \emph{base de datos y esquema
compartidos}: todos los clientes comparten tablas y cada fila lleva el
identificador de la cuenta a la que pertenece. Es el modelo de menor coste de
operación y el que más simplifica las actualizaciones, a cambio de que el
aislamiento dependa por completo de la corrección del control de
acceso~\cite{bezemer2010multitenant}. Por eso el diseño traslada ese control a la base de
datos.

El \textit{tenant} es la cuenta (sección~\ref{sec:ana-modelo-datos}). La
relación de un usuario con una cuenta se resuelve con un pequeño conjunto de
funciones auxiliares sobre las que se escriben todas las políticas:
\texttt{has\_role\_on\_account} (pertenencia), \texttt{has\_permission}
(permiso concreto) y \texttt{has\_more\_elevated\_role} (jerarquía). Las
escrituras destructivas se condicionan al permiso, no a la mera pertenencia,
y la cuenta activa de cada sesión se guarda en el servidor
(\texttt{user\_active\_account}) para que el área privada sepa en qué
contexto trabaja el usuario.

El modelo admite tres modos, elegidos por configuración
(sección~\ref{sec:dis-configuracion}): solo cuentas personales (B2C), solo
cuentas de equipo (B2B) o ambas. En la web de una sola pyme (ADR-020), los
equipos pueden representar, por ejemplo, las tiendas de la empresa, y sus
clientes usar la cuenta personal.

\subsection{Control de acceso del CMS}
\label{sec:dis-rbac-cms}

El CMS tiene su propio control de acceso basado en roles
(\textit{Role-Based Access Control}, RBAC) en el esquema \texttt{cms}. Conviven
así dos modelos de autorización en la misma base de datos
(tabla~\ref{tab:dis-dos-modelos}). La decisión de diseño (ADR-014) fue
conectarlos en lugar de sustituir uno por otro:

\begin{itemize}
  \item \textbf{El super-administrador es la raíz del CMS.} Unos
        \textit{triggers} sobre los usuarios sincronizan a los
        super-administradores con un rol de sistema \emph{Root}, de rango
        máximo, que concede todos los permisos. Al promocionar a un usuario
        recibe acceso, una cuenta activa del CMS y ese rol; al retirarle la
        condición pierde el acceso y su cuenta se desactiva, pero se conserva
        para la auditoría. Un índice único garantiza que solo exista un rol
        raíz.
  \item \textbf{El RBAC del CMS da acceso limitado a otro personal}
        (soporte, gestión de contenidos). Se descartó dejar solo al
        super-administrador porque incumpliría RF-09 («con permisos
        propios»). Para que una instalación nueva no empiece con el RBAC
        vacío, una migración de datos idempotente deja preparados dos roles
        de partida, editables y sin nadie asignado (ADR-026): «Soporte»,
        que solo lee cuentas, miembros, usuarios y la auditoría, y «Gestor
        de contenidos», que gestiona el blog y los archivos a través de un
        grupo de permisos «Contenidos». Se eligió una migración y no el
        \textit{seed} porque este solo se aplica en desarrollo.
\end{itemize}

\begin{table}[htbp]
  \centering
  \footnotesize
  \caption{Los dos modelos de autorización de PymeKit.}
  \label{tab:dis-dos-modelos}
  \begin{tabularx}{\textwidth}{@{}>{\raggedright\arraybackslash}p{2.6cm}XX@{}}
    \toprule
     & \textbf{Plataforma (\texttt{public})} & \textbf{CMS (\texttt{cms})} \\
    \midrule
    Sujetos & Usuarios de las cuentas (clientes) & Personal con cuenta del CMS \\
    Ámbito & Una cuenta (\textit{tenant}) & Toda la base de datos administrada \\
    Roles & Jerarquía por \texttt{hierarchy\_level} & Rango único 0--100, con vigencia \\
    Permisos & Cinco permisos por funcionalidad & De datos (esquema, tabla, columna), de sistema y de almacenamiento; grupos; concesiones y denegaciones por cuenta \\
    Raíz & Propietario principal de cada cuenta & Rol Root, ligado al super-administrador \\
    MFA & Exigido si el usuario lo tiene configurado & Obligatorio por defecto \\
    \bottomrule
  \end{tabularx}
\end{table}

\subsection{Acceso a la consola}
\label{sec:dis-acceso-consola}

Super-administradores y personal del CMS comparten la consola, con límites
distintos (ADR-016): la consola admite a quien es super-administrador
\emph{o} tiene acceso al CMS; las páginas de gestión de la plataforma exigen
super-administrador y redirigen al personal al CMS. La barra lateral solo
muestra las secciones que la API permite, pero esa ocultación es solo de
interfaz: cada ruta y cada \textit{endpoint} vuelven a comprobar el permiso.

% ---------------------------------------------------------------------
\section{Seguridad en profundidad}
\label{sec:dis-seguridad}

El principio que guía el diseño es que \textbf{la base de datos es la
frontera de autorización} (RNF-02): toda comprobación en la interfaz o en el
servidor de aplicación es una primera barrera, útil para la experiencia de
uso, pero nunca suficiente. Si una capa superior falla, las políticas de la
base de datos siguen impidiendo el acceso, en línea con las recomendaciones
de verificación de seguridad de OWASP~\cite{owasp-asvs}.

\subsection{RLS como frontera}
\label{sec:dis-rls}

PostgreSQL permite asociar a cada tabla políticas de seguridad a nivel de
fila (\textit{Row-Level Security}, RLS)~\cite{postgresql-rls}: predicados que
el propio motor añade a cada consulta según el usuario que la ejecuta.
Supabase~\cite{supabase-rls} traslada la identidad del usuario (su JWT) a la
sesión de PostgreSQL, de modo que las políticas deciden con ella. Las reglas
de diseño, recogidas en las directrices de la base de datos, son:

\begin{itemize}
  \item RLS activo en toda tabla. Se revocan los privilegios por defecto a
        los tres roles de la API (\texttt{anon}, \texttt{authenticated} y
        \texttt{service\_role}), incluidos los que ignoran RLS como
        \texttt{TRUNCATE}, y se conceden solo los verbos que la funcionalidad
        usa.
  \item El \texttt{UPDATE} se concede por columnas, dejando fuera las de
        identidad y de pertenencia (\texttt{id}, \texttt{account\_id}), para
        que nadie pueda mover una fila a otra cuenta.
  \item Las políticas de \texttt{UPDATE} repiten en \texttt{WITH CHECK} las
        comprobaciones del \texttt{INSERT}: sin esa cláusula PostgreSQL solo
        valida la fila antigua~\cite{postgresql-createpolicy} (B-05, B-50).
  \item Toda función \texttt{security definer}, que se ejecuta con
        privilegios elevados y se salta RLS, aplica sus propias
        comprobaciones y las justifica en un comentario; las funciones
        internas no son ejecutables por los roles de la API (B-45).
  \item Las llamadas a funciones dentro de un predicado se envuelven en
        \texttt{(select \ldots)} para que se evalúen una sola vez (B-55).
  \item El cliente de servidor que ignora RLS se reserva para casos
        excepcionales, tras validar a mano los permisos.
\end{itemize}

\subsection{Segundo factor obligatorio en la administración}
\label{sec:dis-mfa}

El servicio de autenticación distingue el nivel de garantía de la sesión
(\textit{Authenticator Assurance Level}, AAL)~\cite{supabase-mfa}:
\texttt{aal1} tras la contraseña y \texttt{aal2} tras completar el segundo
factor. El diseño exige \texttt{aal2} en todas las capas que dan acceso a la
administración:

\begin{itemize}
  \item en la base de datos, la función que identifica al super-administrador
        devuelve falso si la sesión no es \texttt{aal2}; unas políticas
        restrictivas exigen \texttt{aal2} a quien tiene un factor
        configurado; y el acceso al CMS lo decide
        \texttt{cms.verify\_admin\_access()}, que comprueba la cuenta activa
        del CMS y el MFA;
  \item en la interfaz, la consola exige siempre \texttt{aal2}: con un factor
        configurado redirige a la verificación y vuelve a la página pedida;
        sin factor responde 404 (ADR-016).
\end{itemize}

\subsection{Middleware más comprobación en la base de datos}
\label{sec:dis-middleware}

La figura~\ref{fig:dis-flujo-cms} muestra el recorrido de una petición a la
API del CMS y las comprobaciones que encuentra en cada capa. Las
\textit{server functions} de la web siguen el mismo patrón: una cadena de
\textit{middleware} tipado (autenticación, pertenencia al equipo, permiso por
funcionalidad, super-administrador) valida la petición y después el cliente de
Supabase con la sesión del usuario deja que RLS tome la decisión final.

\begin{figure}[htbp]
  \centering
  \resizebox{\textwidth}{!}{%
  \begin{tikzpicture}[
      font=\footnotesize,
      paso/.style={draw, rounded corners, align=center, minimum width=3.3cm,
                   minimum height=1.0cm, inner sep=3pt},
      check/.style={align=left, font=\scriptsize, text width=7.6cm},
      arr/.style={-{Stealth[length=2mm]}}
    ]
    \node[paso] (s1) at (0, 0)     {Navegador\\(consola \texttt{/admin/cms})};
    \node[paso] (s2) at (0, -1.6)  {Ruta de servidor\\\texttt{/api/cms/\$}};
    \node[paso] (s3) at (0, -3.2)  {Aplicación Hono\\(\textit{middleware})};
    \node[paso] (s4) at (0, -4.8)  {Servicio Drizzle\\(transacción)};
    \node[paso] (s5) at (0, -6.4)  {PostgreSQL\\(esquema \texttt{cms})};

    \draw[arr] (s1) -- (s2);
    \draw[arr] (s2) -- (s3);
    \draw[arr] (s3) -- (s4);
    \draw[arr] (s4) -- (s5);

    \node[check, right=0.6cm of s1] {Guarda de ruta: sesión viva y
      \texttt{aal2}, acceso de administración (si no: inicio de sesión,
      verificación MFA o 404)};
    \node[check, right=0.6cm of s2] {Reenvío a Hono; el código del CMS se
      importa solo en el servidor};
    \node[check, right=0.6cm of s3] {Cabeceras de seguridad y CSRF propia; JWT
      válido (401), acceso al CMS (403), usuario no bloqueado (en cerrado);
      validación estricta de la entrada con Zod (400)};
    \node[check, right=0.6cm of s4] {\textit{Claims} del usuario y rol
      \texttt{authenticated} fijados con parámetros enlazados; permiso
      concreto comprobado por el servicio};
    \node[check, right=0.6cm of s5] {\texttt{cms.verify\_admin\_access()}
      (cuenta activa, MFA) y políticas RLS: \textbf{decisión final}};
  \end{tikzpicture}}
  \caption{Flujo de una petición a la API del CMS y comprobaciones de
    seguridad en cada capa.}
  \label{fig:dis-flujo-cms}
  % Texto alternativo: diagrama vertical de cinco pasos (navegador, ruta de
  % servidor, aplicación Hono, servicio Drizzle, PostgreSQL). A la derecha de
  % cada paso se indican sus comprobaciones; la última, la base de datos, es
  % la que toma la decisión.
\end{figure}

Dos detalles del diseño merecen mención. Primero, el servicio fija la
identidad del usuario en la transacción mediante parámetros enlazados, nunca
interpolando valores del JWT en el SQL, de modo que las políticas RLS del
esquema \texttt{cms} se aplican igual que a cualquier otro cliente. Segundo,
un super-administrador con sesión \texttt{aal1} supera el \textit{middleware}
(tiene el \textit{claim} de acceso), pero no ve datos, porque la base de
datos exige el segundo factor: es la defensa en profundidad en acción.

\subsection{Fallo en cerrado y errores sin detalles internos}
\label{sec:dis-fallo-cerrado}

Todas las comprobaciones de seguridad \textbf{fallan en cerrado}
(\textit{fail closed}): ante un dato ausente, ilegible o no válido, se
deniega. El diseño lo aplica a la opción que indica si el CMS exige MFA, que
se lee con una función de privilegios elevados para que RLS no pueda
ocultársela a quien decide (B-06); a la comprobación de usuario bloqueado
cuando el servicio de autenticación no responde (B-23); y a la auditoría,
cuyo fallo deshace la escritura auditada (B-41).

Del mismo modo, los errores de la base de datos se clasifican en el servidor:
el cliente solo recibe un estado HTTP y un código estable
(\texttt{RECORD\_*}, \texttt{SETTINGS\_*}, \texttt{AUDIT\_LOG\_*}\ldots),
nunca el texto de PostgreSQL, que podría revelar nombres de tablas o
restricciones (B-20, B-33).

\subsection{Seguridad del entorno de desarrollo}
\label{sec:dis-seguridad-entorno}

La seguridad alcanza también al entorno de trabajo. Los servicios locales de
la base de datos se aíslan de Internet con una regla de cortafuegos específica
para contenedores y la aplicación se consulta mediante un túnel SSH (ADR-010,
tras la incidencia B-04). Los secretos nunca se versionan.

% ---------------------------------------------------------------------
\section{Integración con la pasarela de pago}
\label{sec:dis-pagos}

La facturación se diseña detrás de una \emph{pasarela abstracta}
(\texttt{@pymekit/billing-gateway}) con una estrategia por proveedor: la
aplicación pide «crear una sesión de pago» o «abrir el portal de cliente» sin
conocer el proveedor concreto, que se elige con la variable
\texttt{VITE\_BILLING\_PROVIDER} y se valida al construir la aplicación.
Actualmente se registra Stripe.

Los planes, sus precios y sus características se declaran en
\texttt{billing.config.ts}, validado con un esquema. El flujo tiene dos
sentidos (figura~\ref{fig:dis-arquitectura}): las \textit{server functions}
crean la sesión de pago tras comprobar el permiso \texttt{billing.manage}; y
la pasarela notifica cada cambio a la ruta \texttt{/api/billing/webhook}, que
verifica la firma y sincroniza clientes, suscripciones y pedidos con el
cliente de servicio. Ningún usuario escribe esas tablas con su sesión
(RN-07): la fuente de verdad de un cobro es la pasarela.

Los planes de partida están pensados para una pyme (ADR-025): «Básico»,
«Profesional» y «Empresa», en euros, cada uno con pago mensual y anual. Los
importes son ficticios y los identificadores de precio de la pasarela se
leen de variables de entorno (\texttt{VITE\_STRIPE\_PRICE\_*}); si faltan,
se usan valores de relleno marcados como ejemplo, de modo que la página de
precios funciona desde el primer arranque y la pyme solo tiene que crear sus
precios en Stripe y copiar sus identificadores. Los textos de cada plan son
claves de traducción.
\todo{Completar con la verificación con la pasarela en modo de pruebas (F6/F7).}

% ---------------------------------------------------------------------
\section{Diseño de la interfaz}
\label{sec:dis-interfaz}

\subsection{Web pública y área privada}

La web pública ofrece inicio, precios, blog, preguntas frecuentes, contacto y
textos legales; el área privada, el panel de la cuenta activa, sus ajustes,
miembros y facturación. La identidad visual es propia (logotipo, iconos y
capturas generadas a partir de la propia aplicación, B-46). La interfaz es
adaptable a dispositivos móviles, admite tema claro y oscuro y se construye
con componentes accesibles (RNF-08).

Mientras PymeKit se presenta como TFG, la página de inicio incluye un bloque
académico (logotipo de la escuela, autor, tutor y un recorrido guiado por la
demostración). Depende de un único interruptor de configuración, de modo que
una pyme que reutilice la plataforma lo desactiva sin tocar el código
(sección~\ref{sec:dis-configuracion}).

\subsection{Consola de administración por áreas de negocio}
\label{sec:dis-consola}

La primera versión de la consola separaba «Plataforma», «CMS» y «Recursos» y
agrupaba las tablas por esquema SQL, una organización técnica que no entiende
quien gestiona una pyme. Tras la revisión del autor se rediseñó como una
\textbf{única consola organizada por áreas de negocio} (ADR-020,
figura~\ref{fig:dis-consola}):

\begin{itemize}
  \item un bloque \textbf{Gestión} con Inicio, Gestión de cuentas (solo
        super-administrador) y las herramientas del CMS: Usuarios, Archivos,
        Paneles, Auditoría y Ajustes;
  \item un bloque \textbf{Datos} con una entrada plegable por área (Blog,
        Cuentas, Facturación, Sistema y Otros datos) que contiene sus tablas
        con nombre en español, y una entrada «Todas las tablas».
\end{itemize}

El área de cada tabla, su nombre visible y su descripción son \emph{datos}
de los metadatos del CMS, editables en Ajustes $\rightarrow$ Recursos, de
modo que una pyme puede crear áreas para sus propias tablas sin programar.
Las tablas de esquemas protegidos no aparecen en la navegación.

\begin{figure}[htbp]
  \centering
  \begin{tikzpicture}[
      font=\scriptsize,
      item/.style={anchor=west, inner sep=1pt},
      head/.style={anchor=west, font=\scriptsize\bfseries, inner sep=1pt}
    ]
    \draw[rounded corners] (0,0) rectangle (4.2,-8.6);
    \node[head] at (0.2,-0.35) {Gestión};
    \node[item] at (0.35,-0.8)  {Inicio};
    \node[item] at (0.35,-1.2)  {Gestión de cuentas$^{*}$};
    \node[item] at (0.35,-1.6)  {Usuarios};
    \node[item] at (0.35,-2.0)  {Archivos};
    \node[item] at (0.35,-2.4)  {Paneles};
    \node[item] at (0.35,-2.8)  {Auditoría};
    \node[item] at (0.35,-3.2)  {Ajustes};
    \node[head] at (0.2,-3.8) {Datos};
    \node[item] at (0.35,-4.25) {$-$ Blog};
    \node[item] at (0.75,-4.65) {Entradas};
    \node[item] at (0.75,-5.0)  {Categorías};
    \node[item] at (0.75,-5.35) {Etiquetas};
    \node[item] at (0.35,-5.8)  {$+$ Cuentas};
    \node[item] at (0.35,-6.2)  {$+$ Facturación};
    \node[item] at (0.35,-6.6)  {$+$ Sistema};
    \node[item] at (0.35,-7.0)  {$+$ Otros datos};
    \node[item] at (0.35,-7.5)  {Todas las tablas};
    \node[item, font=\tiny] at (0.2,-8.25) {$^{*}$ solo super-administrador};

    \draw[rounded corners] (4.5,0) rectangle (11.4,-8.6);
    \node[anchor=west, font=\small\bfseries] at (4.8,-0.45) {Entradas};
    \node[anchor=west, text width=6.3cm, font=\tiny] at (4.8,-1.0)
      {Resumen: «Artículos del blog. Solo los publicados se muestran en la web.»};
    \draw (4.8,-1.5) rectangle (11.1,-2.0);
    \node[anchor=west, font=\tiny] at (4.9,-1.75) {Buscar\ldots \quad Filtros \quad Columnas \quad Vistas};
    \foreach \y in {-2.4,-2.9,-3.4,-3.9,-4.4} { \draw[gray] (4.8,\y) -- (11.1,\y); }
    \node[anchor=west, font=\tiny] at (4.9,-2.65) {Título \hspace{1.6cm} Estado \hspace{0.9cm} Publicada};
    \draw (4.8,-5.0) rectangle (11.1,-8.3);
    \node[anchor=west, font=\tiny\bfseries] at (4.9,-5.3) {Editar entrada};
    \node[anchor=west, font=\tiny] at (4.9,-5.8) {Título \textcircled{\scalebox{0.7}{?}}};
    \draw (4.9,-6.05) rectangle (10.9,-6.4);
    \node[anchor=west, font=\tiny] at (4.9,-6.8) {Estado \textcircled{\scalebox{0.7}{?}}};
    \draw (4.9,-7.05) rectangle (10.9,-7.4);
    \node[anchor=east, font=\tiny] at (10.9,-7.9) {[Guardar]};
  \end{tikzpicture}
  \caption{Esquema de la consola: navegación por bloques y áreas (izquierda)
    y una pantalla del explorador con su resumen, el listado y un formulario
    con ayuda contextual «?» (derecha).}
  \label{fig:dis-consola}
  % Texto alternativo: a la izquierda, la barra lateral con el bloque Gestión
  % (Inicio, Gestión de cuentas solo para el super-administrador, Usuarios,
  % Archivos, Paneles, Auditoría y Ajustes) y el bloque Datos con las áreas
  % plegables Blog (desplegada: Entradas, Categorías, Etiquetas), Cuentas,
  % Facturación, Sistema y Otros datos, y Todas las tablas. A la derecha, la
  % pantalla de Entradas con un resumen bajo el título, la barra de búsqueda
  % y filtros, el listado y un formulario de edición con un icono de ayuda
  % junto a cada campo.
\end{figure}
\todo{Sustituir o complementar la figura~\ref{fig:dis-consola} con capturas
reales de la consola (pueden generarse con el script de capturas de
\texttt{apps/e2e/scripts}).}

\subsection{Ayuda contextual y resúmenes}
\label{sec:dis-ayudas}

Para que la consola la pueda usar personal no técnico, cada pantalla muestra
bajo el título un \textbf{resumen} de una o dos frases que explica para qué
sirve; en el explorador, ese resumen es la descripción de la tabla guardada en
sus metadatos. Además, cada campo de los formularios de la consola lleva una
\textbf{ayuda contextual «?»}, un componente reutilizable de la biblioteca de
interfaz que muestra al pulsarlo qué significa el campo y qué efecto tiene.
Ambos textos pasan por el sistema de traducción.

\subsection{Internacionalización}
\label{sec:dis-i18n}

La interfaz solo ofrece español (ADR-021): PymeKit se dirige a pymes
españolas y mantener dos idiomas duplicaría el coste de cada texto. Sin
embargo, se conserva toda la infraestructura de traducción: ningún texto
visible se escribe en el código ---tampoco atributos como \texttt{aria-label}
o \texttt{placeholder}---, sino que se referencia por clave a los mensajes del
idioma (\texttt{packages/i18n/src/messages/es}); la lista de idiomas está en
la configuración y una prueba unitaria comprueba la coherencia de las claves.
Añadir un idioma consiste en crear su carpeta de mensajes y registrarlo,
y se documenta como \emph{punto de extensión}. Los identificadores, las
tablas y las columnas siguen en inglés, como es habitual en el ecosistema
(ADR-003); sus nombres visibles en español son datos del CMS.

% ---------------------------------------------------------------------
\section{Decisiones de diseño principales}
\label{sec:dis-decisiones}

Cada decisión relevante se registró en el momento de tomarla con su contexto,
alternativas y consecuencias. La tabla~\ref{tab:dis-adr} las resume; las
decisiones no se borran, sino que se marcan como sustituidas cuando dejan de
valer.

\begin{longtable}{@{}l>{\raggedright\arraybackslash}p{4.4cm}>{\raggedright\arraybackslash}p{6.2cm}l@{}}
  \caption{Resumen de decisiones de diseño (ADR) y alternativas descartadas.}\label{tab:dis-adr}\\
  \toprule
  \textbf{ADR} & \textbf{Decisión} & \textbf{Alternativas descartadas (motivo)} & \textbf{Req.} \\
  \midrule
  \endfirsthead
  \toprule
  \textbf{ADR} & \textbf{Decisión} & \textbf{Alternativas descartadas (motivo)} & \textbf{Req.} \\
  \midrule
  \endhead
  \bottomrule
  \endfoot
  001 & Partir de componentes base con licencia
    & Desarrollo completo desde cero (inviable en plazo); otras bases abiertas (menos completas en multi-tenant y RLS); CMS genéricos con otra base de datos u otro \textit{backend} & RNF-01, 04 \\
  002 & Monorepo único y una sola base de datos & Dos repositorios con la misma base de datos (CI, despliegue y documentación duplicados) & RF-09, RNF-04 \\
  003 & Comentarios y documentación en español; identificadores en inglés & Todo en español (choca con el ecosistema); todo en inglés (dificulta la evaluación) & RF-13, RNF-05 \\
  004 & Un único CMS, de datos, sin CMS de contenidos basado en ficheros & Mantener dos CMS (complejidad sin aportar al objetivo) & RF-09 \\
  005 & Plataforma genérica sin dominio de negocio concreto & Implementar un dominio de ejemplo (desvía el esfuerzo de la base reutilizable) & RNF-01 \\
  006 & Flujo de desarrollo sin comprobaciones remotas de licencia & Mantenerla (dependencia de red y de un servicio externo) & --- \\
  007 & La memoria no nombra los componentes base & --- & --- \\
  008 & Descartar el servidor MCP y los generadores de código ligados a documentación externa & Conservarlos renombrados (dependen de documentación no incluida) & RNF-04 \\
  009 & La web escucha en el puerto 3100 & Puerto por defecto (ocupado en el equipo de desarrollo, B-02) & RNF-07 \\
  010 & Aislar de Internet los servicios locales & Confiar en el cortafuegos del sistema (los contenedores se lo saltan, B-04) & RNF-02 \\
  011 & CMS integrado: API Hono en \texttt{/api/cms}, pantallas como rutas de la web & Aplicaciones separadas (tres servicios, dos inicios de sesión); híbrida (dos enrutadores) & RF-08--11 \\
  012 & Esquema SQL del CMS llamado \texttt{cms} & Conservar el nombre original (excepción permanente en los controles automáticos) & RF-09 \\
  013 & Una sola pila de librerías en toda la web & Mantener las del CMS (dos formas de hacer lo mismo; dos sistemas de i18n) & RNF-05 \\
  014 & Super-administrador como raíz del CMS; RBAC del CMS para el resto del personal & Solo super-administrador, sin RBAC granular (incumple RF-09) & RF-08, 09 \\
  015 & Endurecer la seguridad del CMS & Aceptar el código sin auditar (brechas confirmadas, B-05--B-09) & RNF-02, 03 \\
  016 & Consola para super-administrador y personal del CMS, siempre con \texttt{aal2} & Dejar entrar con el \textit{claim} de acceso sin comprobar el AAL (B-14) & RF-08, RNF-02 \\
  017 & Blog en la base de datos gestionado desde el CMS & Blog fuera (debilita la gestión de contenidos); CMS de ficheros (dos CMS) & RF-01, 09 \\
  018 & Instantánea del autor en la auditoría & Claves ajenas \texttt{RESTRICT} (impide borrar usuarios); borrado lógico (más complejo) & RF-10 \\
  019 & Paneles con cuadrícula propia, sin librería & Librería de cuadrícula (más dependencias; arrastrar sin teclado no es accesible) & RF-11, RNF-08 \\
  020 & Consola por áreas de negocio, nombres en español como datos y sin esquema de demostración & Traducir nombres en el cliente; áreas fijas en un fichero; quitar equipos y pagos & RF-09, RNF-01, 08 \\
  021 & Interfaz solo en español conservando la infraestructura i18n & Mantener dos idiomas (doble mantenimiento); eliminar la i18n (impide traducir) & RF-13, RNF-01 \\
  022 & Una sola imagen Docker para la web y Supabase externa & Contenedores separados para la API del CMS; Supabase autoalojada en el mismo \textit{compose} & RNF-04, 07 \\
  023 & Comentarios didácticos en español limitados al núcleo de la arquitectura & Reescribir todos los comentarios (coste desproporcionado); no reescribir ninguno & RNF-05 \\
  024 & Rol intermedio «Encargado» en los equipos & Mantener solo propietario y miembro; dar al miembro más permisos & RF-06, RNF-01 \\
  025 & Planes de ejemplo para pymes en euros con ids de precio por variables & Ids fijos en el fichero; planes sin precios & RF-07, RNF-01 \\
  026 & Plantilla de roles del CMS en una migración de datos idempotente & Solo en el \textit{seed} (no llega a producción); roles de sistema no editables & RF-09, RNF-01 \\
\end{longtable}

% ---------------------------------------------------------------------
\section{Configuración reutilizable}
\label{sec:dis-configuracion}

La reutilización por configuración (RNF-01) se apoya en tres mecanismos, de
menor a mayor dinamismo.

\paragraph{Ficheros de configuración tipados.} Todo lo que cambia de un
producto a otro se concentra en \texttt{apps/web/src/config/*.config.ts}
(tabla~\ref{tab:dis-config}). Cada fichero valida su contenido con un esquema
Zod al arrancar, de modo que una configuración incorrecta falla pronto y con
un mensaje claro, en lugar de en tiempo de ejecución.

\begin{table}[htbp]
  \centering
  \footnotesize
  \caption{Ficheros de configuración de la aplicación.}
  \label{tab:dis-config}
  \begin{tabularx}{\textwidth}{@{}>{\raggedright\arraybackslash}p{4.4cm}X@{}}
    \toprule
    \textbf{Fichero} & \textbf{Qué configura} \\
    \midrule
    \texttt{app.config.ts} & Nombre del producto, título, descripción, URL y tema por defecto \\
    \texttt{auth.config.ts} & Métodos de acceso (contraseña, enlace mágico, OAuth) y captcha \\
    \texttt{account-mode.config.ts} & Modo de cuentas: personal, equipos o híbrido \\
    \texttt{feature-flags.config.ts} & Interruptores de funcionalidad derivados del modo y de las variables de entorno \\
    \texttt{billing.config.ts} & Proveedor de pago, productos, planes y precios \\
    \texttt{paths.config.ts} & Rutas de la aplicación \\
    \texttt{navigation.config.tsx}, \texttt{marketing-navigation.config.tsx} & Menús del área privada y de la web pública \\
    \texttt{tfg.config.ts} & Bloque académico de la página de inicio \\
    \bottomrule
  \end{tabularx}
\end{table}

\paragraph{Variables de entorno.} Los valores que dependen de la instalación
se leen de variables de entorno. Las públicas llevan el prefijo
\texttt{VITE\_} y los secretos (claves de la base de datos y de la pasarela)
nunca. Los interruptores de los módulos opcionales son:

\begin{itemize}
  \item \textbf{Equipos}: \texttt{VITE\_ACCOUNT\_MODE}
        (\texttt{personal-only}, \texttt{organizations-only} o
        \texttt{hybrid}) decide qué superficies de cuenta existen, y
        \texttt{VITE\_ENABLE\_TEAM\_ACCOUNTS\_CREATION} y
        \texttt{\ldots\_DELETION} las afinan. El modo es la única fuente de
        verdad: un interruptor granular nunca puede reactivar una superficie
        que el modo desactiva.
  \item \textbf{Facturación}: \texttt{VITE\_BILLING\_PROVIDER} y
        \texttt{VITE\_ENABLE\_PERSONAL\_ACCOUNT\_BILLING} /
        \texttt{VITE\_ENABLE\_TEAM\_ACCOUNTS\_BILLING}; la base de datos
        mantiene el mismo interruptor en su tabla de configuración para que
        las políticas RLS también lo respeten.
  \item \textbf{Bloque TFG}: \texttt{VITE\_SHOW\_TFG\_INFO=false} deja la web
        limpia de cualquier referencia académica.
  \item \textbf{Otros}: notificaciones (y su modo en tiempo real), selector
        de tema, baja de cuentas y métodos de autenticación.
\end{itemize}

\paragraph{Configuración como datos.} Lo que la pyme debe poder cambiar sin
desplegar se guarda en la base de datos y se edita desde la consola: el
nombre visible, la descripción y el área de cada tabla, el formato de sus
columnas y el diseño de sus fichas (metadatos del CMS); los roles y permisos
del personal; y las preferencias globales, como la obligatoriedad del MFA.

Finalmente, el monorepo incluye generadores para crear paquetes nuevos con la
estructura y las dependencias correctas, de modo que la lógica propia de un
producto se añade como un paquete más sin modificar el núcleo. El
generador \texttt{package} (\texttt{pnpm turbo gen package}) pregunta si se
trata de una funcionalidad de negocio (\texttt{packages/features}) o de un
paquete de infraestructura y crea su \texttt{package.json},
\texttt{tsconfig.json} y punto de entrada, tomando las versiones del catálogo
del monorepo. Los generadores heredados para otras plataformas de despliegue
se retiraron porque producían configuración de otro \textit{framework}. Toda
esta configuración ---ficheros, variables de entorno, interruptores de
módulos, identidad visual, roles, planes y tablas propias en el CMS--- está
documentada en la guía \texttt{docs/CONFIGURACION.md} del repositorio.


\newpage{\ }
\cleardoublepage

\chapter{Construcción del Sistema}
% =====================================================================
% Capítulo: Construcción del Sistema
% Fuentes: docs/tfg/{PLAN,PROGRESO,BITACORA,DECISIONES,TRAZABILIDAD}.md,
% MAPA-REFERENCIAS.md (solo como guía; no se nombran las referencias),
% AGENTS.md, .claude/{settings.json,hooks,skills,agents},
% apps/web/supabase/{schemas,migrations}, packages/cms/*,
% packages/function-middleware, apps/web/src. Estado: F3c cerrada.
% =====================================================================

% Estilo de los extractos de código de este capítulo: letra pequeña, marco
% y sustitución de los caracteres acentuados (listings no admite UTF-8
% multibyte de forma nativa).
\lstdefinestyle{pymekit}{
  basicstyle=\ttfamily\scriptsize,
  breaklines=true,
  frame=single,
  columns=fullflexible,
  keepspaces=true,
  showstringspaces=false,
  numbers=left,
  numberstyle=\tiny,
  numbersep=6pt,
  xleftmargin=1.2em,
  literate={á}{{\'a}}1 {é}{{\'e}}1 {í}{{\'i}}1 {ó}{{\'o}}1 {ú}{{\'u}}1
           {Á}{{\'A}}1 {É}{{\'E}}1 {Í}{{\'I}}1 {Ó}{{\'O}}1 {Ú}{{\'U}}1
           {ñ}{{\~n}}1 {Ñ}{{\~N}}1 {ü}{{\"u}}1 {¿}{{?`}}1 {¡}{{!`}}1
           {«}{{\guillemotleft}}1 {»}{{\guillemotright}}1
}

Este capítulo describe cómo se ha construido PymeKit: el entorno y las
herramientas, la estructura del código, qué parte se ha reutilizado y qué
parte es propia, la secuencia de incrementos (fases F0 a F3c), una selección
de extractos de código representativos, el uso de agentes de IA en el
desarrollo y los problemas más relevantes que se resolvieron.

% ---------------------------------------------------------------------
\section{Entorno y herramientas}
\label{sec:imp-entorno}

El desarrollo se ha realizado en un servidor Linux remoto, accedido por SSH,
con las herramientas de la tabla~\ref{tab:imp-herramientas}. La base de datos
local (PostgreSQL, autenticación y almacenamiento) se ejecuta en contenedores
con la CLI de Supabase; como el servidor tiene IP pública, sus puertos se
filtran explícitamente y la aplicación se consulta a través de un túnel SSH
(ADR-010).

\begin{table}[htbp]
  \centering
  \footnotesize
  \caption{Herramientas del entorno de desarrollo.}
  \label{tab:imp-herramientas}
  \begin{tabularx}{\textwidth}{@{}>{\raggedright\arraybackslash}p{3.4cm}X@{}}
    \toprule
    \textbf{Ámbito} & \textbf{Herramienta} \\
    \midrule
    Lenguaje y ejecución & TypeScript en modo estricto, Node.js (versión LTS) \\
    Monorepo & pnpm (con catálogo de versiones compartido) y Turborepo~\cite{turborepo} \\
    Aplicación & TanStack Start, Router y Query; React 19; Vite; Tailwind CSS 4 \\
    CMS & Hono, Drizzle ORM, Zod \\
    Base de datos & Supabase CLI (PostgreSQL, Auth, Storage) en contenedores Docker \\
    Pagos & Stripe en modo de pruebas y su CLI para reenviar \textit{webhooks} \\
    Pruebas & Vitest (unitarias), pgTAP~\cite{pgtap} (base de datos), Playwright~\cite{playwright} (E2E) \\
    Calidad & ESLint, Prettier, comprobación de tipos y de dependencias del monorepo \\
    Control de versiones & Git, GitHub (repositorio privado) y GitHub Actions; mensajes según Conventional Commits~\cite{conventionalcommits} \\
    Asistencia de IA & Claude Code~\cite{claudecode} con un \textit{harness} propio (sección~\ref{sec:imp-agentes}) \\
    Memoria & \LaTeX{} con la plantilla oficial de la escuela \\
    \bottomrule
  \end{tabularx}
\end{table}

La integración continua (\texttt{.github/workflows/ci.yml}) tiene tres
trabajos: \emph{Controles del TFG} (control de marca, bloqueante, e informe de
comentarios en inglés), \emph{TypeScript} (tipos, \textit{lint}, formato y
pruebas unitarias en cada \textit{push} a la rama principal y en cada \textit{pull request}) y
\emph{Supabase + E2E} (base de datos local, pgTAP, \textit{build} de pruebas y
Playwright), que por su coste solo se ejecuta en las \textit{pull requests}
hacia la rama principal o a mano.

% ---------------------------------------------------------------------
\section{Estructura del código}
\label{sec:imp-estructura}

La figura~\ref{fig:imp-arbol} muestra la estructura del repositorio. La
aplicación web es la única aplicación desplegable; todo lo demás son paquetes
que la componen o herramientas de verificación.

\begin{figure}[htbp]
  \centering
\begin{lstlisting}[style=pymekit, numbers=none, frame=none]
apps/
  web/                      TanStack Start + React 19 (una sola aplicación)
    src/routes/             rutas por fichero
      _marketing/           web pública: inicio, precios, blog, legales
      _authenticated/       área privada (sesión viva y MFA)
      admin/                consola: plataforma (super-admin) y cms/ (CMS)
      api/                  rutas de servidor: cms/$, billing/webhook, db/...
    src/lib/                server functions (*.functions.ts), servicios (*.server.ts)
    src/config/             configuración reutilizable (*.config.ts)
    supabase/               schemas/ (00-55), migrations/, tests/database (pgTAP)
  e2e/                      pruebas Playwright de la web y del CMS
packages/
  ui, i18n, shared, supabase, function-middleware, policies, otp, ...
  features/{auth,accounts,team-accounts,admin,notifications}
  billing/{core,gateway,stripe}
  cms/                      27 paquetes @pymekit/cms-*: api, auth, data-explorer,
                            filters, query-builder, audit-logs, dashboards, ...
                            y sus paquetes cliente *-ui
docs/tfg/                   plan, requisitos, decisiones, bitácora, trazabilidad
scripts/tfg/                controles del TFG (marca, comentarios, etiquetas)
.claude/                    harness de desarrollo asistido (skills, agentes, hooks)
memoria/                    esta memoria
\end{lstlisting}
  \caption{Estructura simplificada del repositorio.}
  \label{fig:imp-arbol}
\end{figure}

Las convenciones más importantes, recogidas en las directrices del
repositorio (\texttt{AGENTS.md}), son:

\begin{itemize}
  \item las \textit{server functions} se definen en ficheros
        \texttt{*.functions.ts} con la cadena
        \texttt{.middleware(\ldots).validator(Esquema).handler(\ldots)};
  \item los módulos solo de servidor llevan el sufijo \texttt{.server.ts} y
        un \textit{plugin} de Vite impide que lleguen al navegador;
  \item en el CMS, cada funcionalidad es un paquete con un servicio Drizzle
        (clase y fábrica \texttt{createXService()}), sus rutas Hono y
        funciones puras en \texttt{utils/};
  \item ningún texto visible se escribe en el código (todo pasa por i18n), y
        los registros usan el \textit{logger} compartido, nunca
        \texttt{console.log};
  \item los puntos clave se marcan con comentarios \texttt{[TFG] RF-xx} que
        un script recopila en la matriz de trazabilidad (275 etiquetas en su
        última generación).
\end{itemize}

% ---------------------------------------------------------------------
\section{Reutilización: qué se ha heredado y qué es propio}
\label{sec:imp-reutilizacion}

% PENDIENTE (P-06): redacción neutra; confirmar con el tutor la forma de
% presentar la base reutilizada sin nombrarla (ADR-007).
Como se decidió al planificar el proyecto (ADR-001), PymeKit no se ha escrito
desde cero. Parte de dos bases de código existentes, utilizadas conforme a su
licencia: una \textbf{base SaaS} sobre la misma pila tecnológica (aplicación
web, autenticación, cuentas y equipos, facturación y panel de
super-administración) y un \textbf{CMS de administración de datos} para el
mismo tipo de base de datos (aplicación de página única con su propia API y
su esquema de base de datos). Ambas se han \emph{integrado} en un único
producto, \emph{adaptadas} a sus requisitos y \emph{extendidas} con
funcionalidad nueva. El origen está documentado en el repositorio, que es
privado, y todo el código importado ha seguido el mismo flujo: copiar,
retirar las marcas originales, comentar en español, verificar y registrar.

% PENDIENTE (P-06)
La tabla~\ref{tab:imp-reutilizacion} distingue las tres categorías. La
aportación del TFG no está en reescribir lo que ya existía, sino en la
integración de dos sistemas con arquitecturas distintas, en la corrección de
su seguridad, en la adaptación al caso de una pyme y en la verificación y la
documentación del conjunto.

\begin{longtable}{@{}>{\raggedright\arraybackslash}p{2.0cm}>{\raggedright\arraybackslash}p{12.2cm}@{}}
  \caption{Código reutilizado, adaptado y propio.}\label{tab:imp-reutilizacion}\\
  \toprule
  \textbf{Tipo} & \textbf{Elementos} \\
  \midrule
  \endfirsthead
  \toprule
  \textbf{Tipo} & \textbf{Elementos} \\
  \midrule
  \endhead
  \bottomrule
  \endfoot
  Reutilizado & Biblioteca de componentes de interfaz; clientes de base de datos; \textit{middleware} de las \textit{server functions}; módulos de autenticación, cuentas, equipos, notificaciones, correo y facturación; esquema \texttt{public} de cuentas, roles, invitaciones y facturación con sus pruebas pgTAP; servicios y rutas de la API del CMS; utilidades del CMS (filtros, constructor de consultas, formateadores). \\
  \midrule
  Adaptado & Aplicación web (sin el CMS de contenidos basado en ficheros, el registro de cambios ni la documentación de la base; con comprobación de sesión viva); esquema del CMS renombrado a \texttt{cms} y endurecido (ADR-012, ADR-015); migración de unos 64 formularios y unos 183 ficheros con textos del CMS a la pila de la web (ADR-013); reescritura de las aproximadamente treinta pantallas del CMS como rutas de la web; pruebas E2E heredadas; internacionalización (solo español, ADR-021). \\
  \midrule
  Propio & Integración del CMS en la web (montaje de la API, llamada en proceso durante el SSR, guardas de la consola; ADR-011, ADR-016); pegamento entre super-administrador y rol raíz del CMS (ADR-014); corrección de 38 fallos heredados, 25 de ellos de seguridad (bitácora); blog en la base de datos y rutas públicas (ADR-017); instantánea del autor en la auditoría (ADR-018); cuadrícula propia de los paneles (ADR-019); consola por áreas de negocio, nombres en español como datos, ayuda contextual «?» y resúmenes (ADR-020); bloque académico configurable; identidad visual; 25 de las 26 migraciones añadidas durante el proyecto (la primera incorpora el esquema del CMS de partida) y las pruebas pgTAP de aislamiento escritas en las revisiones de RLS; el \textit{harness} de desarrollo y los controles del TFG. \\
\end{longtable}
\todo{Confirmar con el tutor (P-06) la forma de presentar la parte
reutilizada y, si procede, cuantificarla (ficheros o líneas reutilizados,
adaptados y nuevos).}

% ---------------------------------------------------------------------
\section{Construcción por incrementos}
\label{sec:imp-fases}

El sistema se ha construido en incrementos, cada uno funcional y verificado
antes de pasar al siguiente, con su rama de trabajo y su
\textit{pull request}. La tabla~\ref{tab:imp-fases} resume el contenido de
cada incremento y el estado de las pruebas a su cierre; la planificación
temporal se presenta en el capítulo de Planificación.

\begin{longtable}{@{}l>{\raggedright\arraybackslash}p{9.6cm}>{\raggedright\arraybackslash}p{3.0cm}@{}}
  \caption{Incrementos de construcción y estado de las pruebas al cerrarlos.}\label{tab:imp-fases}\\
  \toprule
  \textbf{Fase} & \textbf{Contenido} & \textbf{Pruebas al cierre} \\
  \midrule
  \endfirsthead
  \toprule
  \textbf{Fase} & \textbf{Contenido} & \textbf{Pruebas al cierre} \\
  \midrule
  \endhead
  \bottomrule
  \endfoot
  F0 & \textit{Harness} de desarrollo: directrices, plan, requisitos, guías, \textit{skills}, agentes, \textit{hooks}, controles del TFG e integración continua mínima & --- \\
  F1 & Base SaaS importada: renombrado del ámbito de paquetes a \texttt{@pymekit/*}, retirada de módulos no necesarios, base de datos local, puerto propio, CI completa & 61 unit.; pgTAP 27/396; E2E 43 \\
  F2.1 & Base de datos del CMS: esquema \texttt{cms}, pegamento super-administrador $\rightarrow$ Root, MFA obligatorio; primera revisión adversarial de RLS y corrección de dos brechas heredadas & pgTAP 57/1.550 \\
  F2.2 & API del CMS: 19 paquetes de servidor, Hono en \texttt{/api/cms}, autenticación con la sesión de la web & 1.130 unit.; 11 E2E de la API \\
  F2.3 & Base de la interfaz: consola compartida, guardas con \texttt{aal2}, cliente RPC, \textit{fetch} isomorfo & E2E 60 \\
  F2.4 & Explorador de datos: listado con filtros, orden, búsqueda, columnas y vistas; ficha con relaciones; creación, edición y borrado & 1.313 unit.; E2E 86 \\
  F2.5 & Exploradores de usuarios y de almacenamiento & E2E 102 \\
  F2.6 & Auditoría y búsqueda global; auditoría no falsificable y redactada según permisos & pgTAP 60/1.652; E2E 108 \\
  F2.6b & Blog en la base de datos gestionado desde el CMS y rutas públicas; formatos legibles & pgTAP 62/1.730; E2E 129 \\
  F2.7 & Ajustes: general, autenticación, miembros, RBAC (roles, grupos y permisos), recursos y diseñador de fichas & pgTAP 67/2.049; E2E 141 \\
  F2.8 & Paneles con \textit{widgets} de métrica, gráfica y tabla & pgTAP 2.060 \\
  F2.9 & Cierre de la integración del CMS: suite E2E completa, \textit{skills} actualizadas & E2E 149 \\
  F3a--b & Desmarcado completo, español por defecto (unas 1.770 claves), correos y plantillas en español, identidad visual, etiquetas legibles de enumerados, sesión caducada & E2E 150 \\
  F3c & Consola por áreas de negocio, nombres en español de las tablas, retirada del esquema de demostración (ADR-020); solo español (ADR-021); ayuda «?» y resúmenes; bloque académico en la página de inicio & 1.691 unit.; pgTAP 70/2.073; E2E 156 \\
\end{longtable}

Las cifras de pgTAP se expresan como ficheros/pruebas. Las fases F4
(comentarios didácticos), F5 (adaptación a pymes), F6 (validación) y F7
(despliegue) quedan fuera de este capítulo en la fecha de redacción.
\todo{Actualizar la tabla~\ref{tab:imp-fases} al cerrar F4--F7.}

% ---------------------------------------------------------------------
\section{Extractos de código significativos}
\label{sec:imp-extractos}

Se presentan cinco extractos breves que ilustran las decisiones del capítulo
de Diseño. Se han recortado los comentarios y las líneas no esenciales.

\subsection{Una política RLS condicionada al permiso}

El extracto~\ref{lst:imp-rls} muestra la política que permite borrar una
invitación. No basta con pertenecer a la cuenta: hace falta además el permiso
\texttt{invites.manage} en ella. La llamada a \texttt{auth.uid()} va dentro de
\texttt{(select \ldots)} para que PostgreSQL la evalúe una sola vez por
consulta y no por cada fila.

\begin{lstlisting}[style=pymekit, language=SQL, caption={Política de borrado de invitaciones (\texttt{07-invitations.sql}).}, label={lst:imp-rls}]
create policy invitations_delete on public.invitations
  for delete to authenticated
  using (
    has_role_on_account (account_id)
    and public.has_permission (
      (select auth.uid ()),
      account_id,
      'invites.manage'::public.app_permissions
    )
  );
\end{lstlisting}

El blog (extracto~\ref{lst:imp-blog}) combina tres mecanismos: se revocan
todos los privilegios, se concede la lectura \emph{por columnas} (la web no
necesita saber qué usuario creó una entrada) y la política solo devuelve lo
publicado. La escritura no se concede a ningún rol de la API: se hace desde
el CMS, con sus funciones que comprueban el permiso.

\begin{lstlisting}[style=pymekit, language=SQL, caption={Lectura pública del blog (\texttt{19-blog.sql}, recortado).}, label={lst:imp-blog}]
alter table public.blog_posts enable row level security;
revoke all on public.blog_posts from anon, authenticated, service_role;
grant select, insert, update, delete on table public.blog_posts to service_role;

grant select (id, slug, title, excerpt, content, status,
              published_at, category_id /* ... */)
  on table public.blog_posts to anon, authenticated;

create policy blog_posts_read_published on public.blog_posts
  for select to anon, authenticated
  using (status = 'published' and published_at <= now());
\end{lstlisting}

\subsection{Una \textit{server function} protegida por \textit{middleware}}

El extracto~\ref{lst:imp-sfn} es la función que crea una sesión de pago para
un equipo; es la que materializa la segunda barrera del caso CU-04. La cadena
de \textit{middleware} gestiona los errores, exige sesión y comprueba el
permiso \texttt{billing.manage} sobre la cuenta indicada (consultando
\texttt{has\_permission} en la base de datos) antes de ejecutar el
manejador. Los datos de entrada se validan con un esquema Zod.

\begin{lstlisting}[style=pymekit, caption={Server function de pago de equipo (\texttt{billing.functions.ts}, recortado).}, label={lst:imp-sfn}]
const billingManageMiddleware = [
  errorMiddleware,
  authMiddleware,
  withFeaturePermission('billing.manage'),
];

export const createTeamAccountCheckoutSession = createServerFn({
  method: 'POST',
})
  .middleware(billingManageMiddleware)
  .validator(TeamCheckoutSchema)
  .handler(async ({ data }) => {
    if (!featureFlagsConfig.enableTeamAccountBilling) {
      throw new Error('Team account billing is not enabled');
    }

    const client = getSupabaseServerClient(); // cliente con RLS
    const service = createTeamBillingService(client);

    return service.createCheckout(data);
  });
\end{lstlisting}

\subsection{Montaje de la API del CMS y una ruta Hono}

La integración del CMS en la web (ADR-011) se reduce a una ruta de servidor
que reenvía cualquier petición bajo \texttt{/api/cms/} a la aplicación Hono
(extracto~\ref{lst:imp-montaje}). La importación dinámica dentro del
manejador garantiza que el código del CMS ---con el controlador de PostgreSQL
y la clave secreta--- nunca forme parte de un fragmento del navegador.

\begin{lstlisting}[style=pymekit, caption={Montaje de la API del CMS (\texttt{routes/api/cms/\$.ts}).}, label={lst:imp-montaje}]
async function handleCmsApiRequest({ request }: { request: Request }) {
  const { getCmsApiApp } = await import('#/lib/cms/cms-api-app.server.ts');
  const app = await getCmsApiApp();

  return app.fetch(request);
}

export const Route = createFileRoute('/api/cms/$')({
  server: {
    handlers: {
      GET: handleCmsApiRequest,
      POST: handleCmsApiRequest,
      // ... PUT, PATCH, DELETE y OPTIONS igual
    },
  },
});
\end{lstlisting}

Cada funcionalidad del CMS registra sus rutas en la aplicación Hono
(extracto~\ref{lst:imp-hono}): valida la entrada, delega la autorización y la
consulta en su servicio y traduce cualquier error a un código estable, sin
exponer el mensaje interno.

\begin{lstlisting}[style=pymekit, caption={Ruta del registro de auditoría (\texttt{get-audit-logs-route.ts}, recortado).}, label={lst:imp-hono}]
export function registerAuditLogsRoute(router: Hono) {
  return router.get(
    '/v1/audit-logs',
    zValidator('query', AuditLogsQuerySchema, invalidAuditLogsInput),
    async (c) => {
      const { cursor, limit, ...filters } = c.req.valid('query');

      try {
        const data = await createAuditLogsService(c).getAuditLogs({
          cursor, limit, filters,
        });
        return c.json(data);
      } catch (error) {
        // 4xx/5xx con código estable (AUDIT_LOG_*); el detalle va al log
        return respondWithAuditLogsError(c, error, { route: 'list' });
      }
    },
  );
}
\end{lstlisting}

El servicio ejecuta cada consulta dentro de una transacción en la que antes
fija la identidad del usuario (extracto~\ref{lst:imp-claims}). Los valores se
pasan como \emph{parámetros enlazados}, nunca concatenados en el texto SQL, y
el rol se limita a los que puede llevar un JWT de usuario. Así, las políticas
RLS del esquema \texttt{cms} se aplican a la API igual que a cualquier otro
cliente.

\begin{lstlisting}[style=pymekit, caption={Identidad del usuario en cada transacción del CMS (\texttt{drizzle-client.ts}, recortado).}, label={lst:imp-claims}]
const role = ALLOWED_DB_ROLES.has(token.role) ? token.role : 'anon';
const claims = JSON.stringify(token);

return rlsClient.transaction(async (tx) => {
  await tx.execute(sql`select set_config('request.jwt.claims', ${claims}, true)`);
  await tx.execute(sql`select set_config('request.jwt.claim.sub', ${sub}, true)`);
  await tx.execute(sql`select set_config('role', ${role}, true)`);

  return await transaction(tx);
});
\end{lstlisting}

\subsection{La migración de las áreas de la consola}

La consola por áreas (ADR-020) no necesitó cambios de estructura: el área de
cada tabla se guarda en el campo \texttt{ui\_config} de sus metadatos, que ya
existía. La migración (extracto~\ref{lst:imp-areas}) carga los valores por
defecto en una tabla temporal y los aplica \emph{solo} donde no hay un valor
propio, de modo que nunca pisa lo que una pyme haya personalizado desde
Ajustes $\rightarrow$ Recursos y puede ejecutarse de nuevo sin efectos.

\begin{lstlisting}[style=pymekit, language=SQL, caption={Áreas de navegación del CMS (\texttt{20261002140000\_cms\_navigation\_groups.sql}, recortado).}, label={lst:imp-areas}]
insert into cms_navigation_defaults
    (schema_name, table_name, navigation_group, display_name, previous_names, ordering)
values
    ('public', 'blog_posts', 'Blog', 'Entradas', array['Blog: entradas'], 100),
    ('public', 'accounts', 'Cuentas', 'Cuentas', array['Cuentas'], 200),
    ('public', 'subscriptions', 'Facturación', 'Suscripciones',
     array['Facturación: suscripciones'], 310),
    ('public', 'config', 'Sistema', 'Configuración',
     array['Configuración de la plataforma'], 400);
    -- ... resto de tablas

-- Área: solo si la tabla todavía no tiene una.
update cms.table_metadata as meta
set ui_config = coalesce(meta.ui_config, '{}'::jsonb)
        || jsonb_build_object('navigation_group', defaults.navigation_group)
from cms_navigation_defaults as defaults
where meta.schema_name = defaults.schema_name
  and meta.table_name = defaults.table_name
  and nullif(btrim(meta.ui_config ->> 'navigation_group'), '') is null;
\end{lstlisting}

% ---------------------------------------------------------------------
\section{Desarrollo asistido por agentes de IA}
\label{sec:imp-agentes}

El proyecto se ha desarrollado con la asistencia de un agente de
programación de IA (Claude Code~\cite{claudecode}). El autor fija el alcance,
toma las decisiones y revisa los resultados; el agente implementa, prueba y
documenta bajo unas reglas explícitas. Para que esa asistencia sea
controlable y reproducible se construyó, antes de escribir código, un
\textbf{\textit{harness}}: el conjunto de directrices, procedimientos y
controles automáticos que gobiernan el trabajo del agente
(tabla~\ref{tab:imp-harness}).

\begin{table}[htbp]
  \centering
  \footnotesize
  \caption{Componentes del \textit{harness} de desarrollo asistido.}
  \label{tab:imp-harness}
  \begin{tabularx}{\textwidth}{@{}>{\raggedright\arraybackslash}p{3.0cm}X@{}}
    \toprule
    \textbf{Componente} & \textbf{Contenido} \\
    \midrule
    Directrices & \texttt{AGENTS.md} raíz (fuente única) y uno por aplicación o paquete: pila, patrones, reglas de seguridad de la base de datos, convenciones e idioma, y una lista de verificación obligatoria al terminar cada cambio \\
    Registros & Plan, progreso, requisitos, decisiones (ADR), bitácora de incidencias, trazabilidad y mapa de reutilización en \texttt{docs/tfg/}, que el agente actualiza en cada cambio \\
    \textit{Skills} & Procedimientos reutilizables: incorporar un módulo de una base de partida, comentar un módulo en español, registrar una decisión, redactar una sección de la memoria, revisión adversarial del \textit{diff}, búsqueda de fallos, revisión de RLS, experto en PostgreSQL y patrones de servicios, formularios y E2E \\
    Subagentes & Revisor de calidad, revisor de comentarios y redactor de la memoria, con instrucciones y herramientas propias \\
    \textit{Hooks} & Al iniciar la sesión, muestran la fase actual y sus pendientes; tras cada edición, ejecutan el control de marca sobre el fichero; antes de terminar, exigen actualizar el progreso si hay cambios de código \\
    Permisos & Lectura permitida y escritura denegada sobre las bases de partida; lectura denegada sobre ficheros de secretos \\
    Controles en CI & Control de marca bloqueante e informe de comentarios en inglés; un script aparte regenera la tabla de etiquetas de trazabilidad \\
    \bottomrule
  \end{tabularx}
\end{table}

El procedimiento de revisión de RLS merece mención aparte porque es el que
más hallazgos ha producido. Su premisa es que leer las políticas no basta:
hay que \emph{demostrar} el aislamiento. Tras un inventario y una auditoría
estática con un catálogo de ataques, cada hallazgo plausible se convierte en
una prueba pgTAP que ejecuta el ataque como el usuario atacante y comprueba
que falla; después, un agente distinto, que solo recibe la política, la
prueba y la afirmación, intenta \emph{refutarla}. Una brecha solo se acepta
si el refutador la reproduce, y un aprobado, solo si no encuentra un vector
abierto. La refutación encontró, por ejemplo, la lectura del catálogo de
roles de PostgreSQL que la primera auditoría había dado por buena (B-08). Sus
resultados se presentan en el capítulo de Pruebas.

La experiencia también mostró los límites del enfoque. La verificación
exhaustiva con varios agentes por paso (implementador, revisor de seguridad y
refutador, y la suite E2E completa en cada paso) encontró brechas reales
(B-45, B-47), pero su coste en \textit{tokens} resultó excesivo, y el autor
pidió reducirlo (B-48). Desde entonces se trabaja en un «modo ahorro»: un
agente por paso, E2E completos solo al cerrar cada fase y revisión de RLS
solo en cambios importantes de permisos. El riesgo se traslada al cierre de
fase, como mostraron dos pruebas desactualizadas que solo detectó la suite
completa (B-51, B-54). Además, ninguna revisión automática es exhaustiva: una
inyección SQL heredada pasó por varias revisiones de agentes sin ser vista y
la encontró la revisión crítica de sus informes (B-24), y algunos fallos solo
los descubrió el autor con pruebas manuales (B-14, B-40, B-52).
\todo{Valorar en las Conclusiones el efecto de la asistencia de IA con los
datos de dedicación (horas del autor frente a tiempo de reloj).}

% ---------------------------------------------------------------------
\section{Problemas relevantes resueltos}
\label{sec:imp-problemas}

Todas las incidencias se registraron en una bitácora en el momento de
producirse, con el formato «qué pasó, causa, solución, evidencia y lección».
Hasta la fase F3c se registraron 55. La tabla~\ref{tab:imp-bitacora} recoge
una selección de las no relacionadas directamente con el control de acceso;
las de seguridad se analizan en el capítulo de Pruebas
(sección~\ref{sec:pru-seguridad}).

\begin{longtable}{@{}l>{\raggedright\arraybackslash}p{4.8cm}>{\raggedright\arraybackslash}p{7.4cm}@{}}
  \caption{Selección de problemas resueltos y lecciones aprendidas (bitácora del proyecto).}\label{tab:imp-bitacora}\\
  \toprule
  \textbf{ID} & \textbf{Problema} & \textbf{Solución y lección} \\
  \midrule
  \endfirsthead
  \toprule
  \textbf{ID} & \textbf{Problema} & \textbf{Solución y lección} \\
  \midrule
  \endhead
  \bottomrule
  \endfoot
  B-01 & El renombrado masivo del ámbito de paquetes no alcanzó una expresión regular donde aparecía escapado & Se corrigió y el control de marca detecta ahora también las variantes escapadas. Un renombrado masivo se verifica buscando variantes, no solo la forma literal. \\
  B-04 & La base de datos local quedó accesible desde Internet porque los contenedores se saltan el cortafuegos & Regla de filtrado específica para contenedores y acceso por túnel SSH (ADR-010). En un servidor con IP pública, «local» no significa privado. \\
  B-10 & Esquema declarativo y migraciones del CMS de partida desincronizados (nueve funciones distintas) & Se alinearon con las migraciones. La comparación de esquemas forma parte de la verificación de cualquier cambio de base de datos. \\
  B-11 & Una prueba heredada se titulaba «MFA obligatorio por defecto», pero su aserción comprobaba que se concedía el acceso & Se reescribió. Una prueba en verde solo demuestra que el sistema hace lo que dice la prueba: hay que leer las aserciones. \\
  B-16 & Ciclo de dependencias al poner componentes y API del explorador en el mismo paquete & Paquete cliente separado. Separar cliente y servidor también evita ciclos. \\
  B-17, B-29 & La API respondía 500 a registros inexistentes porque el envoltorio de errores perdía el código de PostgreSQL & Se conserva la causa y se traduce el código SQLSTATE a 403/404. Sin la causa no se distingue «no existe» de «ha fallado». \\
  B-35 & La herramienta de comparación de esquemas omitía la opción \texttt{security\_invoker} de una vista y los permisos por columna & Esa parte de la migración se escribe a mano y pgTAP comprueba ambas propiedades. Que la comparación salga limpia no demuestra que la migración sea segura. \\
  B-36 & La sincronización de tablas del CMS duplicaba las relaciones en cada ejecución & Función idempotente, cubierta por pgTAP. \\
  B-40, B-46 & El control de marca solo miraba el contenido de los ficheros, no sus rutas, y no podía ver la marca dibujada en imágenes & Se comprueban también las rutas; logotipo y capturas propias generadas con un script. Un control automático solo cubre lo que se le dice que mire. \\
  B-52 & Tras cambiar de rama, el servidor de desarrollo conservaba un árbol de rutas sin el CMS & Reinicio limpiando cachés. Que las pruebas pasen contra la \textit{build} no garantiza que el entorno de desarrollo esté sano. \\
  B-53 & Una sesión revocada pero con firma válida rompía la consola & Comprobación con estado de la sesión en la entrada de las rutas protegidas. Verificar la firma de un JWT no demuestra que la sesión exista. \\
  B-55 & La lista de miembros superaba el tiempo máximo con unos cientos de cuentas (función volátil evaluada por fila) & Envolver la llamada en \texttt{(select \ldots)}; véase la sección~\ref{sec:pru-rendimiento}. \\
\end{longtable}


\newpage{\ }
\cleardoublepage

\chapter{Pruebas del Sistema}
% =====================================================================
% Capítulo: Pruebas del Sistema
% Fuentes: apps/web/supabase/tests/database (pgTAP), *.test.ts(x) (Vitest),
% apps/e2e/tests (Playwright), .github/workflows/ci.yml,
% docs/tfg/{PROGRESO,BITACORA,TRAZABILIDAD,DECISIONES}.md,
% .claude/skills/rls-review. Cifras: recuento de ficheros con grep y última
% verificación registrada (F3c); unitarias confirmadas con `pnpm test:unit`.
% =====================================================================

% Estilo de los extractos de este capítulo (mismo criterio que en
% Construcción: listings no admite UTF-8 multibyte de forma nativa).
\lstdefinestyle{pymekitpru}{
  basicstyle=\ttfamily\scriptsize,
  breaklines=true,
  frame=single,
  columns=fullflexible,
  keepspaces=true,
  showstringspaces=false,
  literate={á}{{\'a}}1 {é}{{\'e}}1 {í}{{\'i}}1 {ó}{{\'o}}1 {ú}{{\'u}}1
           {ñ}{{\~n}}1 {Á}{{\'A}}1 {É}{{\'E}}1 {Ó}{{\'O}}1 {Ú}{{\'U}}1
}

Este capítulo presenta la estrategia de pruebas de PymeKit y sus
resultados: pruebas unitarias, de integración con la base de datos y la API
del CMS, pruebas de sistema funcionales (E2E por caso de uso) y no
funcionales (seguridad y rendimiento), y pruebas de aceptación. Termina con
las cifras globales y la matriz de trazabilidad entre requisitos y pruebas.

% ---------------------------------------------------------------------
\section{Estrategia de pruebas}
\label{sec:pru-estrategia}

La estrategia sigue el criterio de la pirámide de pruebas: muchas pruebas
rápidas y aisladas en la base, y pocas pruebas de extremo a extremo en la
cima. Su rasgo distintivo es que la base de datos tiene su propio nivel de
pruebas, porque en PymeKit es la frontera de autorización (sección~\ref{sec:dis-seguridad}):
una política RLS mal escrita no se detecta con pruebas de la interfaz, ya que
la interfaz no ofrece el camino para explotarla.

\begin{table}[htbp]
  \centering
  \footnotesize
  \caption{Niveles de prueba, herramientas y qué verifica cada uno.}
  \label{tab:pru-niveles}
  \begin{tabularx}{\textwidth}{@{}>{\raggedright\arraybackslash}p{2.6cm}>{\raggedright\arraybackslash}p{2.4cm}X@{}}
    \toprule
    \textbf{Nivel} & \textbf{Herramienta} & \textbf{Qué verifica} \\
    \midrule
    Unitarias & Vitest & Funciones puras y componentes: filtros, constructor de consultas, formateadores, cálculo de la cuadrícula de paneles, clasificación de errores, componentes de interfaz, coherencia de las claves de traducción \\
    Integración (base de datos) & pgTAP~\cite{pgtap} & Políticas RLS, privilegios, funciones y \textit{triggers}, ejecutados como usuarios reales; aislamiento entre cuentas y entre niveles de personal del CMS \\
    Integración (API del CMS) & Playwright (peticiones HTTP) & Autenticación, permisos y códigos de error de la API montada en la web \\
    Sistema (funcionales) & Playwright~\cite{playwright} & Casos de uso completos en el navegador contra la \textit{build} de pruebas \\
    Sistema (no funcionales) & pgTAP, revisión adversarial, revisión manual & Seguridad (intentos de ataque), rendimiento y usabilidad \\
    Aceptación & Revisión del autor y del tutor & Adecuación a la propuesta \\
    \bottomrule
  \end{tabularx}
\end{table}

Todas las pruebas automáticas forman parte de la integración continua
(sección~\ref{sec:imp-entorno}): las unitarias se ejecutan en cada
\textit{push} a la rama principal y en cada \textit{pull request}; pgTAP y E2E, solo en cada \textit{pull request} hacia la rama
principal. Además, cada cambio se cierra con una lista de verificación
obligatoria (tipos, \textit{lint}, formato, control de marca, pruebas
afectadas, revisión adversarial del \textit{diff} y, si toca la base de
datos, revisión de RLS). Las pruebas E2E se ejecutan siempre contra la
\textit{build} de pruebas y no contra el servidor de desarrollo (B-13).

% ---------------------------------------------------------------------
\section{Pruebas unitarias}
\label{sec:pru-unitarias}

Hay 116 ficheros de pruebas unitarias repartidos en 28 paquetes, con 1.691
pruebas superadas y una omitida. La mayor parte corresponde a los paquetes
del CMS, que concentran la lógica algorítmica: paneles (262 pruebas), núcleo
del explorador de datos (234), constructor de consultas (222), formateadores
(218), núcleo de filtros (181) y componentes del explorador (137).

Algunas pruebas unitarias son, en realidad, pruebas de seguridad. Por
ejemplo, la de \texttt{buildParameterizedStatement} reproduce el valor
malicioso que permitía la inyección SQL heredada de la incidencia B-24
(un nombre que contenía la cadena \texttt{\$2}) y comprueba que ahora viaja
como parámetro enlazado; y las de protección de usuarios verifican que nadie
puede actuar sobre sí mismo ni sobre un super-administrador (B-21). Otras
garantizan propiedades transversales, como que todas las claves de traducción
usadas existen en los mensajes en español.

% ---------------------------------------------------------------------
\section{Pruebas de integración}
\label{sec:pru-integracion}

\subsection{Base de datos: RLS y aislamiento multi-tenant}
\label{sec:pru-pgtap}

La batería pgTAP tiene 70 ficheros (68 de pruebas y 2 de utilidades) con
2.073 pruebas. Cada fichero se ejecuta dentro de una transacción que se
deshace al terminar; crea usuarios de prueba, se \emph{autentica} como cada
uno de ellos (fijando su JWT y el rol \texttt{authenticated} en la sesión,
como hace el servicio de autenticación) y comprueba qué puede ver y hacer.
Se organizan en cuatro grupos:

\begin{itemize}
  \item \textbf{Aislamiento de la plataforma} (cuentas personales y de
        equipo, membresías, invitaciones y sus límites, transferencia de
        propiedad, facturación, notificaciones, almacenamiento,
        super-administrador): un usuario solo ve y modifica lo de sus
        cuentas, y solo con el permiso adecuado.
  \item \textbf{Superficie de privilegios}: \texttt{privileges.test.sql}
        comprueba que ningún rol de la API conserva privilegios residuales
        (como \texttt{TRUNCATE}, que ignora RLS) y \texttt{anon-surface.test.sql}
        fija la superficie exacta del rol anónimo ---solo puede leer el blog
        y no ejecuta ninguna función--- y falla si se amplía (B-39).
  \item \textbf{Control de acceso del CMS} (41 ficheros): permisos de datos,
        sistema y almacenamiento, rangos, delegación, denegaciones,
        auditoría, paneles, vistas guardadas, búsqueda global y el pegamento
        con el super-administrador.
  \item \textbf{Regresión de seguridad}: un fichero por revisión adversarial
        (\texttt{cms-isolation}, \texttt{cms-audit-isolation},
        \texttt{cms-members-isolation}, \texttt{cms-rbac-isolation}\ldots) que
        conserva como prueba cada ataque estudiado.
\end{itemize}

El extracto~\ref{lst:pru-tenant} muestra el patrón básico de una prueba de
aislamiento: tras crear una suscripción para una cuenta, un usuario ajeno no
obtiene ninguna fila.

\begin{lstlisting}[style=pymekitpru, language=SQL, caption={Aislamiento de la facturación entre cuentas (\texttt{team-billing-subscriptions.test.sql}, recortado).}, label={lst:pru-tenant}]
select tests.create_supabase_user('foreigner');
select pymekit.authenticate_as('foreigner');

select is_empty(
  $$ select 1 from subscriptions where id = 'sub_test' $$,
  'The account cannot read the other account subscriptions'
);
\end{lstlisting}

Las pruebas de regresión de las revisiones adversariales ejecutan el ataque
tal como se descubrió. El extracto~\ref{lst:pru-ataque} reproduce la escalada
de la incidencia B-05: antes de la corrección la sentencia afectaba a una
fila; ahora debe fallar con el código de permiso denegado (\texttt{42501}).
La misma prueba comprueba a continuación que la operación legítima sigue
funcionando, para que la corrección no sea simplemente «bloquearlo todo».

\begin{lstlisting}[style=pymekitpru, language=SQL, caption={Regresión de la escalada por \texttt{UPDATE} sin \texttt{WITH CHECK} (\texttt{cms-isolation.test.sql}).}, label={lst:pru-ataque}]
select cms_tests.authenticate_as('iso_mid');

-- El ataque que funcionaba antes de la corrección (UPDATE 1)
select throws_ok(
    $$ update cms.role_permissions
       set permission_id = cms_tests.test_uuid(9304)
       where role_id = cms_tests.test_uuid(9202)
         and permission_id = cms_tests.test_uuid(9303) $$,
    '42501',
    null,
    'E1: no puede colgar de un rol inferior un permiso que no posee'
);
\end{lstlisting}

\subsection{API del CMS}
\label{sec:pru-api}

La API del CMS se prueba de dos formas. Sus servicios y utilidades tienen
pruebas unitarias con el cliente de datos simulado. Y la API montada en la
web se prueba con peticiones HTTP reales desde Playwright
(\texttt{cms-api.spec.ts} y las comprobaciones de API de los demás ficheros
del CMS): sin sesión responde 401; con sesión pero sin acceso al CMS, o sin
segundo factor, o con la cuenta del CMS desactivada, 403; una petición con
campos no permitidos (por ejemplo, metadatos de la aplicación al crear un
usuario) se rechaza con 400; y cada rechazo se comprueba por su código de
error estable, no solo por el estado HTTP, porque un 403 no demuestra nada si
no se sabe por qué se produjo (B-26).

% ---------------------------------------------------------------------
\section{Pruebas de sistema}
\label{sec:pru-sistema}

\subsection{Pruebas funcionales por caso de uso}
\label{sec:pru-funcionales}

Las pruebas E2E se organizan en 33 ficheros (unas 164 pruebas declaradas),
19 de ellos del CMS, con 99 pruebas. Se ejecutan en español, el único idioma
de la interfaz, y un fichero específico comprueba que no queda ningún texto
en inglés visible. La tabla~\ref{tab:pru-cu} relaciona cada caso de uso con
las pruebas que lo cubren.

\begin{table}[htbp]
  \centering
  \footnotesize
  \caption{Cobertura E2E de los casos de uso.}
  \label{tab:pru-cu}
  \begin{tabularx}{\textwidth}{@{}l>{\raggedright\arraybackslash}p{5.0cm}>{\raggedright\arraybackslash}X@{}}
    \toprule
    \textbf{CU} & \textbf{Ficheros de prueba} & \textbf{Estado} \\
    \midrule
    CU-01 & \texttt{authentication/*} (registro, acceso, recuperación de contraseña, redirecciones, sesión caducada), \texttt{account} & Cubierto \\
    CU-02 & \texttt{team-accounts/*}, \texttt{invitations} (incluida la aceptación con MFA) & Cubierto con los roles actuales; roles orientados a pymes pendientes (F5) \\
    CU-03 & \texttt{team-billing}, \texttt{user-billing} & Escritos, pero desactivados mientras no se configuren las claves de la pasarela en modo de pruebas (F5/F6) \\
    CU-04 & pgTAP de facturación y permisos (capa de datos) & Parcial: falta la prueba E2E de la interfaz y de la \textit{server function} (F6) \\
    CU-05 & \texttt{admin} (404 sin permiso, MFA obligatorio, bloquear, reactivar, borrar) & Parcial: la prueba de suplantación está desactivada \\
    CU-06 & \texttt{cms-data-explorer*}, \texttt{cms-audit-logs}, \texttt{cms-settings-*}, \texttt{cms-ui}, \texttt{cms-sidebar-areas}, \texttt{cms-field-help}, \texttt{cms-page-summaries} & Cubierto \\
    CU-07 & --- & Pendiente: evaluación de la reutilización (F6, P-01) \\
    CU-08 & \texttt{blog} (lectura anónima, borradores invisibles, Markdown seguro, publicación desde el CMS, \textit{sitemap}) & Cubierto \\
    \bottomrule
  \end{tabularx}
\end{table}

Las pruebas del CMS verifican sistemáticamente los dos lados de cada
permiso: el personal de soporte, con un rol limitado, \emph{no} ve en la
barra lateral ni puede abrir por URL las tablas sin permiso, y el
super-administrador sí. Muchas incluyen además pruebas negativas
específicas derivadas de incidencias: la guarda de clave en la edición
(B-19), un HTML disfrazado de imagen en la subida de ficheros (B-25), el
\textit{widget} de una tabla no legible («sin acceso», B-50) o la consola con
una sesión sin segundo factor (B-14).

\subsection{Pruebas no funcionales: seguridad}
\label{sec:pru-seguridad}

La seguridad se ha verificado con revisiones adversariales en dos niveles:
una revisión del \textit{diff} orientada a la búsqueda de fallos al terminar
cada funcionalidad, y una revisión específica de RLS (descrita en la
sección~\ref{sec:imp-agentes}) en cada cambio relevante de políticas,
privilegios o funciones de privilegios elevados. La tabla~\ref{tab:pru-rls}
resume las revisiones de RLS realizadas.

\begin{table}[htbp]
  \centering
  \footnotesize
  \caption{Revisiones adversariales de RLS y su resultado.}
  \label{tab:pru-rls}
  \begin{tabularx}{\textwidth}{@{}l>{\raggedright\arraybackslash}p{3.4cm}X@{}}
    \toprule
    \textbf{Fase} & \textbf{Alcance} & \textbf{Resultado} \\
    \midrule
    F2.1 & Esquema \texttt{cms} completo tras su incorporación & Dos brechas heredadas confirmadas (B-05, B-08), cinco endurecimientos y cuatro funciones desalineadas (ADR-015). La refutación independiente encontró B-08 \\
    F2.6 & Auditoría y búsqueda global & Aislado; cuatro debilidades menores endurecidas (B-31 a B-34); +31 pruebas \\
    F2.6b & Blog y apertura del esquema \texttt{public} al rol anónimo & Aislado; riesgo latente de privilegios por defecto documentado y vigilado por prueba (B-39) \\
    F2.7a & Miembros del CMS y autor de la auditoría & Aislado; +68 pruebas; una inyección SQL latente de origen propio (B-45) \\
    F2.7b & RBAC del CMS con administradores delegados de distintos rangos & Brecha heredada de escape de denegaciones mediante delegación, reproducida por el refutador y corregida (B-47) \\
    \bottomrule
  \end{tabularx}
\end{table}

La tabla~\ref{tab:pru-hallazgos} recoge los hallazgos de seguridad más
relevantes de todo el proyecto. De las 29 incidencias de seguridad
registradas, 25 procedían del código de partida, 2 del entorno y 2 se
introdujeron durante el proyecto. Todas están corregidas y cubiertas por
pruebas de regresión.

\begin{longtable}{@{}l>{\raggedright\arraybackslash}p{6.6cm}lll@{}}
  \caption{Principales hallazgos de seguridad corregidos (bitácora del proyecto).}\label{tab:pru-hallazgos}\\
  \toprule
  \textbf{ID} & \textbf{Fallo} & \textbf{Origen} & \textbf{Grav.} & \textbf{Detección} \\
  \midrule
  \endfirsthead
  \toprule
  \textbf{ID} & \textbf{Fallo} & \textbf{Origen} & \textbf{Grav.} & \textbf{Detección} \\
  \midrule
  \endhead
  \bottomrule
  \endfoot
  B-04 & Base de datos local accesible desde Internet & entorno & Alta & Diagnóstico de red \\
  B-05 & Escalada en el RBAC del CMS: \texttt{UPDATE} sin \texttt{WITH CHECK} & heredado & Alta & Revisión de RLS \\
  B-06 & La comprobación de MFA del CMS fallaba en abierto & heredado & Alta & Prueba nueva \\
  B-07 & Lectura de los usuarios de autenticación (con \textit{hashes}) por una función de consulta & heredado & Alta & Revisión de RLS \\
  B-08 & Lectura del catálogo de roles de PostgreSQL & heredado & Alta & Refutación \\
  B-14 & La consola se cargaba sin segundo factor & propio & Media & Prueba manual \\
  B-19 & Edición y borrado de varias filas con una condición que no era clave & heredado & Alta & Revisión del \textit{diff} \\
  B-21 & La gestión de usuarios solo comprobaba el acceso al CMS, no el permiso ni el objetivo & heredado & Alta & Revisión de seguridad \\
  B-24 & Inyección SQL por sustitución manual de parámetros & heredado & Alta & Revisión de informes \\
  B-25 & Subidas con el tipo de contenido elegido por el cliente & heredado & Media & Revisión de seguridad \\
  B-27 & Entradas de auditoría falsificables por el personal & heredado & Alta & Revisión de seguridad \\
  B-28 & La búsqueda global podía buscar en los usuarios de autenticación & heredado & Alta & Revisión de seguridad \\
  B-42 & Un miembro podía cambiar sus propios roles & heredado & Alta & Revisión de seguridad \\
  B-44 & Cualquiera con permiso de ajustes podía desactivar el MFA obligatorio & heredado & Media & Revisión de seguridad \\
  B-45 & Función interna de borrado ejecutable directamente (inyección latente) & propio & Alta & Revisión de RLS \\
  B-47 & Escape de denegaciones explícitas delegando en una cuenta títere & heredado & Alta & Revisión de RLS \\
  B-50 & \textit{Widgets} sobre tablas no legibles y apropiación de paneles compartidos & heredado & Media & Revisión del \textit{diff} \\
\end{longtable}

De estos resultados se extraen tres conclusiones sobre el proceso. Primera,
la demostración con pruebas ejecutables y la refutación independiente
encuentran fallos que la lectura estática da por buenos (B-08, B-47).
Segunda, las revisiones se complementan: B-45, introducido en F2.1, no lo vio
la revisión de esa fase y lo encontró la de F2.7a, que buscaba otra cosa.
Tercera, la defensa en profundidad funciona: en B-14 falló la guarda de la
interfaz, pero no se expuso ningún dato porque la API y la base de datos
seguían exigiendo el segundo factor.

\subsection{Pruebas no funcionales: rendimiento}
\label{sec:pru-rendimiento}

No se han realizado pruebas de carga formales; el requisito RNF-09 pide un
rendimiento \emph{razonable} y se verifica por diseño (carga en el servidor,
caché de consultas, paginación en el servidor con recuento aproximado en
tablas grandes y tiempo máximo de 5~s en la búsqueda global, B-32) y con
pruebas puntuales. La incidencia más ilustrativa (B-55) surgió de forma
espontánea: la base de datos de pruebas, que acumula cuentas en cada
ejecución completa de los E2E, llegó a unas 200 cuentas y la lista de
miembros de un equipo empezó a superar el tiempo máximo de sentencia (8~s),
con consultas de más de 6~s. La causa era una función volátil
(\texttt{active\_account\_id()}) en la cláusula \texttt{WHERE}, que PostgreSQL
evaluaba una vez por cada fila de \texttt{accounts}. Envolverla en
\texttt{(select \ldots)} hace que se evalúe una sola vez como
\textit{InitPlan} y la consulta volvió a ser inmediata. El mismo patrón se
aplica de forma sistemática a \texttt{auth.uid()} en las políticas RLS
(extracto~\ref{lst:imp-rls}).
\todo{Si se quiere reforzar RNF-09, medir tiempos de respuesta de las
pantallas principales con un volumen de datos representativo (por ejemplo,
10.000 cuentas) y añadir la tabla de resultados.}

\subsection{Pruebas no funcionales: usabilidad y accesibilidad}
\label{sec:pru-usabilidad}

La usabilidad y la accesibilidad (RNF-08) se han verificado mediante revisión
manual del autor, que dio lugar a varios rediseños (consola por áreas,
nombres en español, ayuda contextual y resúmenes, ADR-020), y mediante
pruebas E2E que, en buena parte, localizan los elementos por su rol y su
nombre accesible, de modo que fallan si un control pierde su etiqueta; la
prueba de solo español revisa también los atributos accesibles.
\todo{Añadir, si se realiza, una evaluación de accesibilidad automatizada
(por ejemplo, con una herramienta de auditoría de páginas) o con usuarios.}

% ---------------------------------------------------------------------
\section{Pruebas de aceptación}
\label{sec:pru-aceptacion}

Las pruebas de aceptación validan que el sistema satisface la propuesta del
TFG desde el punto de vista de su destinatario. Durante el desarrollo, el
autor ha actuado como cliente: probó a mano cada incremento (por ejemplo,
CU-01 al cierre de F1) y de esas pruebas surgieron incidencias (B-14, B-40,
B-52) y cambios de diseño (ADR-017, ADR-020, ADR-021).
\todo{Completar tras F6: sesión de aceptación con el tutor recorriendo CU-01
a CU-08 con un guion (precondición, pasos, resultado esperado y obtenido),
resultado de la evaluación de la reutilización (CU-07, P-01) y conformidad
final.}

% ---------------------------------------------------------------------
\section{Resultados}
\label{sec:pru-resultados}

La tabla~\ref{tab:pru-cifras} resume el estado de las pruebas en la última
verificación completa (cierre de F3c).

\begin{table}[htbp]
  \centering
  \footnotesize
  \caption{Resultados de las pruebas automáticas (cierre de F3c).}
  \label{tab:pru-cifras}
  \begin{tabularx}{\textwidth}{@{}>{\raggedright\arraybackslash}p{3.4cm}rr>{\raggedright\arraybackslash}X@{}}
    \toprule
    \textbf{Tipo} & \textbf{Ficheros} & \textbf{Pruebas} & \textbf{Resultado} \\
    \midrule
    Unitarias (Vitest) & 116 & 1.692 & 1.691 superadas, 1 omitida \\
    Base de datos (pgTAP) & 70 & 2.073 & Todas superadas \\
    E2E (Playwright) & 33 & $\approx$164 & 156 superadas; el resto, condicionadas a configuración (facturación, gestión de equipos y suplantación) \\
    Control de marca & --- & --- & 0 referencias \\
    \bottomrule
  \end{tabularx}
\end{table}

La evolución de las cifras a lo largo de los incrementos se recoge en la
tabla~\ref{tab:imp-fases}: la batería pgTAP pasó de 396 pruebas heredadas a
2.073, y la E2E de 43 a 156, en buena parte por las pruebas de regresión de
seguridad. Quedan algunas pruebas E2E inestables bajo ejecución paralela, que
pasan al reintentarlas o en serie y cuya estabilización está prevista antes
de la validación final.

La tabla~\ref{tab:pru-trazabilidad} cierra el capítulo con la trazabilidad
entre cada requisito y su evidencia de verificación.

\begin{longtable}{@{}l>{\raggedright\arraybackslash}p{8.4cm}>{\raggedright\arraybackslash}p{3.6cm}@{}}
  \caption{Trazabilidad entre requisitos y pruebas (cierre de F3c).}\label{tab:pru-trazabilidad}\\
  \toprule
  \textbf{Req.} & \textbf{Evidencia} & \textbf{Estado} \\
  \midrule
  \endfirsthead
  \toprule
  \textbf{Req.} & \textbf{Evidencia} & \textbf{Estado} \\
  \midrule
  \endhead
  \bottomrule
  \endfoot
  RF-01 & E2E \texttt{authentication}, \texttt{account}, \texttt{blog}; pgTAP \texttt{blog}, \texttt{anon-surface} & Verificado \\
  RF-02 & Unitarias; E2E de todos los flujos & Verificado \\
  RF-03 & E2E \texttt{authentication} & Verificado \\
  RF-04 & pgTAP \texttt{super-admin}, \texttt{cms-verify-admin-access}; E2E \texttt{admin}, \texttt{team-invitation-mfa}, \texttt{cms-settings-auth} & MFA verificado; OAuth configurable sin prueba \\
  RF-05 & E2E \texttt{account} & Verificado \\
  RF-06 & pgTAP de cuentas, membresías, invitaciones y permisos; E2E \texttt{team-accounts}, \texttt{invitations} & Parcial (roles para pymes en F5) \\
  RF-07 & pgTAP de facturación; E2E de facturación desactivados & Parcial (F5/F6) \\
  RF-08 & E2E \texttt{admin}, \texttt{cms-ui} & Parcial (suplantación sin E2E) \\
  RF-09 & pgTAP del CMS (41 ficheros); unitarias del CMS; E2E del CMS (19 ficheros) & Verificado \\
  RF-10 & pgTAP de auditoría e integridad; E2E \texttt{cms-audit-logs} & Verificado \\
  RF-11 & pgTAP de paneles; unitarias de paneles; E2E \texttt{cms-dashboards} & Verificado \\
  RF-12 & pgTAP \texttt{notifications} & Parcial \\
  RF-13 & Unitaria de coherencia de claves; E2E \texttt{i18n-solo-espanol} & Verificado \\
  RNF-01 & Configuración tipada; evaluación de la reutilización & Pendiente (F6) \\
  RNF-02 & pgTAP; revisiones adversariales (tabla~\ref{tab:pru-rls}) & Verificado \\
  RNF-03 & pgTAP de aislamiento de la plataforma y del CMS & Verificado \\
  RNF-04 & Estructura del monorepo; comprobación de tipos y dependencias & Verificado \\
  RNF-05 & CI (tipos, \textit{lint}, formato, control de marca) & Parcial (comentarios en F4) \\
  RNF-06 & CI con las tres suites & Parcial (completar en F6) \\
  RNF-07 & Despliegue de prueba y manual & Pendiente (F7) \\
  RNF-08 & Revisión manual; E2E por rol y nombre accesible & Parcial \\
  RNF-09 & Revisión del diseño; B-55 & Parcial \\
\end{longtable}
\todo{Actualizar la tabla~\ref{tab:pru-trazabilidad} y las cifras de la
tabla~\ref{tab:pru-cifras} tras F4--F7.}


\newpage{\ }
\cleardoublepage

\chapter{Despliegue del Sistema}
\label{cap:despliegue}
% =====================================================================
% Capítulo · Despliegue del Sistema (borrador, fase F7)
% Lo incluye \chapter{Despliegue del Sistema}; aquí solo van secciones.
% Fuentes: apps/web/Dockerfile, docker-compose.yml,
%   docker-compose.supabase-local.yml, .env.docker.example, railway.json,
%   docs/DESPLIEGUE.md y docs/tfg/DECISIONES.md (ADR-011, ADR-022).
% Paquetes necesarios: booktabs, tabularx, listings, tikz
%   (bibliotecas: positioning, arrows.meta, fit).
% =====================================================================

Este capítulo describe cómo se empaqueta y se pone en producción PymeKit. El
objetivo de la fase de despliegue era doble: que cualquier persona pudiera
reproducir la plataforma completa en su equipo con un único comando, y que la
misma configuración sirviera para publicar una instancia de demostración
accesible desde Internet. Ambos objetivos se apoyan en una decisión de diseño
previa: el CMS de datos forma parte de la aplicación web y no es un servicio
aparte (ADR-011), de modo que el sistema que hay que desplegar consta de una
sola aplicación más los servicios gestionados de los que depende.

\section{Arquitectura de despliegue}
\label{sec:despliegue-arquitectura}

La Figura~\ref{fig:despliegue} muestra los nodos que intervienen en una
instalación de PymeKit. El único componente propio es el servidor de la
aplicación, un proceso Node.js sin estado que sirve tanto las páginas
renderizadas en el servidor como las rutas de API (\emph{server functions},
\emph{webhooks} y la API del CMS). Todo el estado persistente reside en la
plataforma de base de datos como servicio (Postgres, autenticación y
almacenamiento de ficheros), y los pagos se delegan en la pasarela externa.

\begin{figure}[htbp]
  \centering
  % Se escala al ancho de línea para que encaje en cualquier plantilla.
  \resizebox{\textwidth}{!}{%
  \begin{tikzpicture}[
      node distance=1.4cm and 2.2cm,
      box/.style={draw, rounded corners, align=center, minimum width=3.4cm,
                  minimum height=1.1cm, font=\small},
      group/.style={draw, dashed, rounded corners, inner sep=0.35cm},
      arr/.style={-{Stealth[length=2.2mm]}, thick},
      lbl/.style={font=\scriptsize, midway, fill=white, inner sep=1pt}
    ]
    \node[box] (browser) {Navegador\\(usuario o administrador)};
    \node[box, right=of browser] (web) {Servidor de la aplicación\\Node.js (contenedor)\\\texttt{/}, \texttt{/admin/cms}, \texttt{/api/*}};
    \node[box, above right=0.6cm and 2.2cm of web] (db) {Postgres\\(RLS, esquema \texttt{cms})};
    \node[box, right=2.2cm of web] (auth) {Auth y Storage};
    \node[box, below right=0.6cm and 2.2cm of web] (stripe) {Pasarela de pagos\\(modo test)};
    \node[group, fit=(db)(auth), label={[font=\scriptsize]above:Plataforma BaaS}] {};

    \draw[arr] (browser) -- node[lbl]{HTTPS} (web);
    \draw[arr] (web) -- node[lbl, sloped]{SQL / REST} (db);
    \draw[arr] (web) -- node[lbl]{REST} (auth);
    \draw[arr] (web) -- node[lbl, sloped]{API} (stripe);
    \draw[arr, dashed] (stripe.west) to[bend left=25] node[lbl, sloped]{\emph{webhook}} (web.south);
    \draw[arr, dashed] (browser.north) to[bend left=20] node[lbl, sloped]{sesión (clave pública)} (auth.north west);
  \end{tikzpicture}%
  }
  \caption{Diagrama de despliegue de PymeKit.}
  \label{fig:despliegue}
\end{figure}

La Tabla~\ref{tab:componentes-despliegue} resume la función de cada componente
y el artefacto que lo configura dentro del repositorio.

\begin{table}[htbp]
  \centering
  \small
  \begin{tabularx}{\textwidth}{@{}lXl@{}}
    \toprule
    \textbf{Componente} & \textbf{Función} & \textbf{Configuración} \\
    \midrule
    Servidor de la aplicación & Renderizado, \emph{server functions}, API del CMS y \emph{webhooks} & \texttt{apps/web/Dockerfile} \\
    Orquestación local & Construye la imagen y la arranca con sus variables & \texttt{docker-compose.yml} \\
    Base de datos & Esquemas, políticas RLS y RBAC del CMS & \texttt{apps/web/supabase/migrations} \\
    Autenticación y ficheros & Sesiones, MFA, URLs de redirección, \emph{buckets} & \texttt{config.toml} / panel \\
    Pasarela de pagos & Suscripciones; notifica cambios por \emph{webhook} & variables \texttt{STRIPE\_*} \\
    Plataforma de alojamiento & Build, arranque y sonda de salud en la nube & \texttt{railway.json} \\
    \bottomrule
  \end{tabularx}
  \caption{Componentes del despliegue y su configuración.}
  \label{tab:componentes-despliegue}
\end{table}

\section{Contenerización de la aplicación}
\label{sec:despliegue-imagen}

La aplicación se empaqueta en una imagen de contenedor definida mediante un
\emph{Dockerfile} de tres etapas. Construir por etapas permite que las
herramientas de compilación, que ocupan cientos de megabytes, no lleguen a la
imagen final:

\begin{enumerate}
  \item \textbf{Recorte del monorepo.} La herramienta de orquestación del
    monorepo calcula el subconjunto de paquetes del que depende la aplicación
    web y separa sus manifiestos del resto del código. Así, la capa que
    instala dependencias solo se invalida cuando cambian estas, y no con cada
    modificación del código fuente.
  \item \textbf{Compilación.} Se instalan las dependencias a partir del
    fichero de bloqueo (instalación reproducible) y se ejecuta la build de
    producción. El servidor resultante queda autocontenido: el empaquetador
    incluye en la carpeta de salida las librerías que el servidor necesita,
    de forma que no hace falta copiar el directorio de dependencias.
  \item \textbf{Ejecución.} La imagen final parte de una distribución mínima
    de Node.js y solo contiene la carpeta de salida. El proceso se ejecuta
    con un usuario sin privilegios y la imagen declara una sonda de salud
    sobre la ruta \texttt{/api/healthcheck}, que el motor de contenedores usa
    para marcar el servicio como sano o reiniciarlo.
\end{enumerate}

Un aspecto que condiciona el diseño es la distinción entre variables
\emph{de compilación} y \emph{de ejecución}. Las variables públicas (prefijo
\texttt{VITE\_}: URL del sitio, URL y clave pública de la base de datos, clave
publicable de la pasarela) se incrustan en el código durante la build, por lo
que se pasan como argumentos de construcción. Los secretos de servidor (clave
secreta de la base de datos, cadena de conexión de la API del CMS, claves de
la pasarela y del correo) se inyectan únicamente al arrancar el contenedor y
nunca quedan dentro de la imagen, lo que permite publicarla o reutilizarla sin
exponer credenciales. Para que los argumentos vacíos no anulen los valores por
defecto versionados en el repositorio, la build solo escribe los que llegan
con valor y falla de forma explícita si falta alguno de los obligatorios.

\section{Orquestación con Docker Compose}
\label{sec:despliegue-compose}

El fichero \texttt{docker-compose.yml} de la raíz describe el servicio de la
aplicación. No levanta la base de datos: la aplicación se conecta a una
instancia ya existente, sea la del entorno de desarrollo local o un proyecto
en la nube, y el mismo fichero sirve para ambos casos. El
Listado~\ref{lst:compose} recoge su parte esencial.

\begin{lstlisting}[language=bash, caption={Servicio de la aplicación en \texttt{docker-compose.yml} (extracto).}, label={lst:compose}, basicstyle=\ttfamily\footnotesize]
services:
  web:
    image: pymekit-web:latest
    build:
      context: .
      dockerfile: apps/web/Dockerfile
      args:            # publicas: se incrustan en la build
        VITE_SITE_URL: ${VITE_SITE_URL:?Falta VITE_SITE_URL}
        VITE_SUPABASE_URL: ${VITE_SUPABASE_URL:?Falta VITE_SUPABASE_URL}
        VITE_SUPABASE_PUBLIC_KEY: ${...:?Falta}
    env_file:          # de servidor: se leen al arrancar
      - path: .env.docker
        required: true
    ports:
      - '${WEB_PORT:-3200}:3000'
    restart: unless-stopped
\end{lstlisting}

Todas las variables se documentan en una plantilla versionada,
\texttt{.env.docker.example}, sin valores secretos; el fichero real con las
credenciales está excluido tanto del control de versiones como del contexto
de construcción de la imagen. Para el entorno local se añade un segundo
fichero de ajuste que ejecuta el contenedor con la red del anfitrión. La razón
es que la URL de la base de datos se incrusta en la build y la usan tanto el
navegador como el servidor: con la red del anfitrión, la dirección local
designa el mismo servicio para ambos y no es necesario mantener dos URLs
distintas.

\section{Despliegue en la nube}
\label{sec:despliegue-nube}

La instancia de demostración combina tres servicios gestionados:

\begin{itemize}
  \item \textbf{Base de datos como servicio.} Se crea un proyecto y se le
    aplican las migraciones del repositorio con la interfaz de línea de
    órdenes de la plataforma (\texttt{supabase db push}), que incluyen tanto
    el esquema de la aplicación como el del CMS. Después se configuran las
    URLs autorizadas de autenticación, un servidor de correo propio y el
    \emph{webhook} de base de datos que avisa a la aplicación cuando se borra
    una suscripción. Los datos de ejemplo (\emph{seed}) no se aplican en este
    entorno.
  \item \textbf{Pasarela de pagos en modo test.} Se registra un
    \emph{webhook} público que apunta a \texttt{/api/billing/webhook}; la
    aplicación verifica la firma de cada evento con el secreto compartido
    antes de actualizar las suscripciones en la base de datos.
  \item \textbf{Plataforma de alojamiento.} El fichero \texttt{railway.json}
    declara la orden de build, la de arranque, la ruta de la sonda de salud y
    la política de reinicio ante fallos. Las variables se definen en el panel
    del servicio y cada integración en la rama principal produce un nuevo
    despliegue.
\end{itemize}

El procedimiento completo, junto con la actualización y las copias de
seguridad, se recoge en el manual de instalación (Anexo~\ref{anx:manual-instalacion}).
% PENDIENTE: crear la etiqueta anx:manual-instalacion al pasar
% docs/DESPLIEGUE.md al anexo de la memoria (F8).

\section{Verificación}
\label{sec:despliegue-verificacion}

La imagen se construyó y se arrancó contra la base de datos del entorno local,
comprobando que la página de inicio y la de inicio de sesión responden con
código 200 y que la sonda de salud informa de que la base de datos es
accesible. La imagen final ocupa unos 265~MB en disco
(medido el 20/10/2026), frente a los varios gigabytes que suponen las dependencias de
desarrollo del monorepo.

% PENDIENTE: completar con la URL y la fecha de la instancia de demostración
% cuando se publique (F7, tareas 3–5), y con la prueba del manual desde cero.

 
    % EPILOGO
    \part{Epílogo}
    \null\vfill
\noindent En esta parte se incluyen las conclusiones, la información relativa a la licencia del software y la bibliografía.
\vfill

% Conclusiones
\chapter{Conclusiones}
\label{cap:conclusiones}
% ------------------------------------------------------------------------------
% Capítulo · Conclusiones
% Fuentes: docs/tfg/{PROGRESO,BITACORA,TRAZABILIDAD,DECISIONES}.md.
% Estado de las fuentes: 20/10/2026 (F0–F3 cerradas; F4–F7 pendientes).
% Las filas marcadas como parciales o pendientes deben actualizarse al cerrar
% las fases F4–F7.
% ------------------------------------------------------------------------------

Este capítulo valora el grado de cumplimiento de los objetivos planteados en
la introducción, recoge las lecciones aprendidas durante el proyecto y
enumera las líneas de trabajo futuro.

\section{Objetivos alcanzados}
\label{sec:concl-objetivos}

El objetivo general era diseñar, implementar y validar una arquitectura
software reutilizable para el desarrollo ágil de aplicaciones SaaS orientadas
a pymes, con \textit{frontend}, \textit{backend}, autenticación, gestión de
usuarios y roles, pagos y un CMS para la administración de contenidos y datos.
PymeKit reúne todos esos componentes en un único monorepo y un único servicio
web, sobre una sola base de datos en la que se aplica la autorización. La
tabla~\ref{tab:concl-objetivos} detalla el resultado de cada objetivo
específico (tabla~\ref{tab:intro-objetivos}).

\todo{Actualizar el estado de OE-4, OE-7, OE-8 y OE-9 al cerrar las fases
F4--F7.}

\begin{table}[htbp]
  \centering
  \footnotesize
  \caption{Grado de cumplimiento de los objetivos específicos.}
  \label{tab:concl-objetivos}
  \begin{tabularx}{\textwidth}{@{}lX>{\raggedright\arraybackslash}p{2cm}@{}}
    \toprule
    \textbf{Obj.} & \textbf{Resultado} & \textbf{Estado} \\
    \midrule
    OE-1 & Contexto analizado, componentes comunes identificados y requisitos definidos: 13 funcionales, 9 no funcionales y 7 casos de uso, trazados al código, a las pruebas y a la memoria. & Cumplido \\
    OE-2 & Monorepo de paquetes desacoplados con exportaciones separadas de cliente y servidor; la API del CMS se monta dentro del servidor web, sin servicios adicionales (ADR-011). & Cumplido \\
    OE-3 & Registro, verificación de correo, contraseña, enlace mágico y MFA; cuentas personales y de equipo con invitaciones, roles y permisos; MFA obligatorio en la consola. Autorización aplicada en la base de datos mediante RLS. & Cumplido; roles adaptados a pymes pendientes (F5) \\
    OE-4 & Suscripciones con Stripe (planes, \textit{checkout}, portal del cliente y \textit{webhooks}) operativas en modo de pruebas. & Parcial: planes de ejemplo y pruebas E2E de pago pendientes (F5, F6) \\
    OE-5 & CMS de datos integrado en la consola: explorador de tablas con CRUD, relaciones y vistas guardadas, gestión de usuarios y archivos, auditoría no falsificable, búsqueda global, RBAC propio y paneles configurables (RF-08 a RF-11). & Cumplido \\
    OE-6 & Modelo multicliente con tablas compartidas y RLS; más de dos mil pruebas pgTAP, incluidas pruebas de ataque entre clientes; las auditorías posteriores a cada corrección concluyeron que el aislamiento se mantiene. & Cumplido \\
    OE-7 & Configuración centralizada en ficheros \texttt{*.config.ts} y variables \texttt{VITE\_*}; consola por áreas configurable desde la propia interfaz; interfaz solo en español con la traducción como punto de extensión documentado. & Parcial: documentación de la configuración (F5) y comentarios didácticos (F4) pendientes \\
    OE-8 & Pruebas unitarias, de base de datos y E2E en integración continua. & Parcial: E2E de CU-01 a CU-06 y evaluación de la reutilización (F6) pendientes \\
    OE-9 & Memoria y manuales redactados; configuración de despliegue en la nube preparada. & En curso (F7, F8) \\
    \bottomrule
  \end{tabularx}
\end{table}

% PENDIENTE (P-06): valoración de la aportación propia frente a la reutilizada.
Más allá de la tabla, el resultado principal del trabajo es haber convertido
dos bases de código independientes, con pilas, sesiones y modelos de permisos
distintos, en un único producto coherente: un solo inicio de sesión, una sola
consola, una sola forma de escribir formularios y textos, y una sola base de
datos cuya seguridad se ha auditado y reforzado. La integración del CMS en la
consola ---la decisión de mayor riesgo del proyecto--- se completó sin
recurrir al recorte previsto.

\section{Lecciones aprendidas}
\label{sec:concl-lecciones}

\subsection{Esfuerzo dedicado}

Hasta la fecha de corte (20/10/2026) el repositorio registra 54
\textit{commits}, 21 decisiones de arquitectura y 55 incidencias documentadas
en la bitácora. Las fases F0 a F3 ---arnés, base SaaS, integración completa
del CMS, desmarcado y español--- se ejecutaron en tres jornadas de trabajo con
un total de 23 horas del autor (tabla~\ref{tab:plan-dedicacion}). Al cierre de
la F2 la base de datos contaba con más de dos mil pruebas pgTAP y la suite E2E
superaba las 140 pruebas.
\todo{Completar con la dedicación total del proyecto y las cifras finales de
pruebas al cerrar F6.}

\subsection{Sobre reutilizar código de terceros}

\begin{itemize}
  \item \textbf{Reutilizar no exime de verificar.} De las 55 incidencias, 29
        son de seguridad y 25 de ellas venían del código de partida: escalada
        de privilegios por un \texttt{UPDATE} sin \texttt{WITH CHECK} (B-05),
        comprobaciones de MFA y de bloqueo que fallaban en abierto (B-06,
        B-23), lectura de tablas internas (B-07, B-08, B-28), inyección SQL
        (B-24), auditoría falsificable (B-27) o escape de denegaciones por
        delegación (B-47). Una base reutilizable solo acelera si se ha
        verificado antes; en caso contrario, multiplica sus fallos en cada
        proyecto que la usa.
  \item \textbf{Un test en verde no demuestra corrección.} Un test heredado
        afirmaba precisamente el comportamiento inseguro (B-11). Hay que leer
        las aserciones, no solo los nombres de las pruebas.
  \item \textbf{El inventario de lo que se porta debe hacerse ruta a ruta.} La
        búsqueda global del CMS no figuraba en el plan inicial porque el
        inventario se hizo por paquetes (B-30).
  \item \textbf{El desmarcado es más que un \texttt{grep}.} La marca aparecía
        en expresiones regulares escapadas (B-01), en nombres de carpetas y
        rutas de importación (B-40) y en imágenes que ningún control textual
        puede leer (B-46). Los controles automáticos deben cubrir todas las
        variantes, y aun así hace falta una revisión visual.
\end{itemize}

\subsection{Sobre seguridad en la base de datos}

\begin{itemize}
  \item \textbf{Fallar en cerrado.} Toda comprobación de seguridad que
        encuentra un error debe denegar (B-06, B-23).
  \item \textbf{«Sin error» no es «éxito».} Una escritura filtrada por RLS
        afecta a cero filas sin producir ningún error, y hay que comprobarlo
        (B-49).
  \item \textbf{Defensa en profundidad.} Las reglas deben aplicarse en la
        base de datos y no solo en la API (B-31); una operación pensada para un
        registro debe comprobar que su condición identifica exactamente uno
        (B-19).
  \item \textbf{Las herramientas de diferencias no ven la seguridad.} La
        generación automática de migraciones omite opciones como
        \texttt{security\_invoker} de las vistas y los privilegios por columna
        (B-35); esas propiedades se verifican con pruebas, no con la
        herramienta.
  \item \textbf{«Local» no significa privado.} En un servidor con IP pública,
        el motor de contenedores puede saltarse el cortafuegos (B-04).
\end{itemize}

\subsection{Sobre el desarrollo asistido por IA}

\begin{itemize}
  \item \textbf{El cuello de botella se desplaza a la verificación.} La
        integración del CMS se estimó en semanas de desarrollo convencional y
        se ejecutó en días (sección~\ref{sec:plan-estimado-real}), pero buena
        parte del tiempo del autor se dedicó a revisar, probar y decidir.
  \item \textbf{La revisión humana sigue siendo necesaria.} El autor detectó
        con una prueba manual que la consola se cargaba sin segundo factor
        (B-14); una inyección SQL heredada pasó desapercibida en varias
        pasadas de los agentes por el mismo fichero (B-24), y otra latente,
        introducida al portar el código, solo la encontró una revisión
        posterior que buscaba otra cosa (B-45).
  \item \textbf{El nivel de revisión se decide por riesgo.} La verificación
        exhaustiva con varios agentes encontró brechas reales, pero su coste
        obligó a ajustarla (B-48); al reducirla, aparecieron pruebas
        desactualizadas en zonas no revisadas (B-51, B-54). La regla resultante
        es ejecutar siempre la suite completa antes de fusionar.
  \item \textbf{Las directrices escritas son el mecanismo de control.} Las
        reglas del arnés (leer el progreso antes de empezar, no reabrir
        decisiones sin un ADR, verificar antes de cerrar, registrar cada
        incidencia) son las que hacen que el trabajo de los agentes sea
        evaluable y que esta memoria pueda apoyarse en evidencias.
  \item \textbf{Cada incremento, su \textit{commit}.} Un corte de conexión
        con un agente trabajando dejó cambios de distinta naturaleza mezclados
        en el árbol de trabajo (B-15).
\end{itemize}

\todo{El autor añade su valoración personal: qué ha aprendido, qué haría de
otra manera y cómo ha cambiado su forma de trabajar con agentes de IA.}

\section{Trabajo futuro}
\label{sec:concl-futuro}

Las líneas de trabajo futuro se dividen entre lo que queda del plan del
proyecto y las mejoras detectadas durante su desarrollo.

\paragraph{Fases pendientes del plan.}
\begin{itemize}
  \item \textbf{F4, comentarios didácticos}: reescribir en español los
        comentarios heredados de todos los módulos, empezando por los
        esquemas de la base de datos y el \textit{middleware} de
        autorización.
  \item \textbf{F5, adaptación a pymes}: roles y permisos de equipo
        orientados a pymes, planes de suscripción de ejemplo, una plantilla de
        permisos inicial del CMS (soporte y gestor de contenidos) y la
        documentación completa de la configuración reutilizable.
  \item \textbf{F6, validación}: pruebas E2E de los casos de uso CU-01 a
        CU-06, pruebas de integración con los \textit{webhooks} de Stripe y la
        evaluación de la reutilización (CU-07).
  \item \textbf{F7, despliegue}: despliegue de demostración en la nube y
        manual de instalación probado desde cero.
\end{itemize}

\paragraph{Mejoras del producto.}
\begin{itemize}
  \item Completar el CMS: renombrar carpetas en el explorador de archivos,
        editores de texto enriquecido y de ficheros en los formularios, una
        interfaz para definir relaciones virtuales y ajustes por tipo de
        formato (moneda, fichero).
  \item Cerrar los pendientes heredados documentados: el total incorrecto de
        la lista de paneles y la denegación de almacenamiento sobre una
        subruta.
  \item Plantillas de correo en español para todos los mensajes y fechas
        del CMS según el idioma activo.
  \item Estabilizar las pruebas E2E inestables bajo carga.
  \item Ofrecer otros idiomas siguiendo el punto de extensión documentado.
  \item Un generador que cree un proyecto nuevo a partir de PymeKit con su
        configuración inicial, para reducir aún más el tiempo de puesta en
        marcha.
  \item Explorar una variante de PymeKit pensada para dar servicio a muchas
        pymes desde una sola instalación, apoyada en las cuentas de equipo
        que ya ofrece.
\end{itemize}


\newpage{\ }
\cleardoublepage

% Licencia
\chapter*{\rlap{Información sobre Licencia}}
\label{cap:licencia}
\phantomsection  % so hyperref creates bookmarks
\addcontentsline{toc}{chapter}{Información sobre Licencia}
%\label{label_fdl}

 \begin{center}

       Información sobre Licencia


\end{center}

\section*{Licencia de la documentación}

\todo{Confirmar con el tutor la licencia de la memoria (propuesta: CC BY-SA 4.0).}
Esta memoria se distribuye bajo la licencia \textit{Creative Commons
Reconocimiento-CompartirIgual 4.0 Internacional} (CC BY-SA
4.0)~\cite{cc-by-sa-4}. Se permite copiar, redistribuir y adaptar el
documento con cualquier finalidad, siempre que se reconozca la autoría, se
indiquen los cambios realizados y las obras derivadas se distribuyan con la
misma licencia.

Esta licencia se aplica únicamente al texto, las tablas y las figuras
originales de la memoria. No se extiende al software descrito en ella, a los
fragmentos de código que se reproducen a título ilustrativo ni a las marcas y
logotipos de terceros que puedan aparecer, que conservan sus propias
condiciones.

\section*{Licencia del software}

% PENDIENTE (P-06): forma de presentar las bases de partida, a confirmar con el tutor.
PymeKit no es un desarrollo íntegramente original. Como se explica en la
introducción y en el capítulo de antecedentes, el software parte de dos bases
de código de terceros distribuidas con \textbf{licencia comercial} ---una base
de aplicación SaaS y un CMS para la base de datos---, que se han integrado,
adaptado y extendido en este trabajo. El autor adquirió ambas con su licencia
correspondiente, que permite usarlas como base de productos propios pero no
publicar ni redistribuir su código fuente.
\todo{Confirmar este resumen con el texto de las licencias adquiridas.}

Por ese motivo:

\begin{itemize}
  \item el repositorio del proyecto es \textbf{privado} y su código fuente
        \textbf{no se publica ni se redistribuye} bajo ninguna licencia
        libre;
  \item el código se pone a disposición del tribunal y del tutor únicamente
        para la evaluación del Trabajo Fin de Grado;
        \todo{Confirmar con el tutor el mecanismo de entrega del código al
        tribunal (acceso al repositorio privado o copia en soporte digital).}
  \item el origen del código reutilizado está documentado internamente en el
        repositorio, junto con el registro de los módulos portados y de las
        modificaciones realizadas sobre ellos; y
  \item las aportaciones propias del trabajo se integran en ese mismo
        repositorio y quedan sujetas a sus mismas condiciones de distribución.
\end{itemize}

Las dependencias de código abierto que usa el proyecto (por ejemplo, React,
TanStack, Supabase, Hono, Drizzle, Tailwind CSS o Playwright) se rigen por sus
respectivas licencias, recogidas en sus paquetes.


\newpage{\ }
\cleardoublepage

% Bibliografía
\thispagestyle{empty}
\setcitestyle{aysep={}}
\bibliographystyle{apalike7}
\bibliography{bibliography}

\newpage{\ }
\cleardoublepage

 
    % BACKMATTER:
    \part{Anexos}
    \null\vfill
\noindent En esta última parte quedarán recogidas los manuales necesarios para el manejo de la aplicación resultado del desarrollo. Si se ha realizado algún tipo de evaluación de la solución proporcionada, más allá de las pruebas del sistema, también deberá venir recogida en un capítulo separado dentro de esta parte. Pueden consultarse diversos tipos de evaluaciones sobre sistemas de información en \cite{hevner2004}: casos de estudio, análisis estático, análisis dinámico, simulación, experimento controlado, etc.
\vfill


\appendix
%\newpage{\ }
%\thispagestyle{empty}
%\addappheadtotoc
%\appendixpage

\chapter{Manual del desarrollador}
% ------------------------------------------------------------------------------
% Anexo · Manual del desarrollador
% Fuentes: package.json (raíz y apps/web), README.md, AGENTS.md,
% apps/web/supabase/AGENTS.md, apps/web/src/config/*.config.ts, apps/web/.env*,
% packages/i18n/README.md, .github/workflows/ci.yml. Estado: 20/10/2026.
% Los bloques lstlisting solo contienen caracteres ASCII (listings no admite
% UTF-8 multibyte sin configuración adicional).
% ------------------------------------------------------------------------------
\label{cap:manual-desarrollador}

Este manual se dirige a quien vaya a evolucionar PymeKit o a construir sobre
él un SaaS nuevo. Explica cómo preparar el entorno, cómo está organizado el
código, cómo se trabaja con la base de datos, cómo se añade una tabla de
negocio que aparezca en la consola, qué se configura sin tocar el núcleo y cómo
se prueba y se despliega. Para la arquitectura y las decisiones que hay detrás
de cada pieza, véanse los capítulos de Diseño y de Construcción del Sistema.

\section{Introducción}
\label{sec:dg-introduccion}

PymeKit es un monorepo con una única aplicación web que contiene el área
pública, el área privada de los clientes, la consola de administración y la
API del CMS, más un conjunto de paquetes reutilizables. Todo comparte una sola
base de datos PostgreSQL gestionada con Supabase, cuyas políticas RLS son la
frontera de seguridad. Las directrices completas para personas y agentes de IA
están en el fichero \texttt{AGENTS.md} de la raíz y en los de cada paquete;
este manual resume lo imprescindible.

\section{Preparación del entorno de trabajo}
\label{sec:dg-entorno}

\subsection{Requisitos}

\begin{itemize}
  \item Node.js 20.10 o superior (se recomienda la versión 24).
  \item pnpm 11 (la versión exacta se fija en el campo
        \texttt{packageManager} del \texttt{package.json} raíz).
  \item Docker, que usa la instancia local de Supabase. El usuario debe poder
        ejecutar Docker sin \texttt{sudo}.
  \item Git y acceso al repositorio privado del proyecto.
  \item Opcional: Stripe CLI (se ejecuta en un contenedor) para recibir los
        \textit{webhooks} de pago en local.
\end{itemize}

\paragraph{Advertencia de seguridad.} Si se desarrolla en un servidor con IP
pública, los puertos que publica Docker pueden quedar expuestos aunque haya un
cortafuegos, porque Docker inserta sus propias reglas. Hay que filtrar en la
cadena \texttt{DOCKER-USER} o enlazar los servicios a \texttt{127.0.0.1} y
acceder mediante un túnel SSH (ADR-010, incidencia B-04).

\subsection{Arranque local}

\begin{lstlisting}[language=bash]
git clone <url-del-repositorio> pymekit && cd pymekit
pnpm install                # dependencias de todo el monorepo
pnpm supabase:web:start     # Postgres, Auth, Storage y Studio en Docker
pnpm dev                    # aplicacion web en http://localhost:3100
\end{lstlisting}

El primer arranque de Supabase aplica todas las migraciones y el
\textit{seed}. Una vez en marcha, quedan disponibles los servicios de la
tabla~\ref{tab:dg-puertos}.

\begin{table}[htbp]
  \centering
  \small
  \caption{Servicios del entorno local.}
  \label{tab:dg-puertos}
  \begin{tabularx}{\textwidth}{@{}l>{\raggedright\arraybackslash}p{4.4cm}X@{}}
    \toprule
    \textbf{Servicio} & \textbf{Dirección} & \textbf{Uso} \\
    \midrule
    Aplicación web & \texttt{http://localhost:3100} & Web pública, área privada, consola y API del CMS (ADR-009) \\
    API de Supabase & \texttt{http://127.0.0.1:54321} & Auth, Storage y API de datos \\
    PostgreSQL & \texttt{127.0.0.1:54322} & Conexión directa a la base de datos \\
    Supabase Studio & \texttt{http://127.0.0.1:54323} & Herramienta de desarrollo para inspeccionar la base de datos \\
    Correo de pruebas & \texttt{http://127.0.0.1:54324} & Bandeja donde llegan los correos de verificación, invitaciones y enlaces \\
    \bottomrule
  \end{tabularx}
\end{table}

\paragraph{Usuarios de desarrollo.} El \textit{seed} crea, con la contraseña
\texttt{testingpassword}, las cuentas \texttt{test@pymekit.test},
\texttt{owner@pymekit.test}, \texttt{member@pymekit.test},
\texttt{custom@pymekit.test}, \texttt{super-admin@pymekit.test} y
\texttt{cms-staff@pymekit.test} (personal del CMS con el rol «Soporte»). Las dos
últimas tienen un factor TOTP verificado cuyo secreto está en
\texttt{apps/web/supabase/seed.sql}; hay que darlo de alta en una aplicación de
códigos para entrar en la consola. Son credenciales públicas de desarrollo,
válidas solo contra la instancia local (B-03).

\paragraph{Variables de entorno.} Los valores públicos van en
\texttt{apps/web/.env} (compartido) y \texttt{.env.development},
\texttt{.env.test} o \texttt{.env.production}. Los secretos
---\texttt{SUPABASE\_SECRET\_KEY}, \texttt{STRIPE\_SECRET\_KEY},
\texttt{STRIPE\_WEBHOOK\_SECRET}, credenciales del correo--- se ponen en
\texttt{apps/web/.env.local}, que está excluido del repositorio, o en las
variables de la plataforma de despliegue. Nunca se suben al repositorio.

\paragraph{Pagos en local.} Para probar la facturación hay que añadir las
claves de prueba de Stripe (\texttt{VITE\_STRIPE\_PUBLISHABLE\_KEY} y los dos
secretos anteriores) y reenviar los \textit{webhooks} a la aplicación con
\texttt{pnpm stripe:listen}, que publica los eventos en
\texttt{/api/billing/webhook}.

\section{Estructura del monorepo}
\label{sec:dg-estructura}

La tabla~\ref{tab:dg-estructura} resume dónde está cada cosa.

\begin{table}[htbp]
  \centering
  \footnotesize
  \caption{Estructura del monorepo.}
  \label{tab:dg-estructura}
  \begin{tabularx}{\textwidth}{@{}>{\raggedright\arraybackslash}p{5.3cm}X@{}}
    \toprule
    \textbf{Ruta} & \textbf{Contenido} \\
    \midrule
    \texttt{apps/web} & Aplicación TanStack Start (React 19): rutas por fichero en \texttt{src/routes}, \textit{server functions} en ficheros \texttt{*.functions.ts}, configuración en \texttt{src/config} \\
    \texttt{apps/web/src/routes/\_marketing} & Web pública: inicio, precios, blog, contacto, páginas legales \\
    \texttt{apps/web/src/routes/\_authenticated} & Área privada: panel, ajustes, miembros, facturación, creación de equipos \\
    \texttt{apps/web/src/routes/admin} & Consola: inicio y gestión de cuentas (\texttt{/admin}) y pantallas del CMS (\texttt{/admin/cms}) \\
    \texttt{apps/web/src/routes/api} & Rutas de servidor: API del CMS (\texttt{/api/cms/\$}), \textit{webhooks} de pagos y de base de datos, estado de salud \\
    \texttt{apps/web/supabase} & Esquemas declarativos, migraciones, \textit{seed}, plantillas de correo y pruebas pgTAP de toda la base de datos \\
    \texttt{apps/e2e} & Pruebas Playwright de la web y del CMS \\
    \texttt{packages/features/*} & Autenticación, cuentas, equipos, administración y notificaciones \\
    \texttt{packages/billing/*} & Núcleo de facturación, pasarela y proveedor Stripe \\
    \texttt{packages/cms/*} & CMS: servicios Drizzle y rutas Hono (solo servidor), componentes de interfaz, filtros, formateadores, permisos \\
    \texttt{packages/function-middleware} & \textit{Middleware} de las \textit{server functions}: usuario autenticado, miembro de equipo, super-admin, rol mínimo y permiso \\
    \texttt{packages/supabase} & Clientes de Supabase y tipos generados de la base de datos \\
    \texttt{packages/ui}, \texttt{packages/i18n} & Componentes de interfaz y textos (catálogos de mensajes) \\
    \texttt{packages/*} (resto) & Correo, plantillas de correo, analítica, monitorización, OTP, políticas, utilidades compartidas \\
    \texttt{tooling}, \texttt{turbo/generators} & Configuración compartida y generadores de paquetes \\
    \texttt{docs/tfg}, \texttt{scripts/tfg} & Documentación interna del proyecto y scripts de control \\
    \bottomrule
  \end{tabularx}
\end{table}

Los paquetes internos se importan con el \textit{scope} \texttt{@pymekit/*}
(los del CMS, \texttt{@pymekit/cms-*}). Los módulos que solo deben ejecutarse
en el servidor llevan el sufijo \texttt{.server.ts}, y las rutas Hono y los
servicios del CMS nunca se importan desde componentes.

\section{Consideraciones generales sobre el desarrollo}
\label{sec:dg-consideraciones}

\begin{itemize}
  \item \textbf{Server functions.} Se definen con \texttt{createServerFn} en
        ficheros \texttt{*.functions.ts}, con la cadena
        \texttt{.middleware(...).validator(esquemaZod).handler(...)}, y se
        llaman desde el cliente con \texttt{useMutation} y
        \texttt{useServerFn}.
  \item \textbf{Clientes de base de datos.} El cliente de servidor y el del
        navegador aplican RLS. El cliente de administración la ignora: se usa
        lo mínimo posible, con validación manual y un comentario que lo
        justifique.
  \item \textbf{Formularios}: TanStack Form con los campos de
        \texttt{@pymekit/ui/field} y validación con Zod, también en el CMS.
  \item \textbf{Textos}: ningún texto visible ---tampoco \texttt{aria-label},
        \texttt{title}, \texttt{alt} ni \texttt{placeholder}--- se escribe en
        el código; se añade una clave en
        \texttt{packages/i18n/src/messages/es/} (sección~\ref{sec:dg-i18n}).
  \item \textbf{Comentarios} en español didáctico: cabecera en cada fichero
        con lógica, documentación de las funciones exportadas y justificación
        de toda decisión de seguridad. Los puntos que se citan en la memoria
        llevan la etiqueta \texttt{[TFG]} con el requisito.
  \item \textbf{Registro}: \texttt{getLogger()} de
        \texttt{@pymekit/shared/logger}; nunca \texttt{console.log} en
        producción.
  \item \textbf{Git}: ramas \texttt{fase-N/tema}, \textit{commits} en español
        con \textit{Conventional Commits} y el módulo como ámbito.
\end{itemize}

\section{Base de datos}
\label{sec:dg-bd}

\subsection{Esquemas, migraciones y tipos}

El esquema se escribe de forma declarativa en
\texttt{apps/web/supabase/schemas/}, en ficheros numerados cuyo prefijo fija el
orden de dependencias: \texttt{00}--\texttt{18} plataforma SaaS,
\texttt{19} blog y \texttt{20}--\texttt{55} CMS (esquema \texttt{cms}). Las
migraciones se generan a partir de los esquemas, nunca a mano:

\begin{lstlisting}[language=bash]
# 1. Editar o crear el fichero de esquema (schemas/NN-nombre.sql)
# 2. Generar la migracion a partir de la diferencia
pnpm --filter web run supabase:db:diff -f nombre_del_cambio
# 3. Aplicarla y regenerar los tipos TypeScript
pnpm --filter web supabase migrations up
pnpm supabase:web:typegen
# 4. Comprobar que no queda diferencia (no debe mencionar el cambio)
pnpm --filter web run supabase:db:diff
\end{lstlisting}

La herramienta de diferencias no genera los \texttt{revoke} por tabla, los
\texttt{grant} por columna ni la opción \texttt{security\_invoker} de las
vistas (B-35): hay que escribirlos a mano en la migración y comprobarlos con
pgTAP. Las migraciones nunca son destructivas. \texttt{pnpm
supabase:web:reset} recrea la base de datos local desde cero, con las
migraciones y el \textit{seed}, y conviene ejecutarlo antes de fiarse de una
ejecución de las pruebas.

\subsection{Reglas de seguridad}

\begin{itemize}
  \item RLS activada en toda tabla nueva, con políticas basadas en las
        funciones auxiliares existentes (\texttt{has\_role\_on\_account},
        \texttt{has\_permission}, \texttt{is\_mfa\_compliant}).
  \item \texttt{revoke all} a \texttt{anon}, \texttt{authenticated} y
        \texttt{service\_role} tras crear la tabla, y después solo los verbos
        que use la funcionalidad. \texttt{UPDATE} siempre por columnas, sin
        incluir \texttt{id} ni \texttt{account\_id}.
  \item Ninguna función \texttt{security definer} sin controles de acceso
        explícitos.
  \item Todo cambio de políticas, \texttt{grant}, funciones
        \texttt{security definer} o vistas lleva pruebas pgTAP y una revisión
        adversarial de RLS.
\end{itemize}

\subsection{Pruebas pgTAP}

Las pruebas están en \texttt{apps/web/supabase/tests/database/*.test.sql} y se
ejecutan con \texttt{pnpm supabase:web:test}. El fichero
\texttt{00000-pymekit-helpers.sql} define el esquema auxiliar
\texttt{pymekit.*}, con funciones como \texttt{pymekit.authenticate\_as(...)}
para simular a un usuario concreto; las del CMS usan además
\texttt{00001-cms-test-helpers.sql}. Una prueba de aislamiento típica se
autentica como miembro de un equipo e intenta leer o modificar datos de otro.

\section{Añadir una tabla de negocio}
\label{sec:dg-tabla}

Este es el caso más habitual al construir un SaaS sobre PymeKit. Como ejemplo,
una tabla de citas (\texttt{appointments}) que pertenece a una cuenta de
equipo y debe gestionarse desde la consola en un área «Agenda».

\paragraph{1. Esquema con RLS.} Se crea
\texttt{apps/web/supabase/schemas/56-appointments.sql} a partir de la
plantilla de tabla de \texttt{apps/web/supabase/AGENTS.md}:

\begin{lstlisting}[language=SQL]
create table if not exists public.appointments (
  id uuid not null default extensions.uuid_generate_v4(),
  account_id uuid references public.accounts(id) on delete cascade not null,
  customer_name varchar(255) not null,
  starts_at timestamptz not null,
  created_at timestamptz not null default now(),
  updated_at timestamptz not null default now(),
  created_by uuid references auth.users on delete set null,
  updated_by uuid references auth.users on delete set null,
  primary key (id)
);

alter table public.appointments enable row level security;
revoke all on public.appointments from anon, authenticated, service_role;
grant select, insert, update, delete on public.appointments to service_role;
grant select, insert on public.appointments to authenticated;

create policy appointments_read on public.appointments for select
  to authenticated using (public.has_role_on_account(account_id));
create policy appointments_insert on public.appointments for insert
  to authenticated with check (public.has_role_on_account(account_id));
create policy restrict_mfa_appointments on public.appointments
  as restrictive to authenticated using (public.is_mfa_compliant());
\end{lstlisting}

La plantilla completa añade además los \textit{triggers} de marcas de tiempo y
de trazabilidad y un índice por \texttt{account\_id}.

\paragraph{2. Migración, tipos y pruebas.} Se genera la migración con
\texttt{supabase:db:diff -f appointments}, se completan a mano los
\texttt{revoke}, se aplica, se regeneran los tipos y se añade
\texttt{tests/database/appointments.test.sql} comprobando que un miembro de
otro equipo no puede leer ni insertar citas ajenas. Si la tabla está en
\texttt{public}, hay que confirmar que \texttt{anon-surface.test.sql} sigue
pasando: el rol anónimo solo puede leer el blog.

\paragraph{3. Registro en el CMS.} Para que la consola conozca la tabla, se
sincronizan sus metadatos. Puede hacerse desde la interfaz («Ajustes»
$\rightarrow$ «Recursos» $\rightarrow$ «Sincronizar tablas») o, para que
quede en todas las instalaciones, en la propia migración:

\begin{lstlisting}[language=SQL]
select cms.sync_managed_tables('public', 'appointments');

update cms.table_metadata
set display_name = 'Citas',
    ui_config = coalesce(ui_config, '{}'::jsonb)
      || jsonb_build_object('navigation_group', 'Agenda')
where schema_name = 'public' and table_name = 'appointments';
\end{lstlisting}

La clave \texttt{navigation\_group} de \texttt{ui\_config} es el área de la
barra lateral: si no existe, se crea. La sincronización conserva lo ya
configurado, de modo que puede repetirse sin perder los cambios. Los nombres
de las columnas, sus descripciones (la ayuda «?»), su formato y el diseño de
la ficha se ajustan después en «Ajustes» $\rightarrow$ «Recursos», sin código
(el ejemplo de la migración \texttt{20261002140000\_cms\_navigation\_groups.sql}
muestra cómo fijarlos también por migración).

\paragraph{4. Permisos del personal.} El super-administrador ve la tabla
inmediatamente. Para el resto del personal, se crean o amplían roles en
«Ajustes» $\rightarrow$ «Roles y permisos» con permisos de datos sobre
\texttt{public.appointments}. El CMS accede a las tablas mediante funciones
del esquema \texttt{cms} que comprueban esos permisos y registran cada escritura
en la auditoría; las políticas RLS de la tabla siguen protegiendo el acceso de
los clientes desde la aplicación.

\paragraph{5. Interfaz del cliente.} Si los clientes también deben ver las
citas, se añade una ruta en \texttt{apps/web/src/routes/\_authenticated}, una
\textit{server function} con \texttt{teamAccountFunctionMiddleware} y los
textos en el catálogo de mensajes.

\section{Configuración reutilizable}
\label{sec:dg-configuracion}

Adaptar PymeKit a un producto nuevo empieza por la configuración, sin tocar el
núcleo (RNF-01). Los ficheros de \texttt{apps/web/src/config} validan sus
valores con Zod al arrancar, de modo que un error de configuración se detecta
de inmediato (tabla~\ref{tab:dg-configuracion}).

\begin{table}[htbp]
  \centering
  \footnotesize
  \caption{Ficheros de configuración y variables principales.}
  \label{tab:dg-configuracion}
  \begin{tabularx}{\textwidth}{@{}>{\raggedright\arraybackslash}p{4.3cm}X@{}}
    \toprule
    \textbf{Fichero} & \textbf{Qué configura (variables)} \\
    \midrule
    \texttt{app.config.ts} & Nombre, título, descripción, URL, tema y colores (\texttt{VITE\_PRODUCT\_NAME}, \texttt{VITE\_SITE\_TITLE}, \texttt{VITE\_SITE\_DESCRIPTION}, \texttt{VITE\_SITE\_URL}, \texttt{VITE\_DEFAULT\_THEME\_MODE}, \texttt{VITE\_THEME\_COLOR}, \texttt{VITE\_THEME\_COLOR\_DARK}); en producción exige HTTPS \\
    \texttt{account-mode.config.ts} & Modo de cuentas: \texttt{personal-only}, \texttt{organizations-only} o \texttt{hybrid} (\texttt{VITE\_ACCOUNT\_MODE}); es la fuente única de qué superficies existen \\
    \texttt{feature-flags.config.ts} & Interruptores: selector de tema, borrado de cuentas personales y de equipo, creación de equipos, facturación personal y de equipo, notificaciones (\texttt{VITE\_ENABLE\_*}) \\
    \texttt{auth.config.ts} & Métodos de acceso: contraseña, enlace mágico, CAPTCHA (\texttt{VITE\_AUTH\_PASSWORD}, \texttt{VITE\_AUTH\_MAGIC\_LINK}, \texttt{VITE\_CAPTCHA\_SITE\_KEY}) \\
    \texttt{billing.config.ts} & Proveedor (\texttt{VITE\_BILLING\_PROVIDER}), planes, periodicidades, precios e identificadores de precio de la pasarela \\
    \texttt{paths.config.ts} & Rutas de la aplicación (inicio de sesión, panel, ajustes, facturación\ldots) \\
    \texttt{navigation.config.tsx}, \texttt{marketing-navigation.config.tsx} & Menús del área privada y de la web pública \\
    \texttt{tfg.config.ts} & Bloque «Proyecto académico» y recorrido de la demostración en la página de inicio (\texttt{VITE\_SHOW\_TFG\_INFO}); una pyme real lo desactiva con \texttt{false} \\
    \bottomrule
  \end{tabularx}
\end{table}

A esta configuración por ficheros se suma la configuración \emph{como datos}
de la consola ---nombres de tablas y columnas, áreas, formatos, diseño de las
fichas y roles del personal---, que se cambia desde la propia interfaz
(sección~\ref{sec:dg-tabla}). El puerto de la aplicación (3100) se fija en
\texttt{vite.config.ts} y en la variable \texttt{PORT}.

\section{Internacionalización}
\label{sec:dg-i18n}

La interfaz solo ofrece el español, pero todos sus textos salen de catálogos
de mensajes (ADR-021). Para añadir un texto, se crea la clave en inglés en
\texttt{packages/i18n/src/messages/es/<namespace>.json} con el valor en español
y se usa con el componente \texttt{Trans} o con las funciones de mensajes. Las
etiquetas de los valores de los enumerados de la base de datos van en
\texttt{common.enums.<enumerado>.<valor>}.

Añadir un idioma es un punto de extensión documentado en
\texttt{packages/i18n/README.md}. En resumen: copiar
\texttt{src/messages/es/} con el código del idioma nuevo y traducir los
valores; registrarlo en \texttt{src/messages/index.ts} y en
\texttt{supportedLocales} de \texttt{src/config.ts}; crear las traducciones de
los correos en \texttt{packages/email-templates}; elegir el idioma de los
mensajes de validación; volver a mostrar el selector de idioma; y ejecutar
\texttt{pnpm --filter @pymekit/i18n test:unit}, cuya prueba comprueba que el
nuevo catálogo tiene exactamente las mismas claves y marcadores que el
español.

\section{Pruebas}
\label{sec:dg-pruebas}

\begin{lstlisting}[language=bash]
pnpm typecheck              # tipos de todos los paquetes
pnpm lint:fix && pnpm format:fix
pnpm test:unit              # pruebas unitarias (Vitest)
pnpm supabase:web:test      # pruebas de base de datos (pgTAP)

# Pruebas E2E: contra la build de test, no contra el servidor de desarrollo
pnpm --filter web build:test
pnpm --filter web start:test &
pnpm --filter web-e2e test
\end{lstlisting}

Las pruebas E2E se lanzan contra la \textit{build} de test porque el inicio de
sesión automatizado falla contra el servidor de desarrollo (B-13); con
\texttt{pnpm --filter web-e2e test:ui} se abre la interfaz gráfica de Playwright.
Antes de dar por terminado un cambio se completa además la verificación del
proyecto: los controles automáticos del repositorio (\texttt{scripts/tfg/}),
revisión del cambio y, si se tocó la base de
datos, revisión de RLS. La integración continua
(\texttt{.github/workflows/ci.yml}) repite los controles, los tipos, el
\textit{lint}, el formato y las pruebas unitarias en cada \textit{push}, y la
base de datos y las E2E en las \textit{pull requests} hacia \texttt{main}.

\section{Instrucciones para la construcción y el despliegue}
\label{sec:dg-construccion}

\begin{lstlisting}[language=bash]
pnpm build                  # construye todos los paquetes y la web
pnpm --filter web start     # sirve la build de produccion (.output/)
\end{lstlisting}

El servidor de producción lee las variables de \texttt{apps/web/.env} y del
entorno, escucha en el puerto de \texttt{PORT} y expone
\texttt{/api/healthcheck} para las comprobaciones de salud. La base de datos de
un proyecto de Supabase en la nube se actualiza con
\texttt{pnpm --filter web supabase:deploy}, que enlaza el proyecto
(\texttt{SUPABASE\_PROJECT\_REF}) y aplica las migraciones pendientes. El
fichero \texttt{railway.json} define la construcción y el arranque en la
plataforma de despliegue. El procedimiento completo, con la configuración de
Supabase y de Stripe en la nube, se describe en el capítulo de Despliegue del
Sistema (capítulo~\ref{cap:despliegue}).

\newpage{\ }
\cleardoublepage

\chapter{Manual de usuario}
% ------------------------------------------------------------------------------
% Anexo · Manual de usuario
% Fuentes: rutas de apps/web/src/routes, navegación de la consola
% (apps/web/src/components/admin/admin-navigation.ts) y textos reales de la
% interfaz (packages/i18n/src/messages/es/*.json). Estado: 20/10/2026.
% ------------------------------------------------------------------------------
\label{cap:manual-usuario}

Este manual explica cómo usar PymeKit desde el navegador. Se organiza por
perfiles de usuario, de menos a más privilegios: visitante, cliente o miembro
de un equipo, personal de la consola y super-administrador. Los nombres de los
menús y los botones se citan tal como aparecen en la interfaz, entre comillas
angulares.

\section{Introducción}
\label{sec:ug-introduccion}

PymeKit es la base de la web de una pyme: una parte pública con página de
presentación y blog, un área privada para los clientes y una consola de
administración para el personal, con un CMS que permite gestionar los datos,
los usuarios y los archivos sin programar. Toda la interfaz está en español.

La tabla~\ref{tab:ug-perfiles} resume qué puede hacer cada perfil y por dónde
entra.

\begin{table}[htbp]
  \centering
  \small
  \caption{Perfiles de usuario de PymeKit.}
  \label{tab:ug-perfiles}
  \begin{tabularx}{\textwidth}{@{}>{\raggedright\arraybackslash}p{3cm}>{\raggedright\arraybackslash}p{3.2cm}X@{}}
    \toprule
    \textbf{Perfil} & \textbf{Entrada} & \textbf{Qué puede hacer} \\
    \midrule
    Visitante & \texttt{/} & Ver la página de presentación, los precios, el blog, las preguntas frecuentes, el formulario de contacto y las páginas legales. \\
    Cliente o miembro de equipo & \texttt{/auth/sign-in} $\rightarrow$ \texttt{/dashboard} & Gestionar su cuenta y su seguridad, crear equipos, invitar a otras personas y contratar un plan. \\
    Personal de la consola & \texttt{/admin} & Consultar o editar los datos, usuarios y archivos que su rol del CMS le permita; ver paneles y auditoría. \\
    Super-administrador & \texttt{/admin} & Todo lo anterior, más la gestión de las cuentas de la plataforma y los ajustes de la consola. \\
    \bottomrule
  \end{tabularx}
\end{table}

\section{Requisitos e instalación}
\label{sec:ug-requisitos}

El usuario final no instala nada: basta un navegador web actual con
JavaScript activado, en ordenador o en móvil (la interfaz se adapta al tamaño
de la pantalla). Para usar el segundo factor de autenticación hace falta una
aplicación de códigos temporales (TOTP) en el móvil, como las que ofrecen los
principales fabricantes de sistemas operativos o los gestores de contraseñas.

La instalación de la plataforma, local o en la nube, la realiza un
desarrollador y se describe en el manual del desarrollador
(anexo~\ref{cap:manual-desarrollador}) y en el capítulo de Despliegue del
Sistema.

\section{Visitante}
\label{sec:ug-visitante}

La parte pública ofrece:

\begin{itemize}
  \item \textbf{Página de inicio} (\texttt{/}): presentación del producto y
        acceso a «Iniciar sesión» y «Registrarse». En la instalación de
        demostración incluye además un bloque «Proyecto académico» con un
        «Recorrido de la demostración» en cinco pasos y la tarjeta de prueba
        para los pagos.
  \item \textbf{Precios} (\texttt{/pricing}): los planes disponibles, con
        periodicidad mensual o anual.
  \item \textbf{Blog} (\texttt{/blog}): las entradas publicadas, con
        categorías y etiquetas. Cada entrada tiene su propia dirección
        (\texttt{/blog/<slug>}). El contenido se gestiona desde la consola
        (sección~\ref{sec:ug-datos}).
  \item \textbf{Contacto} (\texttt{/contact}), \textbf{preguntas frecuentes}
        (\texttt{/faq}) y las páginas legales: términos del servicio,
        política de privacidad y política de \textit{cookies}.
\end{itemize}

\section{Cliente o miembro de un equipo}
\label{sec:ug-cliente}

\subsection{Registro e inicio de sesión}
\label{sec:ug-registro}

\begin{enumerate}
  \item En la página de inicio, pulse «Registrarse». En «Crea una cuenta»,
        introduzca su correo electrónico y una contraseña de al menos 8
        caracteres, repítala y pulse «Registrarse».
  \item Recibirá un correo con un enlace de verificación. Ábralo para
        confirmar la dirección; hasta entonces no podrá entrar.
  \item Para entrar en otra ocasión, pulse «Iniciar sesión», escriba su correo
        y su contraseña y pulse «Iniciar sesión». Si la instalación tiene
        activado el enlace mágico, puede pedir en su lugar un enlace de acceso
        por correo.
  \item Si ha olvidado la contraseña, pulse «¿Has olvidado tu contraseña?»,
        indique su correo y siga el enlace que recibirá para establecer una
        nueva.
\end{enumerate}

Tras entrar llegará a su \textbf{panel} (\texttt{/dashboard}). La barra
lateral da acceso al panel, a los ajustes y, si los hay, a los miembros y a la
facturación de la cuenta activa.

\subsection{Verificación en dos pasos (MFA)}
\label{sec:ug-mfa}

Se recomienda proteger la cuenta con un segundo factor. Es obligatorio para
quien vaya a usar la consola de administración.

\begin{enumerate}
  \item Abra los ajustes de su cuenta (menú de su usuario $\rightarrow$
        «Ajustes de la cuenta»).
  \item En la sección «Autenticación multifactor», pulse «Configurar un nuevo
        factor».
  \item Escanee el código QR con su aplicación de códigos temporales, dé un
        nombre al factor (por ejemplo, el de su móvil) e introduzca el código
        de seis dígitos que muestra la aplicación. Al verificarlo aparecerá el
        aviso «La autenticación multifactor está activada».
  \item A partir de entonces, cada vez que inicie sesión se le pedirá el
        código en la pantalla «Verifica tu cuenta» (\texttt{/auth/verify}).
\end{enumerate}

Desde la misma sección puede desactivar un factor («Desactivar factor»). Si
pierde el acceso a su aplicación, un super-administrador puede retirarle el
factor desde la consola (sección~\ref{sec:ug-superadmin}).

\subsection{Ajustes de la cuenta}
\label{sec:ug-cuenta}

En «Ajustes de la cuenta» puede cambiar su nombre y su foto, actualizar su
correo electrónico («Actualiza tu correo electrónico»), cambiar la contraseña
(«Actualiza tu contraseña») y gestionar sus factores de autenticación. En la
«Zona de peligro», «Eliminar tu cuenta» borra la cuenta personal de forma
irreversible, si la instalación lo permite.

\subsection{Equipos e invitaciones}
\label{sec:ug-equipos}

Un \textbf{equipo} es una cuenta compartida por varias personas: por ejemplo,
una tienda o un departamento de la pyme cliente. Cada equipo tiene sus
miembros, sus roles y, si procede, su propia suscripción. Los datos de un
equipo solo son visibles para sus miembros.

\begin{itemize}
  \item \textbf{Crear un equipo.} En el selector de cuenta de la barra
        lateral, elija «Nuevo equipo» (o vaya a \texttt{/create-team}), escriba
        el «Nombre del equipo» y pulse «Crear equipo». Quien lo crea es su
        propietario.
  \item \textbf{Cambiar de cuenta.} El selector de cuenta muestra «Tus
        equipos»; al elegir uno, el panel, los ajustes, los miembros y la
        facturación pasan a ser los de ese equipo.
  \item \textbf{Invitar.} En «Miembros» (\texttt{/settings/members}), pulse
        «Invitar miembros», escriba uno o varios correos, elija el rol de cada
        persona y pulse «Enviar invitaciones». Las invitaciones aparecen en
        «Invitaciones pendientes», desde donde se pueden reenviar, renovar o
        eliminar.
  \item \textbf{Aceptar una invitación.} La persona invitada recibe un correo
        con un enlace (\texttt{/join}). Si no tiene cuenta, se registra; si ya
        la tiene, inicia sesión y pulsa «Aceptar invitación».
  \item \textbf{Roles.} Por defecto existen dos: «Propietario», con todos los
        permisos (gestionar roles, facturación, ajustes, miembros e
        invitaciones), y «Miembro», que puede usar el equipo pero no
        administrarlo. Un miembro sin permiso de facturación no puede
        gestionar la suscripción: la interfaz no se lo ofrece y la base de
        datos lo rechaza aunque lo intente por otra vía.
  \item \textbf{Otras acciones.} Desde «Ajustes» del equipo se cambian su
        nombre y su logotipo; el propietario puede transferir la propiedad a
        otro miembro o eliminar el equipo, y cualquier miembro puede
        abandonarlo («Abandonar equipo»).
\end{itemize}

\subsection{Facturación y suscripciones}
\label{sec:ug-facturacion}

La pestaña «Facturación» (\texttt{/settings/billing}) muestra la suscripción
de la cuenta activa, personal o de equipo.

\begin{enumerate}
  \item En «Elige tu plan preferido», seleccione un plan y su periodicidad y
        pulse «Continuar al pago».
  \item Complete el formulario de pago de la pasarela. En la instalación de
        demostración los pagos funcionan en \textbf{modo de prueba}: use la
        tarjeta \texttt{4242 4242 4242 4242}, con cualquier fecha de caducidad
        futura y cualquier código CVC. No se realiza ningún cargo real.
  \item Al terminar volverá a PymeKit (\texttt{/settings/billing/return}),
        que confirmará la operación. Unos segundos después, la pestaña
        mostrará el plan contratado y su estado.
  \item Para cambiar de plan, actualizar el método de pago o cancelar,
        pulse «Ir al portal de facturación».
\end{enumerate}

\section{Personal de la consola}
\label{sec:ug-personal}

La consola de administración (\texttt{/admin}) está reservada al
super-administrador y al personal al que se ha concedido acceso al CMS. Cada
miembro del personal tiene un rol del CMS que determina qué ve y qué puede
hacer; las entradas de menú y los botones que su rol no permite no se
muestran.

\subsection{Acceso}

\begin{enumerate}
  \item Inicie sesión con su cuenta habitual.
  \item Verifique el segundo factor. La consola exige una sesión verificada
        en dos pasos: si aún no tiene un factor configurado, la consola le
        mostrará el aviso «Se requiere verificación en dos pasos o una cuenta
        de CMS activa»; configúrelo como se indica en la
        sección~\ref{sec:ug-mfa}.
  \item Vaya a \texttt{/admin}.
\end{enumerate}

\subsection{Organización de la consola}

La barra lateral tiene dos bloques (figura~\ref{fig:ug-consola}):

\begin{itemize}
  \item \textbf{Gestión}: las herramientas «Usuarios», «Archivos»,
        «Paneles», «Auditoría» y «Ajustes» que su rol permita (el
        super-administrador ve además «Inicio» y «Gestión de cuentas»).
  \item \textbf{Datos}: una carpeta por \emph{área de negocio} ---«Blog»,
        «Cuentas», «Facturación», «Sistema» y «Otros datos»--- con las tablas
        de cada área, y al final «Todas las tablas».
\end{itemize}

En la cabecera, el buscador abre la \textbf{búsqueda global}
(\texttt{Ctrl+K} o \texttt{Cmd+K}), que busca un texto en todas las tablas que
usted puede leer. Cada pantalla muestra un breve resumen de para qué sirve, y
junto a los campos de los formularios aparece un icono «?» con una ayuda sobre
su contenido.

\begin{figure}[htbp]
  \centering
  \begin{tikzpicture}[
      font=\small,
      box/.style={draw, rounded corners=2pt, text width=3.95cm, minimum height=0.55cm, anchor=west, inner sep=3pt},
      head/.style={font=\small\bfseries, anchor=west},
    ]
    \node[draw, thick, rounded corners=4pt, minimum width=4.8cm, minimum height=9.6cm, anchor=north west] (frame) at (0,0) {};
    \node[head] at (0.25,-0.4) {Gestión};
    \node[box] at (0.3,-0.95) {Inicio \hfill\textit{(super-admin)}};
    \node[box] at (0.3,-1.55) {Gestión de cuentas \hfill\textit{(super-admin)}};
    \node[box] at (0.3,-2.15) {Usuarios};
    \node[box] at (0.3,-2.75) {Archivos};
    \node[box] at (0.3,-3.35) {Paneles};
    \node[box] at (0.3,-3.95) {Auditoría};
    \node[box] at (0.3,-4.55) {Ajustes};
    \node[head] at (0.25,-5.25) {Datos};
    \node[box] at (0.3,-5.8) {$\triangleright$ Blog};
    \node[box] at (0.3,-6.4) {$\triangleright$ Cuentas};
    \node[box] at (0.3,-7.0) {$\triangleright$ Facturación};
    \node[box] at (0.3,-7.6) {$\triangleright$ Sistema};
    \node[box] at (0.3,-8.2) {$\triangleright$ Otros datos};
    \node[box] at (0.3,-8.8) {Todas las tablas};
    \node[draw, dashed, rounded corners=4pt, minimum width=7.6cm, minimum height=9.6cm, anchor=north west, align=center] (main) at (5.2,0) {};
    \node[align=center, text width=7cm] at (main.center) {Área de trabajo:\\ listado de una tabla, ficha de un\\ registro, explorador de archivos,\\ panel, auditoría o ajustes};
    \node[draw, rounded corners=2pt, anchor=north, minimum width=7.2cm] at ([yshift=-0.2cm]main.north) {Búsqueda global (Ctrl/Cmd+K)};
  \end{tikzpicture}
  \caption{Esquema de la consola de administración: bloques Gestión y Datos.}
  \label{fig:ug-consola}
\end{figure}

\subsection{Explorador de datos}
\label{sec:ug-datos}

Al elegir una tabla en el bloque «Datos» se abre su \textbf{listado}. Las
tablas y columnas aparecen con un nombre en español y, junto a él, su nombre
técnico. Desde el listado se puede:

\begin{itemize}
  \item paginar, ordenar por cualquier columna y buscar por texto;
  \item filtrar por columnas, con filtros adaptados al tipo de dato (fechas
        relativas como «últimos 7 días», valores de una lista o registros
        relacionados con autocompletado);
  \item elegir qué columnas se ven, fijarlas y reordenarlas, y guardar la
        combinación de filtros y columnas como una \textbf{vista guardada}
        para reutilizarla;
  \item seleccionar varios registros y eliminarlos en bloque, o editar un
        valor directamente en la celda (edición en línea).
\end{itemize}

Al pulsar un registro se abre su \textbf{ficha}, que muestra sus campos
agrupados según el diseño configurado y, debajo, los «Registros
relacionados» (por ejemplo, los miembros de una cuenta). Desde la ficha se
puede editar o eliminar el registro, crear un registro relacionado o vincular
y desvincular registros existentes. Para crear un registro nuevo se usa el
botón de alta del listado; el formulario valida los datos antes de guardarlos
y avisa si se intenta salir con cambios sin guardar.

\paragraph{Ejemplo: publicar una entrada del blog.}
Abra «Blog» $\rightarrow$ «Entradas», cree un registro nuevo, rellene el
título, el identificador para la dirección (\textit{slug}), el resumen, el
contenido en Markdown, la categoría, la imagen de portada y el estado de
publicación, y guárdelo. La entrada
aparecerá en \texttt{/blog} en cuanto esté publicada.

Todas las altas, modificaciones y bajas hechas desde la consola quedan
registradas en la auditoría con su autor, la fecha y los valores anteriores y
nuevos.

\subsection{Usuarios}
\label{sec:ug-usuarios}

«Usuarios» lista las personas registradas en la plataforma, con búsqueda y
filtros. La ficha de un usuario muestra sus datos y sus factores de
autenticación. Según los permisos, se puede crear o invitar a un
usuario, bloquearlo o desbloquearlo, enviarle un correo para restablecer la
contraseña o un enlace de inicio de sesión, quitarle un factor de dos pasos,
eliminarlo y conceder o revocar su acceso al CMS. Las acciones también pueden
aplicarse a varios usuarios a la vez. Cada acción pide confirmación.

\subsection{Archivos}
\label{sec:ug-archivos}

«Archivos» permite navegar por los contenedores de almacenamiento
(\textit{buckets}) y sus carpetas, subir archivos («Subir archivos»), crear
carpetas («Nueva carpeta»), renombrar y eliminar archivos y ver una vista
previa de los que la admiten, como imágenes y documentos PDF.

\subsection{Paneles}
\label{sec:ug-paneles}

«Paneles» reúne los paneles propios y los que otros han compartido con
usted. Un panel es una cuadrícula de \textit{widgets} de tres tipos: métrica
(un número, como las suscripciones activas), gráfica y tabla. Para crearlo,
pulse «Nuevo panel», añada \textit{widgets} con «Añadir widget» eligiendo la
tabla, los filtros y el tipo de representación, y colóquelos en la cuadrícula.
Un panel se puede renombrar, compartir con otros roles y eliminar. Cada
\textit{widget} muestra solo los datos que puede leer quien lo está viendo.

\subsection{Auditoría}
\label{sec:ug-auditoria}

«Auditoría» muestra los «Registros de auditoría»: quién hizo cada cambio,
cuándo, sobre qué tabla y registro, y cómo quedaron los datos antes y después.
Se puede filtrar por fecha, autor («Realizado por»), operación, tabla y
gravedad, y abrir el detalle de
una entrada para ver las diferencias. Las entradas las escribe el sistema y
nadie puede modificarlas; los datos que su rol no puede leer aparecen
ocultos.

\section{Super-administrador}
\label{sec:ug-superadmin}

El super-administrador es la cuenta raíz de la plataforma y del CMS: tiene
todos los permisos sin necesidad de configurarlos. Además de las herramientas
anteriores, dispone de:

\subsection{Inicio y gestión de cuentas}

\begin{itemize}
  \item \textbf{Inicio} (\texttt{/admin}): métricas de la plataforma
        ---«Usuarios», «Cuentas de equipo», «Clientes de pago» y «Pruebas»---,
        con la advertencia de que son estimaciones.
  \item \textbf{Gestión de cuentas} (\texttt{/admin/accounts}): listado de
        todas las cuentas, personales y de equipo, con búsqueda («Buscar
        cuenta\ldots») y filtro por tipo. En la ficha de una cuenta se ven sus
        equipos, sus miembros y su suscripción, y se puede:
        \begin{itemize}
          \item «Suspender» o «Reactivar» al usuario;
          \item «Suplantar» al usuario para ver la aplicación como él y darle
                soporte;
          \item «Enviar enlace de restablecimiento de contraseña»;
          \item eliminar la cuenta personal o la de equipo; y
          \item «Crear usuario», con la opción de confirmar su correo
                automáticamente.
        \end{itemize}
        Las acciones delicadas piden escribir la palabra
        \texttt{CONFIRMAR}.
\end{itemize}

\subsection{Ajustes de la consola}
\label{sec:ug-ajustes}

«Ajustes» agrupa la configuración del CMS en varias pestañas:

\begin{itemize}
  \item \textbf{General}: preferencias de la consola, como la zona horaria.
  \item \textbf{Autenticación}: si la consola exige verificación en dos pasos
        a todo el personal. Está activado por defecto y solo la cuenta raíz,
        con la sesión verificada en dos pasos, puede desactivarlo, porque
        afecta a toda la consola.
  \item \textbf{Miembros}: el personal con acceso a la consola. Permite
        «Cambiar rol», activar o desactivar a un miembro y consultar su
        auditoría. Nadie puede cambiar sus propios roles ni asignar un rol de
        rango superior al suyo.
  \item \textbf{Roles y permisos}: roles del CMS (por ejemplo, «Soporte»),
        grupos de permisos y permisos individuales sobre tablas,
        almacenamiento y herramientas. Un permiso puede conceder o denegar
        explícitamente una acción.
  \item \textbf{Recursos}: la configuración de cada tabla en la consola.
        \begin{itemize}
          \item «Sincronizar tablas» registra en el CMS las tablas nuevas de la
                base de datos.
          \item Para cada tabla se puede fijar su nombre visible y su
                descripción, su \textbf{área} en la barra lateral (por ejemplo,
                «Blog», o un área nueva para las tablas propias de la pyme),
                si es visible y su orden.
          \item Para cada columna, su nombre, su descripción ---que se muestra
                como ayuda «?» en los formularios---, si es visible o editable y
                su formato de visualización.
          \item «Registros relacionados» decide qué relaciones se muestran en
                la ficha, y «Diseño del registro» permite organizar los campos
                de la ficha en secciones arrastrándolos.
        \end{itemize}
\end{itemize}

\paragraph{Ejemplo: dar acceso limitado a una persona de soporte.}
\begin{enumerate}
  \item En «Usuarios», abra la ficha de la persona y elija «Conceder acceso al
        CMS».
  \item En «Ajustes» $\rightarrow$ «Miembros», asígnele el rol «Soporte» (o
        cree antes un rol a medida en «Roles y permisos» con los permisos de
        lectura de las tablas que necesite).
  \item Pídale que configure su segundo factor. Al entrar en la consola solo
        verá las áreas, tablas y herramientas que su rol le permite.
\end{enumerate}

\newpage{\ }
\cleardoublepage

\chapter{Evaluación de la reutilización}
\label{cap:evaluacion}
% Se redacta tras el escenario de ejemplo (P-01).

    
    
    

% Fin del documento
\end{document}
